% !TeX program = xelatex
% Anonymous standalone source; bibliography is embedded.
\documentclass[11pt,a4paper]{article}
\usepackage[margin=26mm]{geometry}
\usepackage{amsmath,amssymb,amsthm,mathtools,booktabs,longtable}
\usepackage[T1]{fontenc}
\usepackage{newtxtext,newtxmath,microtype,xcolor}
\usepackage[unicode,hidelinks]{hyperref}
\usepackage{bookmark,enumitem,fancyhdr}
\numberwithin{equation}{section}
\setcounter{tocdepth}{2}
\setlist[enumerate]{leftmargin=2em,itemsep=4pt,topsep=5pt}
\pagestyle{fancy}
\fancyhf{}
\fancyhead[L]{\small High-frequency vacuum limits}
\fancyhead[R]{\small A variational approach}
\fancyfoot[C]{\thepage}
\setlength{\headheight}{14pt}
\renewcommand{\headrulewidth}{0pt}
\hypersetup{pdftitle={High-frequency vacuum limits with bounded first derivatives: a variational approach to the forward Burnett conjecture},pdfauthor={},pdfsubject={Anonymous manuscript; proposed proof}}

\setlength{\emergencystretch}{3em}
\allowdisplaybreaks
\newcommand{\Ric}{\operatorname{Ric}}
\newcommand{\Scal}{\operatorname{Scal}}
\newcommand{\tr}{\operatorname{tr}}
\newcommand{\Id}{\operatorname{Id}}
\newcommand{\supp}{\operatorname{supp}}
\newcommand{\loc}{\mathrm{loc}}
\newcommand{\dd}{\mathrm d}
\newcommand{\R}{\mathbb R}
\newcommand{\im}{\operatorname{Im}}
\newcommand{\Op}{\operatorname{Op}}
\newcommand{\dV}{\dd V_g}
\newcommand{\ip}[2]{\langle #1,#2\rangle}
\newcommand{\f}[2]{\frac{#1}{#2}}
\newcommand{\ls}{\lesssim}
\newcommand{\ii}{\mathrm{i}}
\DeclareMathOperator{\Rea}{Re}
\renewcommand{\Im}{\operatorname{Im}}
\newtheorem{theorem}{Theorem}[section]
\newtheorem{proposition}[theorem]{Proposition}
\newtheorem{lemma}[theorem]{Lemma}
\newtheorem{corollary}[theorem]{Corollary}
\theoremstyle{definition}
\newtheorem{definition}[theorem]{Definition}
\newtheorem{problem}[theorem]{Problem}
\newtheorem{remark}[theorem]{Remark}
\title{High-frequency vacuum limits with bounded first derivatives:\\
a variational approach to the forward Burnett conjecture}
\author{}
\date{}
\begin{document}
\maketitle
\vspace{-2.5em}
\begin{abstract}
We present a proposed proof of the forward Burnett conjecture for smooth
vacuum metrics on a four-dimensional manifold which converge locally
uniformly to a smooth Lorentzian metric and have locally uniformly bounded
first derivatives. No coordinate condition, spacetime symmetry, or bound
relating second derivatives to the amplitude is imposed. The claimed
conclusion is that the limiting Einstein tensor is the second moment of a
nonnegative null measure satisfying the massless Hamilton transport equation.
The quadratic part of the argument constructs this measure on the bundle
of positive cotangent rays and eliminates pure coordinate oscillations.
For transport, the Einstein-Hilbert Hessian reduces the problem to the
complete cubic variation of the action. A diffeomorphism identity, angular
estimates, and scalar normal-form corrections are then used to cancel this
residual. A projected wave equation provides the bounded source required
by the corrections, without estimating the background wave operator of
the original metric perturbation. On globally hyperbolic regions, a
characteristic section produces a full-momentum realization with finite
local second moments. The unrestricted transport argument is the central
claim of this proposed proof; reverse approximation and nonsmooth limits
are outside its scope.
\end{abstract}
\noindent\textbf{Keywords.} Einstein vacuum equations; Burnett conjecture;
microlocal defect measures; massless Vlasov equation; compensated compactness.

\smallskip
\noindent\textbf{2020 Mathematics Subject Classification.} 83C05, 35B27, 35L70, 35S05.

\tableofcontents
\clearpage

\section{Introduction}\label{sec:introduction}

\subsection{Main results and remarks}

We consider the high-frequency limit of the vacuum Einstein equations:
\begin{equation}\label{eq:admissible}
\begin{gathered}
\Ric(g_n)=0,\qquad
h_n\longrightarrow0\ \hbox{in }C^0_{\loc}(M),\\
\sup_n\|Dh_n\|_{L^\infty(K;r)}<\infty
\quad(K\Subset M).
\end{gathered}
\end{equation}
Here $M$ is a smooth, Hausdorff, second countable four-manifold,
$g_n$ are smooth Lorentzian metrics, and $h_n=g_n-g$.
The limit $g$ is smooth and time-oriented, with signature $(-,+,+,+)$.
We fix a smooth auxiliary Riemannian metric $r$ and set $D=\nabla^r$.
The norm in \eqref{eq:admissible} uses $r$ on the compact set $K$.

The cosmological constant is zero and $8\pi G_{\mathrm N}=1$.
Replacing $Dh_n$ by $Dg_n$ gives the same class. Changing $r$, or applying
a fixed smooth coordinate change, also leaves the hypotheses unchanged.
No estimate for $D^2g_n$ is assumed.

Let $\pi:T^*M\setminus0\to M$ be the bundle projection, and put
\begin{equation}\label{eq:hamiltonian}
\begin{gathered}
S_r^*M=\{(x,\xi)\in T^*M:|\xi|_r=1\},\qquad
P(x,\xi)=\f12g^{ab}(x)\xi_a\xi_b,\\
H_P=g^{ab}\xi_b\partial_{x^a}
-\f12\partial_{x^c}g^{ab}\xi_a\xi_b\partial_{\xi_c}.
\end{gathered}
\end{equation}
The unit cosphere represents positive rays, namely the quotient by
$\xi\mapsto t\xi$ for $t>0$; it retains both time orientations.
For a measure $\mu$ on this section, the Einstein moment identity is
\begin{equation}\label{eq:moment}
\int_M G(g)_{ab}S^{ab}\,\dV
=\int_{S_r^*M}S^{ab}(x)\xi_a\xi_b\,\dd\mu,
\quad S\in C_c^\infty(M;\operatorname{Sym}^2TM),
\end{equation}
where $G(g)=\Ric(g)-\f12\Scal(g)g$. Its transport equation is
\begin{equation}\label{eq:transport}
\int_{S_r^*M}H_Pa\,\dd\mu=0
\end{equation}
for every smooth $a:T^*M\setminus0\to\R$ with compact base support and
$a(x,t\xi)=t\,a(x,\xi)$ for $t>0$. The Hamilton derivative is taken
before restriction to the unit section; it is homogeneous of degree two.
Here ``positively homogeneous'' refers to positive scaling and imposes
no sign condition on $a$.

Our main result is stated as follows.

\begin{theorem}[Forward kinetic limit in the smooth-limit class]
\label{thm:main}
Assume \eqref{eq:admissible}. After passage to one subsequence, there exists
a nonnegative Radon measure $\mu$ on $S_r^*M$ such that:
\begin{enumerate}[label=\textup{(\roman*)}]
\item $\mu$ is locally finite over the base and supported on $P=0$;
\item the Einstein moment identity \eqref{eq:moment} holds, and
$\Scal(g)=0$;
\item the same measure satisfies \eqref{eq:transport} for every smooth
positively one-homogeneous test with compact base support.
\end{enumerate}
Moreover, on every globally hyperbolic open subspacetime $U\subset M$
there is a nonnegative Radon measure $\nu$ on $T^*U\setminus0$, supported
on its null cone, with
\begin{equation}\label{eq:full-moment-bound}
\int_{\pi^{-1}(K)}|\xi|_r^2\,\dd\nu<\infty
\qquad(K\Subset U),
\end{equation}
and
\begin{equation}\label{eq:full-equations}
\begin{gathered}
\int_U G(g)_{ab}S^{ab}\,\dV
=\int_{T^*U\setminus0}S^{ab}\xi_a\xi_b\,\dd\nu,\\
\int_{T^*U\setminus0}H_Pb\,\dd\nu=0
\qquad(b\in C_c^\infty(T^*U\setminus0)).
\end{gathered}
\end{equation}
The first equality holds for all compactly supported smooth symmetric
contravariant tensors $S$ on $U$.
\end{theorem}

\begin{remark}[Status of the main claim]\label{rmk:status}
This version presents a complete proposed proof of Theorem~\ref{thm:main}.
The decisive assertion is the unrestricted transport estimate in
Sections~\ref{sec:variational-transport}-\ref{sec:completion}, not merely
the existence of positive null second moments. The theorem and its technical lemmas are presented for independent
mathematical assessment.
\end{remark}

\begin{remark}[What the conclusion would settle]\label{rmk:scope}
Theorem~\ref{thm:main}, if its proof is correct, establishes the full
forward conclusion for the class \eqref{eq:admissible}, in arbitrary
coordinates and without symmetry. It is not restricted to a prescribed
finite family of oscillation phases. The kinetic object is a measure,
not necessarily a smooth density. No uniqueness of this measure is
asserted. The theorem does not address the reverse problem of constructing
vacuum approximations of prescribed kinetic solutions, nonsmooth limiting
metrics, concentration beyond the stated derivative bound, nonzero
cosmological constant, or dimensions other than four. Global hyperbolicity
is used only for the full-momentum realization, not for the global ray
measure or its transport equation.
\end{remark}

The proof also gives two independently stated intermediate results.
Theorem~\ref{thm:quadratic} constructs a positive null measure with the
Einstein moments directly from the full metric-gradient defect, including
its coordinate part. Proposition~\ref{prop:general-projection} gives a
scalar wave equation with bounded source for a fixed projection of the
metric perturbation. Neither result by itself establishes Vlasov
transport. The remaining step is exactly the vanishing of the cubic
residual in \eqref{eq:var-cubic-residual}.
Two consequences of the quadratic construction require no transport:
Corollary~\ref{cor:base-marginal} identifies the smooth base marginal
of the physical measure, and Corollary~\ref{cor:strong} characterizes
vacuum limits by strong convergence of the projected physical gradients.

\subsection{Previous results}

Weak limits of vacuum spacetimes need not satisfy the vacuum Einstein
equations. Although the metric perturbation may converge uniformly to
zero, its first derivatives can oscillate and leave a nonzero quadratic
stress in the limiting equation. The forward Burnett conjecture asks
whether this stress always has the dynamics of collisionless massless
matter. A kinetic description has two distinct components: a positive
null second-moment representation of the Einstein tensor and propagation
of its angular measure by the null geodesic flow. The latter is stronger
than covariant conservation of the stress tensor.

The weak-limit framework originates in Burnett's work
\cite{Burnett1989}. Green-Wald \cite{GW2011} established trace and energy
properties of the effective stress in their framework. Huneau-Luk
\cite{HLelliptic} proved a kinetic limit theorem with $\mathbb U(1)$
symmetry in elliptic gauge, using trilinear compensated compactness.
Guerra-Teixeira da Costa \cite{GTC} subsequently weakened the regularity
assumptions in that setting. Without symmetry, Huneau-Luk
\cite{HLwave} proved the conjecture in generalized wave coordinates under
quantitative convergence and derivative assumptions. These results identify
nonlinear transport, rather than the formal weak limit of the Einstein
equation, as the main analytical difficulty.

This paper addresses the forward problem with a smooth limiting metric
and only a uniform local bound on first derivatives. The main argument
is variational: it uses the complete third variation of the
Einstein-Hilbert action to retain the symmetry between all three
oscillatory factors. This permits a separate treatment of the physical
polarizations and the coordinate directions in the microlocal defect.



The hypothesis comparison requires some care. In
\cite[Assumptions 1.2]{HLwave}, a parameter $\lambda_n\to0$ satisfies
\begin{equation}\label{eq:literature-scale}
|g_n-g|\leq\lambda_n,\qquad |\partial g_n|\leq C,
\qquad |\partial^2g_n|\leq C\lambda_n^{-1}.
\end{equation}
There are also quantitative assumptions on
$H^a(k)=k^{bc}\Gamma^a_{bc}(k)$: convergence of $H(g_n)$ to $H(g)$ at a
positive power of $\lambda_n$, together with a first derivative bound
on $H(g_n)$. Remark 1.3 of that paper distinguishes its quadratic
conclusions from transport when the second-derivative and source-derivative
bounds are removed, with the remaining wave-gauge assumptions retained.
Our proposed transport argument uses neither a
coordinate condition nor \eqref{eq:literature-scale}.

The theorem of Guerra-Teixeira da Costa \cite{GTC} is formulated in
$\mathbb U(1)$ symmetry and elliptic gauge. Its published
Hypotheses 1.2(c) retain a product bound involving the uniform difference
of the inverse base metrics and the $L^4$ norm of second derivatives of
the reduced wave-map fields. Their first-derivative assumptions on those
fields are also weaker than the full-metric $W^{1,\infty}_{\loc}$ bound
used here. Accordingly, Theorem~\ref{thm:main} should not be described as
weakening every hypothesis of that result. The proposed advance concerns
removal of the symmetry, coordinate, and derivative-scale restrictions
in the overlapping smooth-limit class.

The distinction between positivity and transport is essential.
Covariant conservation of a positive null second moment tests only a
small subclass of the angular transport equation. Appendix~\ref{sec:rotating}
gives an explicit illustration. Conversely, failure of a commonly used
auxiliary estimate is not a counterexample to the kinetic limit theorem:
the exact flat sequence in Appendix~\ref{sec:shear} has an unbounded
background wave operator but has zero physical defect measure.

\subsection{Sketch of the proof}

The argument has three stages. First, the exact connection-difference
formula and microlocal compactness give a tensor-valued derivative defect
whose range lies in the kernel of the linearized Ricci symbol. Away from
the null cone that kernel is pure gauge and contributes no effective
stress. On the cone the relevant contraction is nonnegative on the
two-dimensional physical polarization space. This yields the global
measure $\mu$.

Second, set $E(g)^{ab}=\sqrt{|\det g|}\,G(g)^{ab}$.
A scalar order-one test of the density-valued Einstein-Hilbert Hessian gives
\begin{equation*}
4\int_{S_r^*M}H_Pa\,\dd\mu
=-\lim_n\Im
\ip{E''_g[h_n,h_n]}{B_ah_n}.
\end{equation*}
Here the formula is localized in one coordinate chart, the pairing uses
Lebesgue density, and $B_a$ is a scalar self-adjoint pseudodifferential
operator of order one with principal symbol $a$.
The entire third variation is retained. For frozen coefficients its
Fourier vertex is multiplied by the cyclic factor
$a(p)+a(q)+a(r)$, where $p+q+r=0$. Near a common null direction this
factor gains two angular powers. When a gauge factor is present,
diffeomorphism invariance converts the vertex into physical characteristic
factors and an exact scalar phase difference.

Third, Hermitian normal-form corrections handle the remaining separated
interactions. They are paired with a scalar equation for the projected
physical field. The coefficients of that equation are the original
inverse metrics and remain uniformly Lipschitz. The final estimate is
\begin{equation*}
\left|4\int_{S_r^*M}H_Pa\,\dd\mu\right|
\leq C_{\kappa,p_*}\delta^{\beta_*}
 +C_{\delta,p_*}\ell^{\alpha_*}+C_\delta\ell,
\end{equation*}
For any fixed finite exponent $p_*>3$, the exponents are
\begin{equation}\label{eq:optimized-exponents}
\beta_* = \f{2(p_*-3)}{p_*-2},\qquad
\alpha_* = \f{2(p_*-3)}{3(p_*-2)}.
\end{equation}
For example, $p_*=4$ gives $(\beta_*,\alpha_*)=(1,1/3)$,
and $p_*=12$ gives $(9/5,3/5)$.
The parameter $\kappa$ fixes the separation between low/high and comparable
frequencies. The sequence limit is taken first, at fixed angular width
$\delta$ and spatial patch size $\ell$. Sending $\ell\to0$ and then $\delta\to0$
gives the asserted transport equation. No sequence-dependent upper
frequency cutoff or second-derivative estimate enters this limit order.

Section~\ref{sec:quadratic} constructs the physical measure.
Section~\ref{sec:analytic} supplies the analytic estimates used for
cubic energies and normal forms, and Section~\ref{sec:general-projection}
establishes the projected wave equation.
Section~\ref{sec:variational-transport} identifies and estimates the
variational residual. Section~\ref{sec:completion} completes the
localization argument and proves Theorem~\ref{thm:main}.
The appendices give exact examples and auxiliary identities.

\begin{samepage}
\subsection{Notations}

\begin{description}[leftmargin=1.5em,labelsep=0.5em,itemsep=3pt]
\item[$\blacktriangleright$] $a,b,c\in\{0,1,2,3\}$ are tensor indices; repeated indices are summed.
\item[$\blacktriangleright$] $\tr_g$ is the Lorentzian trace; symmetric contractions sum over all ordered pairs.
\item[$\blacktriangleright$] $|\cdot|_r$ is the positive tensor length induced by the auxiliary metric $r$.
\item[$\blacktriangleright$] $\dV=\sqrt{|\det g|}\,\dd x$ is the geometric volume density.
\item[$\blacktriangleright$] $\ip{u}{v}=\int_{\R^4}u(x)\cdot\overline{v(x)}\,\dd x$ is linear in $u$ in a coordinate chart.
\item[$\blacktriangleright$] $L^p_x$, $H^s_x$, and $W^{s,p}_x$ denote the corresponding spaces on $\R^4$.
\item[$\blacktriangleright$] $P$ and $H_P$ are the Hamiltonian and vector field in \eqref{eq:hamiltonian}.
\item[$\blacktriangleright$] $\lambda_p=|p|$, $\lambda_q=|q|$, and $\lambda_r=|r|$ are Euclidean frequency magnitudes.
\item[$\blacktriangleright$] $A\ls B$ means $A\leq CB$ with the dependencies specified below.
\end{description}
\end{samepage}
Component lengths use an orthonormal frame for the specified positive metric;
the corresponding $L^2$ norms use the sum of squared component lengths.
Localized coordinate fields are extended by zero when a full-space norm is used.
Constants may depend on fixed compact sets, the smooth background,
the test symbol, and the first-derivative bound in \eqref{eq:admissible}.
They never depend on an unassumed bound for higher derivatives of $g_n$.

\section{Preliminaries}\label{sec:quadratic}

\begin{theorem}[Positive null second moments, without transport]\label{thm:quadratic}
Assume \eqref{eq:admissible}. After extraction of one subsequence, there is a globally defined nonnegative Radon measure $\mu$ on $S_r^*M$, locally finite over compact base sets and supported on $P=0$, such that \eqref{eq:moment} holds. In particular $\Scal(g)=0$. The construction is compatible with fixed smooth coordinate changes. The transport equation is not part of this intermediate theorem.
\end{theorem}

\subsection{Exact connection difference and compactness}

Use $\nabla=\nabla^g$. The difference $\nabla h_n-Dh_n$ is a smooth coefficient times $h_n$, hence is strongly zero in local $L^\infty$. Define
\begin{equation*}
C_n{}^a{}_{bc}=\f12g_n^{ad}(\nabla_bh_{n,cd}+\nabla_ch_{n,bd}-\nabla_dh_{n,bc})
=L(\nabla h_n)^a{}_{bc}+E_n{}^a{}_{bc},
\end{equation*}
where $L$ uses $g^{ad}$ in place of $g_n^{ad}$. On a fixed compact set $K\Subset M$, inverses are uniformly bounded for all sufficiently large $n$, and
\begin{equation}\label{eq:E}
\|E_n\|_{L^\infty(K)}\leq C\|h_n\|_{L^\infty(K)}\|\nabla h_n\|_{L^\infty(K)}\longrightarrow0.
\end{equation}
All constants here and below may depend on the fixed compact sets, the smooth background, and the given first-derivative bound, but not on higher derivatives of $g_n$.

The exact curvature-difference formula gives
\begin{equation}\label{eq:ricci-exact}
0=\Ric(g)_{bc}+\nabla_aC_n{}^a{}_{bc}-\nabla_cC_n{}^a{}_{ba}
+C_n{}^a{}_{ae}C_n{}^e{}_{bc}-C_n{}^a{}_{ce}C_n{}^e{}_{ba}.
\end{equation}
Distributional differentiation of $h_n\to0$ and the derivative bound imply $\nabla h_n\rightharpoonup^*0$ in local $L^\infty$. Consequently $C_n\rightharpoonup0$ in local $L^2$. Testing \eqref{eq:ricci-exact} against a fixed smooth compactly supported tensor, with density $\dV$, and integrating the derivatives by parts yields
\begin{equation}\label{eq:weak-ricci}
\Ric(g)_{bc}=\mathop{\mathrm{w\!-!lim}}_n
\bigl(L(\nabla h_n)^a{}_{ce}L(\nabla h_n)^e{}_{ba}
-L(\nabla h_n)^a{}_{ae}L(\nabla h_n)^e{}_{bc}\bigr).
\end{equation}
Every replacement of $C_n$ by $L(\nabla h_n)$ has a strongly vanishing local $L^\infty$ error by \eqref{eq:E}. Thus no second derivative has been discarded from a product with a sequence-dependent test.

Let $\mathcal Lh=\nabla_aL(\nabla h)^a{}_{bc}-\nabla_cL(\nabla h)^a{}_{ba}$. Equation~\eqref{eq:ricci-exact} writes $\mathcal Lh_n$ as the sum of a bounded local $L^2$ sequence and the derivative of a strongly zero local $L^2$ sequence. On a bounded chart, with nested compact cutoffs, the first part is precompact in $H^{-1}$ by Rellich duality, and the second part is strongly zero in $H^{-1}$. Derivatives of the fixed cutoffs are bounded local $L^2$ terms. Since independently $\mathcal Lh_n\to0$ distributionally, all its $H^{-1}$ cluster points are zero. Therefore
\begin{equation}\label{eq:ricci-Hminus}
\mathcal Lh_n\longrightarrow0\quad\text{strongly in }H^{-1}_{\loc}.
\end{equation}
This does not say that the quadratic terms in \eqref{eq:weak-ricci} vanish.

\subsection{A common microlocal measure and the Ricci constraint}

Apply the standard positive microlocal-defect compactness theorem
\cite[Theorem~5.5]{HLelliptic} to the global tensor $u_n=Dh_n$, using
the positive $r$ fiber inner product and density $\dV$.
Here defect tests are classical order-zero operators: their smooth
principal symbols are homogeneous of degree zero away from the origin,
with fixed infrared regularization and an order-minus-one remainder.
A countable relatively compact atlas, a compact exhaustion, and
countable dense families of these principal symbols on compact
cosphere sets give a single diagonal subsequence. Local $L^2$ bounds
and order-zero operator bounds, using finitely many symbol derivatives,
extend its quadratic limits to every such smooth compactly based
principal symbol. Localized order-minus-one remainders are compact
and have zero limit. This produces a positive Hermitian bundle-valued
measure $M_E$ on $S_r^*M$, where
$E=T^*M\otimes\operatorname{Sym}^2T^*M$.
No simultaneous convergence for arbitrary nonhomogeneous
$S^0_{1,0}$ symbols is asserted. The homogeneous projections and
normalized phase symbols used below belong to this defect-test class;
general multipliers used only in norm estimates need no defect limit.

The localization fact needed here, also stated in \cite[Theorem~5.6]{HLelliptic}, can be checked directly: if a smooth first-order system $A$ satisfies $Au_n\to0$ in $H^{-1}_{\loc}$, apply a fixed elliptic inverse of order $-1$ and insert nested cutoffs. The resulting order-zero operator applied to $u_n$ tends strongly to zero in $L^2$; the cutoff commutators have order $-1$ and are compact. Its principal symbol is $a_1(x,\xi)/|\xi|_r$. Positivity then implies that the range of the density of $M_E$ lies in $\ker a_1$ almost everywhere.

Apply this fact first to the curl of $Dh_n$. Its right-hand side is the fixed curvature of $D$ acting on $h_n$, and tends strongly to zero. Thus the density takes values in tensors $u_{\alpha,ab}=\xi_\alpha k_{ab}$. Apply it next to \eqref{eq:ricci-Hminus}, viewed as a first-order constraint on $Dh_n$. Smooth lower-order terms are strongly zero in $H^{-1}$, because they are weakly zero and bounded in $L^2$ on a bounded chart. The remaining constraint is
\begin{equation}\label{eq:ricci-symbol}
\mathcal R_\xi(k)_{ab}=\f12\bigl(s k_{ab}+\xi_a\xi_b t-\xi_aq_b-\xi_bq_a\bigr)=0,
\quad s=g^{ab}\xi_a\xi_b,\quad t=\tr_gk,\quad q_b=\xi^ak_{ab}.
\end{equation}
An overall Fourier sign does not change its kernel.

Since $|\xi|_r=1$, the map $j_\xi:k\mapsto\xi\otimes k$ is an $r$-isometry. There is therefore a positive Hermitian measure $\sigma$ on symmetric two-tensors such that $M_E=j_\xi\sigma j_\xi^*$. Its density has range in $\ker\mathcal R_\xi$. This operator identity uses $r$ to identify each fiber with its dual. In the coordinate formulas below, $\sigma_{ab,cd}$ denotes the corresponding covariance tensor: if $G_F$ is the fiber Gram matrix, its matrix is $\sigma_{\rm op}G_F^{-1}$, rather than the endomorphism matrix $\sigma_{\rm op}$ itself. This convention also applies when changing the auxiliary metric. For a coordinate chart $U$ containing the base support of $A$,
\begin{equation}\label{eq:sigma}
\lim_n\ip{\partial_\alpha h_{n,ab}}{A\partial_\beta h_{n,cd}}_{L^2(U,\dV)}
=\int_{S_r^*U} a_0(x,\xi)\xi_\alpha\xi_\beta\,\dd\sigma_{ab,cd},
\end{equation}
for classical order-zero $A$ with real scalar principal symbol $a_0$;
complex tests follow by conjugate-linearity in $A$ in the stated
inner-product convention. Tensor contractions here sum over all ordered index pairs, so the off-diagonal multiplicities of symmetric tensors are included. Replacing $D$ by coordinate derivatives changes only strongly zero terms. Tensor component changes and conjugation of a fixed pseudodifferential operator transform the principal symbol naturally; lower-order terms have zero quadratic limit. The local descriptions thus describe one global object, rather than unrelated representations glued by a partition of unity.

\subsection{The full quadratic algebra}

At a fixed point, set
\begin{equation*}
c(k)^a{}_{bc}=\f12(\xi_bk_c{}^a+\xi_ck_b{}^a-\xi^ak_{bc}),\quad
B_{bc}(k,k)=c(k)^a{}_{ce}c(k)^e{}_{ba}-c(k)^a{}_{ae}c(k)^e{}_{bc}.
\end{equation*}
Write $K_2=k_{ab}k^{ab}$ and $(k^2)_{bc}=k_b{}^ak_{ac}$. Expanding the first product, its nine terms, before division by four, are
\begin{equation*}
\begin{gathered}
\xi_b\xi_cK_2,\quad \xi_c(kq)_b,\quad-\xi_c(kq)_b,\\
\xi_b(kq)_c,\quad q_bq_c,\quad-s(k^2)_{bc},\\
-\xi_b(kq)_c,\quad-s(k^2)_{bc},\quad q_bq_c.
\end{gathered}
\end{equation*}
Also $c(k)^a{}_{ae}=\xi_et/2$. It follows that
\begin{equation}\label{eq:B-algebra}
4B_{bc}(k,k)=\xi_b\xi_cK_2+2q_bq_c-2s(k^2)_{bc}
-t(\xi_bq_c+\xi_cq_b-sk_{bc}).
\end{equation}
The associated Hermitian forms follow by polarization of this real identity.

If $s\ne0$, \eqref{eq:ricci-symbol} gives $k=\xi\otimes v+v\otimes\xi$ with $v=(q-t\xi/2)/s$. For this pure gauge tensor, $c(k)^a{}_{bc}=\xi_b\xi_cv^a$, and both products in $B$ equal $\xi_b\xi_c(\xi^av_a)^2$. Thus $B$ vanishes on the whole kernel, including all mixed terms by polarization.

If $s=0$, \eqref{eq:ricci-symbol} is equivalent to $q=t\xi/2$, and \eqref{eq:B-algebra} reduces to
\begin{equation}\label{eq:B-null}
B_{bc}(k,k)=\f14\xi_b\xi_c F(k,k),\quad
F(k,k)=k_{ab}k^{ab}-\f12(\tr_gk)^2.
\end{equation}
For every $k$ in that kernel and every gauge tensor $\ell=\xi\otimes v+v\otimes\xi$,
\begin{equation*}
F(k,\ell)=2q^av_a-(\tr_gk)\xi^av_a=0.
\end{equation*}
This verifies physical/gauge cross cancellation, not just pure-gauge self cancellation.

Choose a null frame $L=\xi^\sharp,\underline L,e_1,e_2$ with $g(L,\underline L)=-2$ and orthonormal screen. The constraint gives $k_{LL}=k_{LA}=0$ and $k_{11}+k_{22}=0$. Gauge additions independently set $k_{L\underline L},k_{\underline LA},k_{\underline L\underline L}$ to zero. The remaining tensor is a trace-free symmetric screen tensor, so $F(k,k)=\sum_{A,B}|k_{AB}|^2\geq0$. This proves nonnegativity on the physical quotient and invariance under the chosen representative.

\subsection{Construction, positivity, and compatibility}

Let $Z=\{P=0\}\cap S_r^*M$. Define
\begin{equation}\label{eq:mu-invariant}
\dd\mu=\f14\mathbf1_Z\bigl(g^{ac}g^{bd}-\f12g^{ab}g^{cd}\bigr)\dd\sigma_{ab,cd}.
\end{equation}
Restriction to the closed set $Z$ is a Borel restriction of a measure; no nonsmooth pseudodifferential projection is being asserted. Off $Z$, the unrestricted contraction need not be nonnegative and must not be used.

Dominate $\sigma$ by its positive $r$ trace measure. Its positive Hermitian density can be diagonalized in each finite-dimensional fiber, and every nonzero eigenvector belongs to $\ker\mathcal R_\xi$. On $Z$, each contribution to \eqref{eq:mu-invariant} is nonnegative by \eqref{eq:B-null}. Smooth coefficients are bounded on compact unit-cosphere sets, proving local finiteness. Equation~\eqref{eq:weak-ricci}, together with \eqref{eq:B-algebra} off $Z$ and \eqref{eq:B-null} on $Z$, proves
\begin{equation*}
\int_M\Ric(g)_{ab}S^{ab}\dV=\int_{S_r^*M}S^{ab}\xi_a\xi_b\,\dd\mu.
\end{equation*}
Testing with $S^{ab}=\phi g^{ab}$ gives $\Scal(g)=0$, proving Theorem~\ref{thm:quadratic}.

For another auxiliary metric $r'$, let $s_0(x,\xi)=|\xi|_{r'}$ on $S_r^*M$ and $F_0(x,\xi)=(x,\xi/s_0)$. The covariance of the full angular symbol functional in \eqref{eq:sigma}, and the factorization $M_E=\xi\otimes\xi\otimes\sigma$, give
\begin{equation}\label{eq:norm-change}
\sigma'=F_{0*}(s_0^2\sigma),\qquad \mu'=F_{0*}(s_0^2\mu).
\end{equation}
This argument uses all angular symbols. Equality of stress second moments alone would not establish this transformation, since those moments do not determine a measure uniquely.

\begin{corollary}[Base marginal and angular disintegration]
\label{cor:base-marginal}
Under the hypotheses of Theorem~\ref{thm:quadratic}, put
$e_r=r^{ab}G(g)_{ab}$. Then $e_r$ is smooth and nonnegative, and
\begin{equation}\label{eq:base-marginal}
\pi_*\mu=e_r\,\dV.
\end{equation}
There is a measurable family of nonnegative finite measures $\mu_x$ on
the unit null covectors at $x$, defined for almost every $x$, such that
\begin{equation*}
\dd\mu(x,\xi)=\dd\mu_x(\xi)\,\dV(x),\qquad
\mu_x(S_{r,x}^*M)=e_r(x),\qquad
G(g)_{ab}(x)=\int_{S_{r,x}^*M}\xi_a\xi_b\,\dd\mu_x(\xi).
\end{equation*}
No angular regularity or determination by the moment tensor alone is asserted.
For a fixed $\mu$, disintegration has its usual almost-everywhere uniqueness.
In particular,
the physical measure has no concentration in its base marginal, and it
vanishes over every open vacuum region of $g$.
\end{corollary}
\begin{proof}
In \eqref{eq:moment} take $S^{ab}=\phi r^{ab}$, with
$\phi\in C_c^\infty(M)$. On the unit section $r^{ab}\xi_a\xi_b=1$,
which proves \eqref{eq:base-marginal}. Positivity and smoothness give
$e_r\geq0$ everywhere. Disintegrate the locally finite positive measure
over its base marginal, first on a compact exhaustion; the fiber
probability measures, multiplied by $e_r(x)$, give $\mu_x$ relative to
$\dV$. Set $\mu_x=0$ where $e_r=0$. Testing each tensor component of
\eqref{eq:moment} gives the last identity almost everywhere.
None of these assertions uses transport.
\end{proof}

\section{Analytic estimates for cubic energies}\label{sec:analytic}

We collect estimates whose constants are independent of the number of
dyadic frequency scales. All fields in this section take values in a
fixed finite-dimensional real vector space, complexified when Fourier
transforms are used. All coordinate variables range over $\R^4$; fields supported in a chart
are extended by zero. In proofs, $L^p_x=L^p(\R^4,\dd x)$ and
$W^{j,p}_x=W^{j,p}(\R^4)$. Here $\partial$ is the full coordinate gradient;
$D_j=-\ii\partial_j$ and $b(D)$ are Fourier multiplier notation.
Smooth coefficients are extended with fixed nested cutoffs when global
norms are used. Pairings use Lebesgue density and the fixed Euclidean
fiber inner product. The metric $m$ is a
constant Lorentz form in a compact family of nondegenerate metrics, and
$q_m(\xi)=m^{ab}\xi_a\xi_b$; for covectors we write
$m(p,\xi)=m^{ab}p_a\xi_b$. All infrared regularizations are fixed.
The constants in the continuity, transfer, and correction estimates are
uniform on bounded sets of the finitely many normalized symbol derivatives
used in their proofs. In the angular estimate those bounds are uniform
in the cap index and in $\delta$.

For $f\in\mathcal S(\R^4)$, the Fourier convention is
\begin{equation*}
\begin{gathered}
\widehat f(\xi)=\int_{\R^4}e^{-\ii x\cdot\xi}f(x)\,\dd x,\\
f(x)=(2\pi)^{-4}\int_{\R^4}e^{\ii x\cdot\xi}\widehat f(\xi)\,\dd\xi.
\end{gathered}
\end{equation*}
A bilinear symbol $b$ acts initially on Schwartz fields by
\begin{equation*}
T_b(f,g)(x)=(2\pi)^{-8}\int_{\R^4\times\R^4}
 e^{\ii x\cdot(p+\xi)}b(p,\xi)\widehat f(p)\widehat g(\xi)
 \,\dd p\,\dd\xi.
\end{equation*}
Tensor components are contracted according to $b$; we also write $B_b=T_b$.
The corresponding trilinear form, with bilinear rather than Hermitian
contraction in its last factor, is
\begin{equation*}
\begin{aligned}
\mathcal T_b(f,g,w)
&=\int_{\R^4} T_b(f,g)(x)w(x)\,\dd x\\
&=(2\pi)^{-8}\int_{\R^4\times\R^4}
 b(p,q)\widehat f(p)\widehat g(q)\widehat w(-p-q)\,\dd p\,\dd q.
\end{aligned}
\end{equation*}

\subsection{First-order commutators and cubic continuity}

\begin{lemma}[Lipschitz-small coefficients in a symmetric energy pairing]\label{lem:commutator}
Let $u$ be a compactly supported smooth vector-valued function, let $r^{ab}=r^{ba}$ be real smooth scalar coefficients in the field indices, and let $B=B^*$ be a fixed properly supported matrix pseudodifferential operator of order one with smooth $S^1_{1,0}$ symbol. Then
\begin{equation}\label{eq:lip-estimate}
\bigl|2\Im\ip{r^{ab}\partial_{ab}u}{Bu}\bigr|
\leq C_B\|r\|_{W^{1,\infty}(\R^4)}\|u\|_{H^1(\R^4)}^2.
\end{equation}
The constant is independent of second derivatives of $r$ or $u$.
The same estimate holds for matrix coefficients $R^{ab}$ with
$(R^{ab})^*=R^{ba}$ if $B$ is scalar at every symbol order.
\end{lemma}

\begin{proof}
Use the inner product linear in its first entry and initially Lebesgue density. Set $L_r=-\partial_a(r^{ab}\partial_b)$, which is formally self-adjoint. The difference $r^{ab}\partial_{ab}+L_r=-(\partial_ar^{ab})\partial_b$ obeys \eqref{eq:lip-estimate} directly. For the other part,
\begin{equation*}
2\Im\ip{L_ru}{Bu}=\ip{\ii [L_r,B]u}{u},
\end{equation*}
and the exact operator identity is
\begin{equation*}
[L_r,B]=-\partial_a\circ r^{ab}\circ[\partial_b,B]
-\partial_a\circ[r^{ab},B]\circ\partial_b-[\partial_a,B]\circ r^{ab}\circ\partial_b.
\end{equation*}
In the first and third terms, move the outside first-order operator to the other factor; $[\partial_a,B]$ has order one, giving a bound by $C\|r\|_{L^{\infty}_x}\|u\|_{H^1_x}^2$. In the middle term, move $\partial_a$ and use the order-one Calder\'on commutator bound
\begin{equation*}
\|[r^{ab},B]v\|_{L^2_x}\leq C_B\|r\|_{W^{1,\infty}_x}\|v\|_{L^2_x}.
\end{equation*}
This is the classical Lipschitz commutator theorem, stated for $S^1_{1,0}$ operators in \cite[Theorem A]{BenyiOh}; it applies entry by entry to the fixed matrix operator. It uses no higher coefficient derivative. Conjugation by the square root of a fixed positive smooth density reduces that case to the one proved, introducing only bounded lower-order terms.
\end{proof}


The matrix-coefficient version also holds when
$(R^{ab})^*=R^{ba}$ and $B$ is scalar at every symbol order. In that case
$-\partial_a(R^{ab}\partial_b)$ is self-adjoint and the same proof
applies entry by entry. Scalarity prevents an algebraic matrix
commutator in $[R^{ab},B]$.

For the next lemma, let $A^{ab}$ be constant linear maps from the field
space to its endomorphisms, with
$(A^{ab}(H))^{\mathsf T}=A^{ba}(H)$ for every real field value $H$.
Let $\mathcal Q$ be any fixed constant-coefficient bilinear expression
in first derivatives, and let $B=a(D)\Id$ have a real smooth homogeneous
degree-one symbol. Define
\begin{equation*}
\mathcal E(u)=2\Im\left\langle
A^{ab}(u)\partial_{ab}u+\mathcal Q(\partial u,\partial u),Bu
\right\rangle.
\end{equation*}

\begin{lemma}[Continuity of the complete cubic energy]\label{lem:cubic-continuity}
For smooth compactly supported fields,
\begin{equation}\label{eq:cubic-continuity}
|\mathcal E(u)-\mathcal E(v)|
\leq C\|\partial(u-v)\|_{L^3(\R^4)}
(\|\partial u\|_{L^3(\R^4)}+\|\partial v\|_{L^3(\R^4)})^2.
\end{equation}
With fixed inhomogeneous regularizations, the corresponding estimate
holds with $W^{1,3}(\R^4)$ norms. The estimate applies to the complete
polarized energy; it does not estimate isolated unsymmetrized pairings.
\end{lemma}

\begin{proof}
The finite-exponent first-order commutator estimate needed here is
\begin{equation*}
\|[b(D),f]v\|_{L^r_x}\leq C_b\|\partial f\|_{L^p_x}\|v\|_{L^s_x},
\qquad 1<p,s,r<\infty,\qquad r^{-1}=p^{-1}+s^{-1},
\end{equation*}
for a smooth homogeneous degree-one $b$. Here is a direct global
Coifman-Meyer decomposition, including arbitrarily small product
frequencies. Let $\eta$ be the frequency of $f$, and choose a smooth
degree-zero cutoff $\chi(\eta,\xi)$ equal to one when
$|\eta|\leq|\xi|/8$ and supported where $|\eta|\leq|\xi|/4$.
Put $\varrho=1-\chi$, $D_j=-\ii \partial_j$, and
$c_j(\zeta)=b(\zeta)\zeta_j/|\zeta|^2$ for $\zeta\ne0$.
Writing $T_m$ for the bilinear multiplier with symbol $m$, define
\begin{equation*}
 m_j=\chi\int_0^1\partial_jb(\xi+t\eta)\,\dd t,
 \qquad M_j=\varrho\f{\eta_j+\xi_j}{|\eta|},
 \qquad N_j=\varrho\f{\xi_j}{|\eta|}.
\end{equation*}
All three are global order-zero Coifman-Meyer symbols. On the support
of $\chi$, the integration segment stays a fixed fraction of
$|\eta|+|\xi|$ away from zero. On the support of $\varrho$,
$|\eta|\geq c(|\eta|+|\xi|)$; hence $M_j,N_j$ extend smoothly by zero
near $\eta=0$ on the joint unit sphere. These facts bound every fixed
number of Cartesian symbol derivatives, including on the axes and at
$\eta+\xi=0$. The exact identity is
\begin{equation*}
[b(D),f]v
=\sum_j T_{m_j}(D_jf,v)
 +\sum_j c_j(D)T_{M_j}(|D|f,v)
 -\sum_j T_{N_j}(|D|f,c_j(D)v).
\end{equation*}
Indeed the first sum is the integral formula for the difference of
$b$ on the $\chi$ region, while the other sums use
$b(\zeta)=\sum_j\zeta_jc_j(\zeta)$ on its complement.
The bilinear Coifman-Meyer theorem, the linear Mikhlin theorem,
and the Riesz-transform bounds give
\begin{equation*}
\begin{aligned}
\|[b(D),f]v\|_{L^r_x}
&\leq\sum_{j=0}^3\bigg(
 \|T_{m_j}(D_jf,v)\|_{L^r_x}
 +\|c_j(D)T_{M_j}(|D|f,v)\|_{L^r_x}
 +\|T_{N_j}(|D|f,c_j(D)v)\|_{L^r_x}\bigg)\\
&\leq C_b\bigg(\sum_{j=0}^3\|D_jf\|_{L^p_x}
                   +\||D|f\|_{L^p_x}\bigg)\|v\|_{L^s_x}
 \leq C_b\|\partial f\|_{L^p_x}\|v\|_{L^s_x}.
\end{aligned}
\end{equation*}
Only finitely many derivatives of $b$ on the unit sphere enter the
constant. The multiplier on the product is kept outside the bilinear symbol;
no differentiability at zero of $c_j$ and no summation over independent
low/high gaps is required.

For a real field $H$ put
$L_H=-\partial_a(A^{ab}(H)\partial_b)$. The operator is self-adjoint, and
$A^{ab}(H)\partial_{ab}=-L_H-(\partial_aA^{ab}(H))\partial_b$.
Since $B$ commutes with coordinate derivatives,
\begin{equation*}
[L_H,B]=-\partial_a[A^{ab}(H),B]\partial_b.
\end{equation*}
Apply the commutator estimate with $p=s=3$, $r=3/2$,
and move the outer derivative by integration by parts. H\"older's inequality gives
\begin{equation*}
\begin{aligned}
\big|2\Im\ip{A^{ab}(H)\partial_{ab}u}{Bu}\big|
&\leq \sum_{a,b}\|[A^{ab}(H),B]\partial_bu\|_{L^{3/2}_x}
                    \|\partial_au\|_{L^3_x}
       +C\|\partial H\|_{L^3_x}\|\partial u\|_{L^3_x}\|Bu\|_{L^3_x}\\
&\leq C\|\partial H\|_{L^3_x}\|\partial u\|_{L^3_x}^2.
\end{aligned}
\end{equation*}
Polarize the two $u$ factors together before estimating the cross terms.
The semilinear energy is bounded by H\"older's inequality in the same spaces.

Let $\mathfrak E$ be the real symmetric trilinear polarization of $\mathcal E$.
The preceding bounds give
\begin{equation*}
|\mathfrak E(H,u,v)|
\leq C\|\partial H\|_{L^3_x}\|\partial u\|_{L^3_x}\|\partial v\|_{L^3_x}.
\end{equation*}
Since $\mathcal E(u)=\mathfrak E(u,u,u)$, telescoping its three factors yields
\begin{equation*}
\begin{aligned}
|\mathcal E(u)-\mathcal E(v)|
&\leq|\mathfrak E(u-v,u,u)|+|\mathfrak E(v,u-v,u)|
       +|\mathfrak E(v,v,u-v)|\\
&\leq C\|\partial(u-v)\|_{L^3_x}
       \bigg(\|\partial u\|_{L^3_x}+\|\partial v\|_{L^3_x}\bigg)^2.
\end{aligned}
\end{equation*}
\end{proof}


The same argument gives a useful coefficient estimate. For a real smooth
scalar $f$, apply the commutator bound to $fA^{ab}(H)$.
The product rule and the same divergence-form estimate give
\begin{equation*}
\begin{split}
\left|2\Im\langle f A^{ab}(H)\partial_{ab}u,Bu\rangle\right|
\leq C\bigl(&\|f\|_{L^{\infty}_x}\|\partial H\|_{L^3_x}
+\|\partial f\|_{L^{\infty}_x}\|H\|_{L^3_x}\bigr)\|\partial u\|_{L^3_x}^2.
\end{split}
\end{equation*}
Polarization gives the corresponding estimate with distinct factors.
Thus an $O(\ell)$ coefficient difference gives an $O(\ell)$ leading
error; derivatives of that difference multiply an undifferentiated
field, which tends to zero at fixed localization scale. Smooth spatial
coefficients of a pseudodifferential test are treated by a summable
spatial and angular symbol expansion using this estimate.

\begin{lemma}[Quantitative freezing of a complete cubic energy]
\label{lem:freezing}
Work in a fixed relatively compact chart in $\mathbb R^4$ with a fixed
finite-dimensional real field space. Let $\mathcal T$ be a real symmetric
trilinear action form whose principal cubic polynomial is, up to its
fixed normalization,
\begin{equation}\label{eq:freezing-statement-1}
 P_C(z)=\int_{\R^4} C_{IJK}^{ab}(x)z^I\partial_a z^J\partial_b z^K\,\dd x,
 \qquad C_{IJK}^{ab}=C_{IKJ}^{ba},
\end{equation}
and whose remaining terms have at most one field derivative in total.
All coefficients are fixed and smooth. Let $B$ be a scalar self-adjoint
order-one test with real odd principal symbol $a(x,\xi)$, chosen so that
$Q=-\ii B$ preserves real fields. Use one fixed infrared regularization.
For $x_0$ in a fixed compact subset let $\mathcal T_0$ retain just the
principal coefficients frozen at $x_0$, and put $Q_0=-\ii a(x_0,D)$.
For real $z\in C_c^\infty(B(x_0,R\ell))$, $0<\ell\leq\ell_0\leq1$,
\begin{equation}\label{eq:freezing-quantitative}
 \big|\mathcal T(z,z,Qz)-\mathcal T_0(z,z,Q_0z)\big|
 \leq A\ell\|\partial z\|_{L^3(\R^4)}^3+B_\ell\|z\|_{L^3(\R^4)}\|\partial z\|_{L^3(\R^4)}^2.
\end{equation}
Here $A$ depends only on $R$, the field dimension, and finitely many
fixed smooth coefficient and normalized test-symbol seminorms.
It is independent of $\ell$, $x_0$, the number of patches, and the
frequencies of $z$. For each fixed $\ell$, $B_\ell$ is finite and
uniform in $x_0$, with the same independence from the frequencies of
$z$; it may depend on further fixed smooth seminorms and on $\ell^{-1}$.
No second derivative of $z$ is present in the estimate.
\end{lemma}

\begin{proof}
Set $R^{ab}_{JK}(z)=C_{IJK}^{ab}z^I$. We first estimate
$\mathcal E_{C,Q}(z)=\dd P_C(z)[Qz]$.
The third derivative $P_C'''(z,z,Qz)$ is exactly
$2\mathcal E_{C,Q}(z)$, so this choice changes only a fixed numerical
factor in the action form. Expanding the three test placements gives
$\partial_bQz=Q\partial_bz+[\partial_b,Q]z$.
The combined $(a,J),(b,K)$ symmetry and $Q^*=-Q$ imply
\begin{equation*}
\begin{aligned}
2\int_{\R^4}(\partial_a z)^J R^{ab}_{JK}(z)Q(\partial_b z)^K\,\dd x
&=-\int_{\R^4}(\partial_a z)^J[Q,R^{ab}_{JK}(z)](\partial_b z)^K\,\dd x.
\end{aligned}
\end{equation*}
Thus the apparent second derivative cancels before estimation:
\begin{equation*}
\begin{split}
\mathcal E_{C,Q}(z)
={}&\int_{\R^4} C_{IJK}^{ab}(Qz)^I\partial_a z^J\partial_b z^K\,\dd x\\
 &-\int_{\R^4} \partial_a z^J[Q,R^{ab}_{JK}(z)]\partial_b z^K\,\dd x\\
 &+2\int_{\R^4} R^{ab}_{JK}(z)\partial_a z^J[\partial_b,Q]z^K\,\dd x.
\end{split}
\end{equation*}

For a real odd degree-one symbol $b$ put $Q_b=-\ii b(D)$ and
$Q_{f,b}=(fQ_b+Q_bf)/2$, where $f$ is real. Let $\mathfrak m_N(b)$
be a fixed finite normalized symbol seminorm controlling the $L^3$
multiplier bound and the $(3,3,3/2)$ commutator bound proved above;
the derivative count $N$ is fixed once and for all.
Since $f$ is scalar,
\begin{equation*}
\begin{gathered}
[Q_{f,b},R]=\f12\big(f[Q_b,R]+[Q_b,R]f\big),\\
[\partial_a,Q_{f,b}]=\f12\big((\partial_af)Q_b+Q_b(\partial_af)\big).
\end{gathered}
\end{equation*}
The product rule gives
\begin{equation*}
\|\partial(Cz)\|_{L^3_x}
\leq\|C\|_{L^\infty_x}\|\partial z\|_{L^3_x}
    +\|\partial C\|_{L^\infty_x}\|z\|_{L^3_x}.
\end{equation*}
Using this estimate in the commutator identity and applying H\"older's
inequality gives
\begin{equation*}
\begin{split}
|\mathcal E_{C,Q_{f,b}}(z)|\leq c\mathfrak m_N(b)\big[&
 \|C\|_{L^{\infty}_x}\|f\|_{L^{\infty}_x} \|\partial z\|_{L^3_x}^3\\
 &+(\|C\|_{L^{\infty}_x}\|\partial f\|_{L^{\infty}_x}
       +\|\partial C\|_{L^{\infty}_x}\|f\|_{L^{\infty}_x})\|z\|_{L^3_x}\|\partial z\|_{L^3_x}^2\\
 &+\|C\|_{L^{\infty}_x}\|\partial^2f\|_{L^{\infty}_x} \|z\|_{L^3_x}^2\|\partial z\|_{L^3_x}\big].
\end{split}
\end{equation*}
Every differentiated coefficient multiplies an undifferentiated field.
Since $z$ is supported in $B(x_0,R\ell)$, Poincar\'e's inequality gives
\begin{equation*}
\|z\|_{L^3_x}\leq c_R\ell\|\partial z\|_{L^3_x}.
\end{equation*}

Insert nested cutoffs $\theta_\ell$ equal to one near $\supp z$.
Replacing $Q-Q_0$ by $\theta_\ell(Q-Q_0)\theta_\ell$ leaves its action
energy unchanged. Its principal symbol is $-\ii s_\ell$, where the
real underlying symbol is
$s_\ell=\theta_\ell^2(a(x,\xi)-a(x_0,\xi))$.
After $y=(x-x_0)/\ell$, every fixed $y$ and angular derivative of
$s_\ell(x_0+\ell y,\omega)$ is $O(\ell)$, uniformly in $x_0$:
the zeroth derivative uses the mean-value theorem, and each positive
$y$ derivative of $a$ supplies a factor $\ell$.
A real spatial Fourier expansion on a fixed $y$ cube therefore gives
$s_\ell=\sum_k f_{k,\ell}(x)b_{k,\ell}(\xi)$, with real $f_{k,\ell}$
and real odd degree-one $b_{k,\ell}$, such that for every fixed $L$
\begin{equation*}
\begin{gathered}
\mathfrak m_N(b_{k,\ell})\leq c\ell\mathcal A_{L,N}(a)(1+|k|)^{-L},\\
\|\partial^rf_{k,\ell}\|_{L^\infty_x}\leq c_r\ell^{-r}(1+|k|)^r
\quad(0\leq r\leq2).
\end{gathered}
\end{equation*}
Here $\mathcal A_{L,N}(a)$ is a finite fixed smooth-symbol seminorm;
the spatial profiles include fixed rescaled support cutoffs.
Taking $L>7$ makes the following sums uniform in every truncation:
\begin{equation*}
\begin{gathered}
\sum_k\mathfrak m_N(b_{k,\ell})\|f_{k,\ell}\|_{L^\infty_x}
\leq c\ell\mathcal A_{L,N}(a),\\
\sum_k\mathfrak m_N(b_{k,\ell})
 (\|\partial f_{k,\ell}\|_{L^\infty_x}+\ell\|\partial^2f_{k,\ell}\|_{L^\infty_x})
 \leq c\mathcal A_{L,N}(a).
\end{gathered}
\end{equation*}
The sum of the skew-adjoint operators $Q_{f_{k,\ell},b_{k,\ell}}$
has principal symbol $-\ii s_\ell$. Its difference from
$\theta_\ell(Q-Q_0)\theta_\ell$ is order zero by smooth symbol
calculus. At each fixed $\ell$ its $L^3$ and $W^{1,3}$ operator norms
are uniformly bounded in $x_0$. An order-zero test $R_0$ contributes
at most $c_\ell \|z\|_{L^3_x}\|\partial z\|_{L^3_x}^2$: expand the three action placements and use
$\|R_0z\|_{L^3_x}\leq c_\ell \|z\|_{L^3_x}$ and
$\|\partial R_0z\|_{L^3_x}\leq c_\ell(\|\partial z\|_{L^3_x}+\|z\|_{L^3_x})$. This also treats the fixed infrared
and quantization remainders. Thus the weighted sums above bound
test freezing by $c\ell\|C\|_{L^{\infty}_x}\mathcal A_{L,N}(a)\|\partial z\|_{L^3_x}^3+c_\ell \|z\|_{L^3_x}\|\partial z\|_{L^3_x}^2$.

For coefficient freezing, extend $C-C(x_0)$ with a nested cutoff.
Its supremum is at most $c\ell\|\partial C\|_{L^{\infty}_x}$, and its first derivative
is at most $c\|\partial C\|_{L^{\infty}_x}$. The separated estimate with $f=1$ and
$b=a(x_0,\cdot)$ applies. In particular one may take the leading
constant to be a fixed multiple of
\begin{equation*}
\|C\|_{L^{\infty}_x}\mathcal A_{L,N}(a)
+\|\partial C\|_{L^{\infty}_x}\sup_{x_0}\mathfrak m_N(a(x_0,\cdot)).
\end{equation*}
Finally, lower action terms containing $z^2\partial(Qz)$ are integrated once
by parts; all resulting terms have an undifferentiated $z$ and obey
the $B_\ell \|z\|_{L^3_x}\|\partial z\|_{L^3_x}^2$ bound. All their coefficients belong to the fixed
smooth background. Collecting the coefficient and test differences gives
\begin{equation*}
\begin{aligned}
\big|\mathcal T(z,z,Qz)-\mathcal T_0(z,z,Q_0z)\big|
&\leq c\sum_k|\mathcal E_{C,Q_{f_{k,\ell},b_{k,\ell}}}(z)|
       +c|\mathcal E_{C-C(x_0),Q_0}(z)|
       +c_\ell\|z\|_{L^3_x}\|\partial z\|_{L^3_x}^2\\
&\leq A\ell\|\partial z\|_{L^3_x}^3
       +B_\ell\|z\|_{L^3_x}\|\partial z\|_{L^3_x}^2.
\end{aligned}
\end{equation*}
This proves \eqref{eq:freezing-quantitative}.
\end{proof}

\subsection{Angular localization and normalized trilinear symbols}

Fix once and for all $0<\kappa<1/16$, which separates low/high from
comparable frequencies. Here and below a low/high boundary is understood
up to the following fixed numerical constants. Choose
$\vartheta\in C^\infty([0,\infty);[0,1])$ equal to one on $[0,1/4]$
and zero on $[1,\infty)$. For a nonzero triad $p+q+r=0$ set
$\lambda_p=|p|$, $\lambda_q=|q|$, $\lambda_r=|r|$, and define
\begin{equation}\label{eq:analytic-triad-partition}
L_p=\vartheta\!\left(\f{\lambda_p^2}{\kappa^2(\lambda_q^2+\lambda_r^2)}\right),
\quad L_q,L_r\ \hbox{cyclically},\qquad
L_{\rm comp}=1-L_p-L_q-L_r.
\end{equation}
At most one low cutoff is nonzero. Indeed $L_p,L_q\ne0$ would give
$\lambda_p^2+\lambda_q^2<\kappa^2(\lambda_p^2+\lambda_q^2+2\lambda_r^2)\leq5\kappa^2(\lambda_p^2+\lambda_q^2)$.
Thus all four cutoffs are nonnegative and sum to one. On $L_p\ne0$,
$\lambda_p/\lambda_q\leq2\kappa$ and $1-2\kappa\leq \lambda_r/\lambda_q\leq1+2\kappa$;
on $L_{\rm comp}\ne0$, $\min(\lambda_p,\lambda_q,\lambda_r)/\max(\lambda_p,\lambda_q,\lambda_r)\geq\kappa/2$.
The low piece is invariant under interchange of its high slots and
under high adjunction $(p,q,r)\mapsto(-p,-r,-q)$.

For the three-physical low pieces we also fix the angular partition
explicitly. Put
\begin{equation*}
\mathfrak d(p,q)=\f{|p\wedge q|^2}{\lambda_p^2\lambda_q^2},\qquad
N_p=\vartheta\!\left(\f{\mathfrak d(p,q)}{\delta^2}\right)
    \vartheta\!\left(\f{\mathfrak d(p,r)}{\delta^2}\right).
\end{equation*}
Use $L_pN_p$ for the near piece and $L_p(1-N_p)$ for the far piece.
These cutoffs have the same high-slot symmetries. On the far support,
at least one of the two squared projective separations is at least
$\delta^2/4$; the identity
$\mathfrak d(p,r)=(\lambda_q^2/\lambda_r^2)\mathfrak d(p,q)$ shows that both are
bounded below by a fixed multiple of $\delta^2$.

For $0<\delta<\delta_0$ set
$\eta=c\kappa\delta^2$, with $c>0$ sufficiently small and fixed.
Cover the projective null directions by $N_\delta\leq C\delta^{-2}$ caps
of radius $\delta$, with bounded overlap, using both signed lifts to the
two sheets. Take smooth multipliers $P_\gamma$ supported in these caps
and in $|q_m(\xi)|\leq C\eta|\xi|^2$, forming a partition on a smaller
cone neighborhood. These multipliers are constructed from fixed smooth
profiles in the parabolically rescaled cap coordinates, with uniform
bounds for every required normalized derivative. Bounded enlargements
$P_\gamma^+$ will always be used on fields when a symbol is
supported in a cap. They equal one on that support, including its
thin-cone profile, and have bounded frequency overlap.

For clarity, the cap localization of a trilinear form uses only one
partition: insert $\sum_\gamma P_\gamma(p)=1$ in one designated slot.
On the near support, the triad geometry puts both other slots in a
fixed enlargement of that same projective cap. Choose $P_\gamma^+$
equal to one on all three resulting slot supports, retaining both
signed lifts. Thus each localized form acts on
$(P_\gamma^+f,P_\gamma^+g,P_\gamma^+w)$. A maximal $\delta$-separated
net on the two-dimensional projective null set gives
$N_\delta\leq C\delta^{-2}$ and uniformly bounded overlap even after
this fixed enlargement. In particular, no independent sum over three
cap indices is introduced.

The individual bounds, and the angular sequence estimate needed here,
are
\begin{equation}\label{eq:analytic-cap-bounds}
\begin{gathered}
\|P_\gamma f\|_{L^p_x}\leq C_p\|f\|_{L^p_x}\quad(1<p<\infty),\\
\bigg(\sum_\gamma\|P_\gamma f\|_{L^3_x}^3\bigg)^{1/3}
\leq C_{p_*} N_\delta^{1/(3(p_*-2))}\|f\|_{L^3_x}
\quad(3<p_*<\infty).
\end{gathered}
\end{equation}
The same estimates hold for $P_\gamma^+$.
To verify the uniform first bound, rotate the cap center and use null
coordinates $(s,t,z)$ in frequency space, with two screen coordinates
$z$ and $q_m=-st+|z|^2$. On the cap,
$|z|\ls\delta|s|$ and $|t|\ls\delta^2|s|$.
The invertible linear change
$(s,t,z)\mapsto(s,t/\delta^2,z/\delta)$ puts its symbol in a fixed cone
with uniform Mikhlin seminorms. Its determinant cancels under conjugation
on $L^p$. Constants may depend on the fixed $\kappa,c$ and the enlargement,
but not on $\delta$. Give the finite cap set $\Gamma_\delta$ counting
measure and regard $Tf(x,\gamma)=P_\gamma f(x)$ as a linear operator
into functions on $\R^4\times\Gamma_\delta$. Plancherel and bounded
frequency overlap give
\begin{equation*}
\|Tf\|_{L^2(\dd x\times\#)}^2
=\sum_\gamma\|P_\gamma f\|_{L^2_x}^2\leq C\|f\|_{L^2_x}^2.
\end{equation*}
For every fixed finite $p_*>3$, the individual cap estimates give
\begin{equation*}
\|Tf\|_{L^{p_*}(\dd x\times\#)}^{p_*}
=\sum_\gamma\|P_\gamma f\|_{L^{p_*}_x}^{p_*}
\leq C_{p_*}^{p_*}N_\delta\|f\|_{L^{p_*}_x}^{p_*}.
\end{equation*}
Riesz-Thorin interpolation between these two ordinary $L^p$ estimates
uses $1/3=(1-\theta_*)/2+\theta_*/p_*$, hence
$\theta_*=p_*/(3(p_*-2))$. Its counting factor is
$N_\delta^{\theta_*/p_*}=N_\delta^{1/(3(p_*-2))}$, proving
\eqref{eq:analytic-cap-bounds}. This is an $L^3(\dd x\times\#)$
estimate; no angular $\ell^2$ bound at $L^3$ or $L^\infty$ multiplier
estimate is needed. All exponents remain finite.

Consider a nonzero triad $(p,\xi,-p-\xi)$ whose three frequencies satisfy
$|q_m(\zeta)|\leq C\eta|\zeta|^2$, and choose $p$ to have least
Euclidean magnitude. Write $\lambda_p=|p|$, $\lambda_q=|\xi|$; the two high magnitudes
are comparable. The identity
\begin{equation*}
q_m(p+\xi)=q_m(p)+q_m(\xi)+2m(p,\xi)
\end{equation*}
implies, on the comparable piece $\lambda_p/\lambda_q\geq\kappa/2$, that the three projective directions are
in boundedly adjacent $\delta$ caps: their squared angular separation is
at most $C\eta/\kappa$. Indeed the cone bounds first give
$|m(p,\xi)|/(\lambda_p\lambda_q)\leq C\eta/\kappa$; on two near-null sheets the absolute
normalized Lorentz product controls squared projective angle, up to
$O(\eta)$. On the low pieces $\lambda_p/\lambda_q\leq2\kappa$, use the near/far
partition above. In the separated region
\begin{equation}\label{eq:analytic-phase-gap}
\Delta(p,\xi):=q_m(\xi+p)-q_m(\xi)
=2m(p,\xi)+q_m(p),\qquad
|\Delta(p,\xi)|\geq c_1\delta^2\lambda_p\lambda_q.
\end{equation}
Here $|m(p,\xi)|\geq c_2\delta^2\lambda_p\lambda_q$, whereas
$|q_m(p)|\leq C\eta \lambda_p^2$ is absorbable by the fixed choice of $c$.
Only the angular comparison identifies antipodal directions; signed
frequencies and the exact phase are retained in every Fourier calculation.


The following transfer statement separates the symbol calculation from
its operator estimate. ``Normalized cap variables'' means the variables
obtained from the common parabolic linear change used above, with each
radial frequency divided by its own magnitude. Smoothness in the ratio
$\rho=\lambda_p/\lambda_q$ is required down to $\rho=0$, not only on a finite range
of frequency gaps.

\begin{lemma}[All-scale transfer and angular summation]\label{lem:analytic-near}
Suppose a translation-invariant trilinear Fourier multiplier on
$p+q+r=0$ has, after factoring the three first-derivative weights,
a symbol $m$ homogeneous of degree zero under simultaneous positive
dilation of its three frequencies. On each piece of a fixed smooth
homogeneous low/high partition, require
$m=M(\omega,\theta,\rho)$, where $\omega=p/\lambda_p$, $\theta=q/\lambda_q$, and
$\rho=\lambda_p/\lambda_q$. Let $\mathfrak A_{100}$ bound its $C^{100}$ extension to
$\rho=0$ in a fixed finite angular atlas, together with the
$C^{100}$ extension norms of the comparable pieces to fixed normalized
annular boxes. Fixed factorwise infrared cutoffs are allowed separately.
For the three differentiated factors $f,g,w$, the form satisfies
\begin{equation}\label{eq:analytic-near-statement-2}
|\mathcal T_m(f,g,w)|
\leq C_\kappa\mathfrak A_{100}
\|f\|_{L^3(\R^4)}\|g\|_{L^3(\R^4)}\|w\|_{L^3(\R^4)}.
\end{equation}
The constant is uniform in
finite smooth dyadic truncations of the factors and does not depend on their
highest frequency or number of scales.

If in addition the symbol is supported on adjacent null caps and equals
$\delta^2$ times a symbol with $\mathfrak A_{100}\leq A$ uniformly in
normalized cap variables, including the ratio endpoint, then, for
every fixed $p_*>3$,
\begin{equation}\label{eq:analytic-near-bound}
|\mathcal T^{\mathrm{near}}(f,g,w)|
\leq C_{\kappa,p_*}A\delta^{\beta_*}\|f\|_{L^3(\R^4)}\|g\|_{L^3(\R^4)}\|w\|_{L^3(\R^4)}.
\end{equation}
Thin-cone multipliers are included in the individual cap symbols. The
right-hand side contains no norm of the full thin-cone projection.
\end{lemma}

\begin{proof}
We give the symbol-to-operator step. More generally, suppose a trilinear
symbol, after factoring its three first-derivative weights, is $\delta^2$
times a uniformly smooth symbol in normalized cap variables, smooth also
in $\rho=\lambda_p/\lambda_q$ down to zero on the low/high part.
We claim that the contribution of one enlarged cap satisfies
\begin{equation*}
|\mathcal T_{m,\gamma}(f,g,w)|\leq C_\kappa A\delta^2
\|P_\gamma^+f\|_{L^3_x}\|P_\gamma^+g\|_{L^3_x}\|P_\gamma^+w\|_{L^3_x}.
\end{equation*}
Indeed, the common parabolic change acts on the differentiated factors;
componentwise one may take $f=|D_x|u$, using Riesz transforms to compare
its norm with $\|\partial u\|_{L^3_x}$. Let
$\mathcal A=E_\gamma\operatorname{diag}(1,\delta^2,\delta,\delta)$
map a fixed enlarged cone to the physical cap, where $E_\gamma$ is a
uniformly bounded and invertible null frame. For $F(y)=f(\mathcal A^{-\mathsf T}y)$
and likewise $G,W$, the change of variables gives
\begin{equation*}
\mathcal T_m(f,g,w)=|\det\mathcal A|^{-1}
 \mathcal T_{m(\mathcal A\cdot,\mathcal A\cdot)}(F,G,W),
\qquad \|F\|_{L^3_x}=|\det\mathcal A|^{1/3}\|f\|_{L^3_x}.
\end{equation*}
Thus the determinant factors cancel exactly. On the enlarged fixed cone
$|s|\geq c|\zeta|$, so
$b_\delta(\zeta)=|\mathcal A\zeta|/|\zeta|$ and its reciprocal have
uniform homogeneous symbol seminorms. This is a cone-restricted
statement, not uniform invertibility of $\mathcal A$ on all frequencies.
Write $p'=\mathcal A^{-1}p$, $q'=\mathcal A^{-1}q$,
$\lambda_p'=|p'|$, $\lambda_q'=|q'|$, $\omega'=p'/\lambda_p'$, $\theta'=q'/\lambda_q'$.
If $\rho'=\lambda_p'/\lambda_q'$ is the transformed radial ratio, then
\begin{equation*}
\rho=\rho'\f{b_\delta(\omega')}{b_\delta(\theta')}.
\end{equation*}
This preserves all required endpoint and annular derivative bounds.
Numerical low/high thresholds may change by fixed factors; their
transition pieces have finitely many relative radial shifts. It remains
to prove a fixed-cone estimate for the transformed normalized symbol.
In that argument we write $M,f,g,w,\omega,\theta$ for the transformed
symbol, factors and angles, while $\rho$ retains the physical ratio.
Radial annuli may be taken in the transformed magnitudes, which are
uniformly comparable to the physical ones on this cone.

For one already partitioned low/high symbol put
$M_0(\omega,\theta)=M(\omega,\theta,0)$. Choose a fixed smooth ratio
cutoff $\chi_0$ equal to one near zero and supported in a sufficiently
small low/high region. The exact split is
\begin{equation*}
M=\chi_0(\rho)M_0+\rho R,\qquad
R=\f{M-\chi_0(\rho)M_0}{\rho}.
\end{equation*}
Near zero, $R=\int_0^1\partial_\rho M(\omega,\theta,t\rho)\,\dd t$.
Away from zero the quotient is harmless, so $R$ has a uniform
$C^{99}$ bound and fixed low/high support. Fixed transition regions
are again included among the finitely many relative annular shifts.
No vanishing of $M_0$ is assumed. A separated angular expansion of
$M_0$ has coefficients bounded
by $C\mathfrak A_{100}(1+|a|+|b|)^{-40}$; its angular multiplier norms
grow at most as $(1+|a|)^6$ and $(1+|b|)^6$. These fixed derivative
counts suffice in dimension four and are not claimed to be optimal.
There is a cutoff correction in identifying its radial paraproduct.
On $\lambda_q'\sim2^j$, choose
$S_{<j-C}$ with symbol $\psi(2^{C-j}p')$, where $\psi=1$ on the
unit ball and is supported in the ball of radius two. A fixed large
$C$ puts its support where $\chi_0(\rho)=1$, uniformly in the
two angles. Both cutoffs equal one when $\lambda_p'/\lambda_q'$ is sufficiently
small. Their difference has $c_\kappa\leq \lambda_p/\lambda_q\leq C_\kappa$, so only
finitely many relative annular shifts occur; these are estimated below.
The remaining endpoint terms are exactly paraproducts of the form
\begin{equation*}
\sum_j\int_{\R^4}
(S_{<j-C}R_1f)(\Delta_jR_2g)(\widetilde\Delta_jw)\,\dd x.
\end{equation*}
Here $\widetilde\Delta_j$ is a fixed thick annulus equal to one on
the possible product frequencies $-p'-q'$; low/high separation makes
this insertion exact. The maximal-function bound, the two square-function
bounds, and H\"older's inequality give
\begin{equation*}
\begin{aligned}
\bigg|\sum_j\int_{\R^4}
 (S_{<j-C}R_1f)(\Delta_jR_2g)(\widetilde\Delta_jw)\,\dd x\bigg|
&\leq\int_{\R^4}\sup_j|S_{<j-C}R_1f|
 \bigg(\sum_j|\Delta_jR_2g|^2\bigg)^{1/2}
 \bigg(\sum_j|\widetilde\Delta_jw|^2\bigg)^{1/2}\,\dd x\\
&\leq C\|R_1f\|_{L^3_x}\|R_2g\|_{L^3_x}\|w\|_{L^3_x}.
\end{aligned}
\end{equation*}
The angular coefficients sum the polynomial norms of $R_1,R_2$.
No separate sum over low/high gaps is applied to $M_0$.

For $\rho R$, fix a radial gap $s\geq C$. The factor $\rho$ is bounded
by $C2^{-s}$. More precisely, on transformed annuli
$p'=2^{j-s}\zeta$, $q'=2^j\upsilon$ one has
$\rho=2^{-s}|\mathcal A\zeta|/|\mathcal A\upsilon|$.
Every required annular derivative of this expression is $O(2^{-s})$;
the remaining symbol has smooth bounds uniform in $s$ and $\delta$.
Before separating the frequency variables, insert a thick product-frequency
annular cutoff equal to one on the original block support. The two
high magnitudes are comparable, so the product-frequency cutoff is a fixed
enlargement of the $j$th annulus, uniformly in $s$. Keep this cutoff
on the third factor during the following expansion.
Extend the normalized two-factor symbol to a fixed eight-dimensional
box and expand it in a Fourier series. Its coefficients obey
\begin{equation*}
|c_{s,a,b}|\leq C\mathfrak A_{100}\,2^{-s}
(1+|a|+|b|)^{-40},\qquad a,b\in\mathbb Z^4.
\end{equation*}
Each separated term consists of annular
multipliers translated at their own scales. Their square-function bounds
grow at most as $(1+|a|)^6$ or $(1+|b|)^6$, uniformly in $s$: randomize
the annular signs, apply Mikhlin using bounded annular overlap, and use
Khintchine's inequality. More explicitly, on a multiplier at scale
$2^{j-s}$, each normalized derivative of
$\exp(\ii 2^{-(j-s)}a\cdot\xi)$ costs at most $C(1+|a|)$.
Only boundedly many randomized annuli meet any nonzero frequency,
independently of $s$ and the finite index set. The required derivatives
are all controlled by $\mathfrak A_{100}$.
Use
\begin{equation*}
\sum_j|f_jg_jw_j|
\leq\bigg(\sum_j|f_j|^2\bigg)^{1/2}
     \bigg(\sum_j|g_j|^2\bigg)^{1/2}
     \bigg(\sum_j|w_j|^2\bigg)^{1/2}
\end{equation*}
and H\"older. The Fourier coefficients are absolutely summable with
these multiplier losses, since
\begin{equation*}
\sum_{a,b\in\mathbb Z^4}
(1+|a|+|b|)^{-40}(1+|a|)^6(1+|b|)^6
\leq C\sum_{n\in\mathbb Z^8}(1+|n|)^{-28}<\infty.
\end{equation*}
Let $\mathcal T_s$ be the part with transformed radial gap $s$.
The annular expansion and these square-function bounds yield
\begin{equation*}
\begin{aligned}
|\mathcal T_s(f,g,w)|
&\leq C\sum_{a,b\in\mathbb Z^4}|c_{s,a,b}|
 (1+|a|)^6(1+|b|)^6
 \|f\|_{L^3_x}\|g\|_{L^3_x}\|w\|_{L^3_x}\\
&\leq C\mathfrak A_{100}2^{-s}
 \|f\|_{L^3_x}\|g\|_{L^3_x}\|w\|_{L^3_x}.
\end{aligned}
\end{equation*}
Summing every radial gap gives
\begin{equation*}
\begin{aligned}
|\mathcal T_{\rho R}(f,g,w)|
&\leq\sum_{s\geq C}|\mathcal T_s(f,g,w)|\\
&\leq C\mathfrak A_{100}\sum_{s\geq C}2^{-s}
 \|f\|_{L^3_x}\|g\|_{L^3_x}\|w\|_{L^3_x}
 \leq C\mathfrak A_{100}\|f\|_{L^3_x}\|g\|_{L^3_x}\|w\|_{L^3_x}.
\end{aligned}
\end{equation*}
On the comparable
region there are only $O_\kappa(1)$ relative radial shifts; the same
annular expansion and three square functions apply directly. Here too
the thick product-frequency annulus must be inserted before the expansion: a
single separated Fourier term need not retain the original exclusion
of $p'+q'=0$. The third cutoff retains that exclusion and the third-factor
square-function bound term by term. These
arguments initially apply to finite radial sums and give constants
independent of their length, hence also their limits. More precisely,
finite truncations of the factors $\sum_{j\in E}\epsilon_j\Delta_j$, with
$|\epsilon_j|\leq1$, have uniformly bounded Mikhlin norms by annular
overlap and may be absorbed into each factor. No arbitrary independent
weight on every pair of dyadic indices is being allowed. They cover each
low-slot choice and the boundedly many adjacent caps and high-index shifts.

H\"older's inequality in the cap index and \eqref{eq:analytic-cap-bounds}
give the full target estimate:
\begin{equation*}
\begin{aligned}
|\mathcal T^{\mathrm{near}}(f,g,w)|
&\leq C_\kappa A\delta^2\sum_{\gamma\in\Gamma_\delta}
 \|P_\gamma^+f\|_{L^3_x}\|P_\gamma^+g\|_{L^3_x}\|P_\gamma^+w\|_{L^3_x}\\
&\leq C_\kappa A\delta^2
 \prod_{F\in\{f,g,w\}}\bigg(\sum_{\gamma\in\Gamma_\delta}
                    \|P_\gamma^+F\|_{L^3_x}^3\bigg)^{1/3}\\
&\leq C_{\kappa,p_*}A\delta^2N_\delta^{1/(p_*-2)}
       \|f\|_{L^3_x}\|g\|_{L^3_x}\|w\|_{L^3_x}\\
&\leq C_{\kappa,p_*}A\delta^{2(p_*-3)/(p_*-2)}
       \|f\|_{L^3_x}\|g\|_{L^3_x}\|w\|_{L^3_x}.
\end{aligned}
\end{equation*}
This proves \eqref{eq:analytic-near-bound}. Omitting the cap localization
and its prefactor in the radial argument proves the first assertion.
\end{proof}

\subsection{Scalar corrections and Lipschitz coefficients}

Consider a translation-invariant bilinear operator $T_Hu=T_t(H,u)$, with low-frequency
factor $H$ and high-frequency factor $u$. Assume its symbol is supported in
$|p|<C\kappa|\xi|$, with $C\kappa<1/2$, and vanishes outside this
region after a smooth
cutoff. On that support suppose its
Fourier-shift symbol has the form
\begin{equation}\label{eq:analytic-extension}
t(p,\xi)=\lambda_\xi M(\omega,z,\theta),\qquad
\lambda_\xi=|\xi|,\quad \omega=p/|p|,\quad z=p/\lambda_\xi,\quad \theta=\xi/\lambda_\xi,
\end{equation}
where $M$ extends smoothly to independent variables
$(\omega,z,\theta)\in S^3\times\mathcal K$, where $\mathcal K$ is a
fixed compact $(z,\theta)$ domain containing $\{0\}\times S^3$.
Low/high cutoffs are smooth in $|z|^2$ and high frequencies are
regularized near zero. Fixed cutoffs at low absolute frequency are
placed on $H$. This independent-variable extension is a hypothesis;
smoothness in the ratio $|p|/|\xi|$ alone does not imply it.
The symbol may depend on a fixed angular parameter $\delta$;
constants $C_\delta$ below depend on finitely many of its seminorms.

Expand only the independent variable $\omega$ into spherical harmonics:
\begin{equation}\label{eq:analytic-bs-expansion}
M(\omega,z,\theta)
=\sum_\nu Y_\nu(\omega)c_{\nu,\delta}(z,\theta).
\end{equation}
For every finite derivative order in $(z,\theta)$ and every $N$, the
coefficients have bounds $C_{N,\delta}(1+\nu)^{-N}$, with any fixed
enumeration whose eigenvalues and multiplicities grow polynomially.
Extend the low angular factors smoothly outside the supporting region and
move $Y_\nu(p/|p|)$ onto the low-frequency factor, writing $R_\nu H=Y_\nu(D)H$.
For each fixed $1<s<\infty$, the $L^s$ multiplier norms of $R_\nu$ grow
at most polynomially in $\nu$. The remaining bilinear symbol
\begin{equation*}
b_{\nu,\delta}(p,\xi)
=\lambda_\xi c_{\nu,\delta}(p/\lambda_\xi,\xi/\lambda_\xi)
\end{equation*}
with the smooth low/high cutoff belongs to $BS^1_{1,0}$ after high-frequency
infrared regularization. Any fixed infrared cutoff in the \emph{low}
frequency is placed on $H$ (and hence on $R_\nu H$), not in this
Cartesian symbol. Keeping such a cutoff inside the symbol would fail the
required low-derivative bounds when $|p|$ is fixed and $\lambda_\xi$ is large.
The smooth $z=0$ extension requires no low cutoff inside the symbol. Indeed each Cartesian $p$ derivative costs
$\lambda_\xi^{-1}$, as does each $\xi$ derivative; since $|p|\ls \lambda_\xi$,
\begin{equation*}
|\partial_p^\alpha\partial_\xi^\beta b_{\nu,\delta}|
\leq C_{N,\alpha,\beta,\delta}(1+\nu)^{-N}
       (1+|p|+|\xi|)^{1-|\alpha|-|\beta|}.
\end{equation*}
Thus \eqref{eq:analytic-bs-expansion} is an absolutely summable expansion in
every required operator seminorm. In particular it is not an assertion
that the unexpanded low angular symbol is Cartesian smooth at $p=0$.

The dependence on finitely many derivatives can be made explicit.
For a translation-invariant order-one bilinear symbol set
\begin{equation*}
 \mathfrak p_K(b)=
 \max_{|\alpha|+|\beta|\leq K}
 \sup_{p,\xi}(1+|p|+|\xi|)^{|\alpha|+|\beta|-1}
 |\partial_p^\alpha\partial_\xi^\beta b(p,\xi)|.
\end{equation*}
Fix a finite derivative threshold $K$ large enough for the
order-zero Coifman-Meyer estimate and the two exponent triples
$(4,4,2)$ and $(4,2,4/3)$ in \cite[Theorem~1]{BenyiOh}.
That theorem's commutator norm is bounded by
$C\mathfrak p_K(b)\|\partial d\|_{L^{\infty}_x}$.
Here is a functional-analytic justification of the uniform finite
threshold, without assuming a quantitative version of that theorem.
Let $X$ be the Fr\'echet space of translation-invariant $BS^1_{1,0}$
symbols, $Y$ the Banach space of Lipschitz functions modulo constants,
and $Z_r=\mathcal B(L^4(\R^4)\times L^r(\R^4),L^s(\R^4))$, where
$(r,s)=(4,2)$ or $(2,4/3)$. Write
$C(b,d)(f,g)=B_b(f,dg)-dB_b(f,g)$.
The cited theorem supplies an element of $Z_r$ for every $(b,d)$.
To verify its distributional continuity, set
$\Lambda^2=1+|p|^2+|\xi|^2$ and use the exact decomposition
\begin{equation*}
b=b_0+\sum_{j=0}^3p_jb_j+\sum_{j=0}^3\xi_jc_j,
\qquad
b_0=\f b{\Lambda^2},\quad
b_j=\f{p_jb}{\Lambda^2},\quad
c_j=\f{\xi_jb}{\Lambda^2}.
\end{equation*}
All these coefficient symbols belong to $BS^0_{1,0}$. The ordinary
Coifman-Meyer estimate therefore bounds $B_b$ by one derivative on
either factor, using a fixed finite symbol seminorm. Choose $d(0)=0$.
Then $|d(x)|\leq\|d\|_Y|x|$, so $dg\in W^{1,r}$ and
$d\phi\in L^{s'}$ for Schwartz $g,\phi$.
Consequently the pairing of $C(b,d)(f,g)$ with any Schwartz $\phi$
depends continuously on $b\in X$ and $d\in Y$; its bound uses
$\mathfrak p_{K_0}(b)\|d\|_Y$ and fixed weighted Schwartz seminorms
of the three factors. This identifies every graph limit in $Z_r$.
The closed graph theorem now makes $b\mapsto C(b,d)$ and
$d\mapsto C(b,d)$ continuous separately. The latter continuity
makes $\{b\mapsto C(b,d):\|d\|_Y\leq1\}$ pointwise bounded;
uniform boundedness on the Fr\'echet space $X$ makes this family
equicontinuous. Hence one finite $\mathfrak p_K$ bounds the whole
family. Take the larger threshold for the two exponent triples and
for the order-zero estimate. This proves the asserted uniform bound.

Index real orthonormal harmonics on the round $S^3$ by degree
$l\geq0$ and $1\leq k\leq(l+1)^2$, with eigenvalue $l(l+2)$.
Their coefficients are
\begin{equation*}
c_{l,k,\delta}(z,\theta)
=\int_{S^3}M(\omega,z,\theta)Y_{l,k}(\omega)\,\dd S(\omega).
\end{equation*}
Integration by parts with $(1-\Delta_{S^3})^6$ gives
\begin{equation*}
 \|c_{l,k,\delta}\|_{C^K(\mathcal K)}
 \leq C(1+l)^{-12}\|M\|_{C^{K+12}(S^3\times\mathcal K)}.
\end{equation*}
The Mikhlin theorem and $H^7(S^3)\hookrightarrow C^5(S^3)$ give
\begin{equation*}
 \|Y_{l,k}(D)\|_{L^s(\R^4)\to L^s(\R^4)}\leq C_s(1+l)^7,
 \qquad 1<s<\infty.
\end{equation*}
Combining these estimates with the Cartesian chain rule bounds the
full operator constant by
\begin{equation*}
 C\|M\|_{C^{K+12}(S^3\times\mathcal K)}
 \sum_{l\geq0}(l+1)^2(l+1)^7(l+1)^{-12}<\infty.
\end{equation*}
In particular one fixed finite norm controls the full series and all
of its angular partial sums. Truncating at harmonic degree $L$ leaves
an operator-norm tail at most $C(1+L)^{-2}\|M\|_{C^{K+12}(S^3\times\mathcal K)}$ in each
of the spaces used below. Smooth upper-frequency truncations at
scale $\lambda_{\max}$ may be placed symmetrically on the high-frequency
factor and product. Their differentiated supports have
$|\xi|\asymp\lambda_{\max}$, so the same seminorms are uniform above
the fixed infrared scale; convergence of these truncations is proved below.


\begin{lemma}[Bounds for the correction and its rough commutator]
\label{lem:analytic-correction-bounds}
An operator satisfying \eqref{eq:analytic-extension} obeys
\begin{equation}\label{eq:analytic-correction-bounds}
\begin{gathered}
\|T_Hu\|_{L^2(\R^4)}\leq C_\delta\|H\|_{L^4(\R^4)}\|\partial u\|_{L^4(\R^4)},\\
\|[\partial_a,T_H]u\|_{L^2(\R^4)}
\leq C_\delta\|\partial_aH\|_{L^4(\R^4)}\|\partial u\|_{L^4(\R^4)},\\
\|[d,T_H]v\|_{L^2(\R^4)}
\leq C_\delta\|\partial d\|_{L^{\infty}(\R^4)}\|H\|_{L^4(\R^4)}\|v\|_{L^4(\R^4)}
\end{gathered}
\end{equation}
for real scalar Lipschitz $d$. The analogous bound for the formal adjoint
pieces holds. Inhomogeneous variants add only lower-order norms of $u$.
\end{lemma}

\begin{proof}
For the first estimate divide each high symbol in
\eqref{eq:analytic-bs-expansion} by the high-frequency weight $|\xi|$. Its support is low/high
and its order-zero Cartesian derivatives have the Coifman-Meyer bounds.
The low/high paraproduct theorem gives $L^4\times L^4\to L^2$ after
putting $|D|u$ in the high-frequency factor. Equivalently, the radial proof bounds a
lowpass maximal function times a high square function, using Plancherel
for the annular product and H\"older with exponents $4,4$. Smooth normalized
annular expansions handle its ratio dependence. The two high frequencies
remain comparable on taking the adjoint, so its high derivative can also
be divided by the high-frequency weight. The harmonic coefficient bounds therefore give
\begin{equation*}
\begin{aligned}
\|T_Hu\|_{L^2_x}
&\leq\sum_{l=0}^\infty\sum_{k=1}^{(l+1)^2}
       \|B_{b_{l,k,\delta}}(Y_{l,k}(D)H,u)\|_{L^2_x}\\
&\leq C\|M\|_{C^{K+12}(S^3\times\mathcal K)}
       \sum_{l=0}^\infty(l+1)^2(1+l)^{-12}(1+l)^7
       \|H\|_{L^4_x}\|\partial u\|_{L^4_x}\\
&\leq C_\delta\|H\|_{L^4_x}\|\partial u\|_{L^4_x}.
\end{aligned}
\end{equation*}
Translation invariance gives $[\partial_a,T_H]=T_{\partial_aH}$, hence
\begin{equation*}
\|[\partial_a,T_H]u\|_{L^2_x}
=\|T_{\partial_aH}u\|_{L^2_x}
\leq C_\delta\|\partial_aH\|_{L^4_x}\|\partial u\|_{L^4_x}.
\end{equation*}

For the last estimate apply \cite[Theorem 1]{BenyiOh}: a symbol in
$BS^1_{1,0}$, commuted with a Lipschitz multiplier in either bilinear factor,
defines a bilinear Calder\'on-Zygmund operator. Its finite-exponent
boundedness includes $L^4\times L^4\to L^2$, with constant controlled by
the Lipschitz seminorm and finitely many symbol seminorms. Apply it in
the high-frequency factor of each $b_{\nu,\delta}(R_\nu H,\cdot)$. Rapid coefficient
decay sums the polynomial multiplier norms of $R_\nu$:
\begin{equation*}
\begin{aligned}
\|[d,T_H]v\|_{L^2_x}
&\leq\sum_\nu\|[d,B_{b_{\nu,\delta}}(R_\nu H,\cdot)]v\|_{L^2_x}\\
&\leq C\|\partial d\|_{L^\infty_x}\|v\|_{L^4_x}
       \sum_\nu\mathfrak p_K(b_{\nu,\delta})\|R_\nu H\|_{L^4_x}\\
&\leq C_\delta\|\partial d\|_{L^\infty_x}\|H\|_{L^4_x}\|v\|_{L^4_x}.
\end{aligned}
\end{equation*}
For clarity, if $A_Fu=B_b(F,u)$, its high adjoint has symbol
\begin{equation*}
 b^\dagger(p,\xi)=\overline{b(-p,\xi+p)}^{\mathsf T},
 \qquad A_F^*u=B_{b^\dagger}(\overline F,u).
\end{equation*}
On its support the high-frequency factor and product remain comparable.
Consequently $b^\dagger/|\xi|$ has the same finite order-zero bounds,
and the low harmonic is only replaced by its conjugate at $-\omega$.
Thus the first estimate also holds for the high adjoint. Moreover,
\begin{equation*}
 [\partial_a,T_H]=T_{\partial_aH},\qquad
 [\partial_a,T_H^*]=[\partial_a,T_H]^*.
\end{equation*}
For the rough commutator, $d$ is real scalar and
$[d,T_H^*]=-[d,T_H]^*$. The second exponent triple and duality give
\begin{equation*}
\begin{aligned}
\|[d,T_H^*]v\|_{L^2_x}
&=\sup_{\substack{w\in\mathcal S(\R^4)\\\|w\|_{L^2_x}=1}}
       |\langle[d,T_H^*]v,w\rangle|\\
&\leq\sup_{\|w\|_{L^2_x}=1}
       \|v\|_{L^4_x}\|[d,T_H]w\|_{L^{4/3}_x}\\
&\leq C_\delta\|\partial d\|_{L^\infty_x}
       \|H\|_{L^4_x}\|v\|_{L^4_x}.
\end{aligned}
\end{equation*}
The multiplier $d$
occurs only in the high slot, so no commutator with a low angular
transform is introduced. If an exact energy identity is first proved
for a finite angular sum, use its Hermitian part; the bounds are
uniform and the full Hermitian correction is recovered in operator
norm.

\end{proof}

For real symmetric scalar Lipschitz coefficients $d^{ab}=d^{ba}$, set
$L_d=-\partial_a(d^{ab}\partial_b)$.
Its exact operator commutator is
\begin{equation*}
[L_d,T_H]=
-\partial_a\circ d^{ab}\circ[\partial_b,T_H]
-\partial_a\circ[d^{ab},T_H]\circ\partial_b
-[\partial_a,T_H]\circ d^{ab}\circ\partial_b.
\end{equation*}
Move the outer derivative in the first two terms and use the adjoint
bound in the third. Lemma~\ref{lem:analytic-correction-bounds} gives
\begin{equation}\label{eq:analytic-quartic-bound}
\begin{aligned}
|\ip{\ii [L_d,T_H]u}{u}|
&\leq\sum_{a,b}\bigg(
 \|d\|_{L^\infty_x}\|[\partial_b,T_H]u\|_{L^2_x}\|\partial_au\|_{L^2_x}
 +\|[d^{ab},T_H]\partial_bu\|_{L^2_x}\|\partial_au\|_{L^2_x}\\
&\hspace{35mm}
 +\|d\|_{L^\infty_x}\|\partial_bu\|_{L^2_x}
                         \|[\partial_a,T_H]^*u\|_{L^2_x}\bigg)\\
&\leq C_\delta\bigg(
 \|d\|_{L^\infty_x}\|\partial u\|_{L^2_x}
                      \|\partial H\|_{L^4_x}\|\partial u\|_{L^4_x}\\
&\hspace{23mm}
 +\|\partial d\|_{L^\infty_x}\|\partial u\|_{L^2_x}
                      \|H\|_{L^4_x}\|\partial u\|_{L^4_x}\bigg).
\end{aligned}
\end{equation}
No second derivative of $d$ or $u$ occurs in this estimate.

The relevant energy identity is exact. If $T_H=T_H^*$ and
$P_0=-m^{ab}\partial_{ab}$, define $S_H=\ii [P_0,T_H]$. Its Fourier-shift
symbol is $\ii \Delta t$, where
$\Delta=q_m(\xi+p)-q_m(\xi)$. Conversely, a Hermitian numerator
$\sigma$ can be divided by $\ii \Delta$ on a region where $\Delta\ne0$,
provided the resulting quotient has the extension
\eqref{eq:analytic-extension}. Thus
\begin{equation}\label{eq:analytic-exact-correction}
\ii [P_0,T_H]=S_H.
\end{equation}
For $P_n=P_0+L_d$, formal Hermitian symmetry gives
\begin{equation}\label{eq:analytic-far-pairing}
\langle S_Hu,u\rangle
=2\Im\langle P_nu,T_Hu\rangle
-\langle \ii [L_d,T_H]u,u\rangle.
\end{equation}
Consequently an $L^2$ bound for $P_nu$, smallness of $H$ in $L^4$,
and the commutator estimate \eqref{eq:analytic-quartic-bound} control this
correction using only first derivatives of $d$ and of the fields.
The application below constructs two such corrections and verifies
their independent-variable extensions separately.

There is no operator-domain approximation implicit in this application.
For each fixed $n$ and each fixed localization scale, the original field
$z=\chi h_n$ is smooth and compactly supported. Its projected physical
and gauge fields are Schwartz: their complete Fourier symbols are smooth
of order zero and vanish near frequency zero. The same holds after the
fixed screen and cone multipliers are applied. A low angular multiplier
also preserves this property because the low-frequency factor retains its fixed
infrared cutoff. The coefficient $d$ used below is smooth and compactly
supported. The low/high support and polynomial bounds on symbol
derivatives also show that $T_Hu$ is Schwartz. Hence the Fourier identities, formal adjoints, and integrations
by parts above are valid on these fields for each $n$. Their Schwartz
seminorms need not be uniform in $n$; only the norms on the right sides
of the proved estimates are used when passing to the sequence limit.
In particular, let $A_L$ be the spherical-harmonic sum through degree
$L$ in \eqref{eq:analytic-bs-expansion}, and set
$T_L=(A_L+A_L^*)/2$. The independent-variable extension need not obey
the high-adjoint identity away from the actual frequency graph, so
$A_L$ itself need not be Hermitian. For the full Hermitian $T_H$,
the tail estimate and its adjoint version give
\begin{equation*}
\|(T_L-T_H)u\|_{L^2_x}
\leq C_\delta(1+L)^{-2}\|H\|_{L^4_x}\|\partial u\|_{L^4_x},
\end{equation*}
with the same tail factor in the derivative and rough-commutator
bounds. At fixed $n$, $P_0u\in L^2(\R^4)$, and Hermitian symmetry gives
\begin{equation*}
\big|\langle \ii [P_0,T_L-T_H]u,u\rangle\big|
\leq2\|P_0u\|_{L^2_x}\|(T_L-T_H)u\|_{L^2_x}\longrightarrow0.
\end{equation*}
This removes angular partial sums in operator norm at each fixed $n$.
The same angular-tail bound applies to all three rough-commutator pairings.

For upper-frequency truncation, choose real even $\chi\in C_c^\infty(\R^4)$,
equal to one near zero, and put
\begin{equation*}
J_{\lambda_{\max}}=\chi(D/\lambda_{\max}),\qquad
T_{L,\lambda_{\max}}=J_{\lambda_{\max}}T_LJ_{\lambda_{\max}}.
\end{equation*}
For fixed $n,L$, the Schwartz fields and the correction bound give
\begin{equation*}
\begin{aligned}
\|(T_{L,\lambda_{\max}}-T_L)u\|_{L^2_x}
&\leq\|J_{\lambda_{\max}}T_L(J_{\lambda_{\max}}-I)u\|_{L^2_x}
       +\|(J_{\lambda_{\max}}-I)T_Lu\|_{L^2_x}\\
&\leq C_\delta\|H\|_{L^4_x}
       \|\partial(J_{\lambda_{\max}}-I)u\|_{L^4_x}
       +\|(J_{\lambda_{\max}}-I)T_Lu\|_{L^2_x}
 \longrightarrow0.
\end{aligned}
\end{equation*}
The same argument applies to $\partial_aH$, $\partial_bu$ and
$d\partial_bu$, giving convergence of the three rough-commutator pairings.
Since $P_0u\in L^2(\R^4)$ at fixed $n$, Hermitian symmetry gives
\begin{equation*}
\big|\langle \ii [P_0,T_{L,\lambda_{\max}}-T_L]u,u\rangle\big|
\leq2\|P_0u\|_{L^2_x}
        \|(T_{L,\lambda_{\max}}-T_L)u\|_{L^2_x}
\longrightarrow0.
\end{equation*}
Thus upper-frequency cutoffs converge strongly at fixed $n,L$;
then $L\to\infty$ recovers the full correction. These limits precede
$n\to\infty$. No uniform operator-norm tail in $\lambda_{\max}$,
uniform bound for $P_0u_n$, or sequence-dependent cutoff is required.

\section{A bounded physical source in arbitrary coordinates}\label{sec:general-projection}

The correction identity \eqref{eq:analytic-far-pairing} requires a scalar
wave equation with a bounded source. We obtain it by removing the
symmetrized-gradient directions through a fixed Fourier projection.
This preserves the physical angular measure while retaining the original
Lipschitz inverse-metric coefficients in the wave operator.

\begin{proposition}[A projected source without a coordinate condition]
\label{prop:general-projection}
Under \eqref{eq:admissible}, fix a chart $U\Subset\R^4$ and
$\chi\in C_c^\infty(U)$. Extend $\chi h_n$ by zero to $\R^4$. There is a fixed real even order-zero matrix Fourier multiplier
$\Pi$ such that $u_n=\Pi(\chi h_n)$ has the following properties.
For every fixed $1<p<\infty$, $\partial u_n$ is bounded in $L^p(\R^4)$ and
$u_n\to0$ in $L^p(\R^4)$ and $L^\infty(\R^4)$. If $a_n^{ab}$ is a smooth uniformly
Lipschitz extension to $\R^4$ of $g_n^{ab}$ from a neighborhood of $\supp\chi$, then
\begin{equation}\label{eq:general-bounded-source}
\big\|-\partial_a(a_n^{ab}\partial_bu_n)\big\|_{L^2(\R^4)}\leq C_\chi.
\end{equation}
The constant may depend on the uniform Lipschitz bound of the chosen
extension family; no bound on its second derivatives is required.
Over the region where $\chi=1$, the derivative defect of $u_n$ is null,
Lorenz polarized, and has exactly the physical scalar measure
\eqref{eq:mu-invariant}.
\end{proposition}

\begin{proof}
In local Fourier calculations, $x\in\R^4$ and
$L^p_x=L^p(\R^4,\dd x)$; $H^1_x$ and $W^{1,\infty}_x$
have the same domain. Use an auxiliary Euclidean inner product on coordinate components, and
for $\xi\ne0$ set
\begin{equation*}
\begin{gathered}
P_\xi=I-\f{\xi\xi^{\mathsf T}}{|\xi|_E^2},\\
\Pi_E(\xi)H=P_\xi H P_\xi.
\end{gathered}
\end{equation*}
This is the orthogonal projection onto the complement of
$\{\xi\otimes v+v\otimes\xi:v\in\R^4\}$. Indeed
\begin{equation*}
H-P_\xi H P_\xi=\xi\otimes v+v\otimes\xi,\qquad
v=\f{H\xi}{|\xi|_E^2}
 -\f{\xi^{\mathsf T}H\xi}{2|\xi|_E^4}\xi.
\end{equation*}
Multiply the symbol by a fixed smooth radial cutoff vanishing near zero
and equal to one at high frequencies, obtaining $\Pi$. Even after this
regularization,
\begin{equation*}
\Pi(\partial_\mu V_\nu+\partial_\nu V_\mu)=0
\end{equation*}
for every compactly supported distributional covector $V$.

Define $H^\rho(k)=k^{ab}\Gamma^\rho_{ab}(k)$. The coordinate Ricci
formula, without a gauge condition, is
\begin{equation*}
2\Ric(k)_{\mu\nu}
=-k^{ab}\partial_{ab}k_{\mu\nu}
 +2k_{\rho(\mu}\partial_{\nu)}H^\rho(k)
 +H^\rho(k)\partial_\rho k_{\mu\nu}
 +\mathcal Q_{\mu\nu}(k)(\partial k,\partial k).
\end{equation*}
Here $\mathcal Q(k)$ is a smooth coefficient bilinear form in first
derivatives, obtained by collecting the Christoffel-symbol products;
its detailed coefficients are not needed for this estimate.
The potentially unbounded gauge term can be rewritten as
\begin{equation*}
2k_{\rho(\mu}\partial_{\nu)}H^\rho(k)
=2\partial_{(\mu}\big(k_{\nu)\rho}H^\rho(k)\big)
 -2(\partial_{(\mu}k_{\nu)\rho})H^\rho(k).
\end{equation*}
Here $H(g_n)$ is bounded because it uses only first derivatives.
Multiply the vacuum equation by $\chi$, subtract the smooth-background
second derivatives, and move $\chi$ inside the symmetrized gauge gradient.
The resulting derivative is killed exactly by $\Pi$; the cutoff error is
bounded. Every other term is a bounded product of first derivatives, smooth
background derivatives, inverse metrics, or fixed cutoff derivatives.
Writing $z_n=\chi h_n$ and $Q_n=a_n^{ab}\partial_{ab}$ therefore gives
\begin{equation*}
\|\Pi Q_nz_n\|_{L^2_x}\leq C_\chi.
\end{equation*}
This is not a bound for $Q_nz_n$ itself.

For a scalar Lipschitz coefficient $f$ and an $L^2(\R^4)$ field $w$, the exact
identity
\begin{equation*}
[\Pi,f]\partial_bw=[\Pi\partial_b,f]w-\Pi((\partial_bf)w)
\end{equation*}
and the first-order Calder\'on commutator bound imply
\begin{equation*}
\|[\Pi,f]\partial_bw\|_{L^2_x}
\leq C\|\partial f\|_{L^\infty_x}\|w\|_{L^2_x}.
\end{equation*}
Apply this with $w=\partial_a z_n$ to $[Q_n,\Pi]z_n$.
In particular the full global source is given explicitly by
\begin{equation}\label{eq:general-explicit-source}
-\partial_a(a_n^{ab}\partial_bu_n)
=-\Pi Q_nz_n+[\Pi,a_n^{ab}]\partial_{ab}z_n
 -(\partial_aa_n^{ab})\partial_bu_n.
\end{equation}
Each term on the right is bounded in $L^2(\R^4)$, proving
\eqref{eq:general-bounded-source}. The second derivatives of $z_n$
occur only inside the controlled commutator. These are global estimates
on the localized chart extended to Euclidean space; the source includes
the noncompact tails of $\Pi z_n$.

Finite-$L^p$ multiplier bounds give the stated norms. Uniform convergence
also follows: $z_n\to0$ in every fixed noninteger $C^\alpha$, $0<\alpha<1$,
by interpolation between uniform smallness and the Lipschitz bound, and
$\Pi$ is bounded on these spaces. No $L^\infty$ bound for $\partial u_n$
is asserted.

The original defect takes values in $\ker\mathcal R_\xi$.
Subtracting a gauge tensor preserves that kernel. Off the null cone the
kernel is entirely gauge and $\Pi_E$ kills it. On the cone its image
remains Lorenz polarized. Finally \eqref{eq:B-null} and its gauge cross
identity show that the physical contraction is unchanged by this
subtraction. The equality is an equality of full angular measures,
not merely of second moments.
\end{proof}

The bounded source in \eqref{eq:general-bounded-source} still contains
quadratic terms in the original $h_n$ and the commutator with $\Pi$.
The removed gauge field can have a non-null defect; it cannot be discarded
by strong convergence of its derivative.
The variational identity in the next section accounts for that field
and cancels the complete nonlinear energy, including all its frequency
placements.

\section{The variational transport identity and cubic cancellation}
\label{sec:variational-transport}

We derive the transport residual for the global physical measure and
analyze its frozen cubic symbol using the density-valued variational
equation and its complete cubic symmetry. The proof is completed in Section~\ref{sec:completion}. The scalar source needed for the nonlinear
correction is supplied by Proposition~\ref{prop:general-projection}; no
coordinate normalization is constructed or assumed.

\begin{theorem}[Transport under the original first-derivative hypotheses]
\label{thm:variational-transport}
Under \eqref{eq:admissible}, the invariant nonnegative measure of
Theorem~\ref{thm:quadratic} satisfies \eqref{eq:transport} for every smooth
positively one-homogeneous test with compact base support. 
No coordinate condition and no second-derivative scale are assumed.
\end{theorem}

Fix a coordinate chart $U\Subset\R^4$. Localized fields are extended
by zero to $\R^4$, and the subscripts $x$ use $\R^4$ with Lebesgue measure. Use lower metric
components as affine field coordinates. Pairings use fixed Lebesgue density
and ordered symmetric-index summation, so that off-diagonal components have
their usual multiplicity two. Equivalently one can use an orthonormal basis
of the space of symmetric tensors. A partition of the compact base support
of the test gives the global assertion after the local proof.

\subsection{The density Hessian and the exact cubic residual}

For a smooth metric $k$ on $U$ and variations in
$C_c^\infty(U;\operatorname{Sym}^2\R^4)$, set
\begin{equation*}
\mathcal S(k)=-\int_U R(k)\sqrt{|\det k|}\,\dd x,\qquad
E(k)^{ab}=\sqrt{|\det k|}\,G(k)^{ab}.
\end{equation*}
The sign is chosen so that $\mathcal S'(k)=E(k)$ for variations of the
\emph{lower} metric components. Boundary terms are immaterial for these
variations. The Hessian $E'_g$ is formally self-adjoint for every smooth
background, including a background that is not vacuum. This follows by
commuting two affine variations of the action. One may equivalently use
its local first-order Christoffel-symbol Lagrangian, modulo a divergence.

Put $w=\sqrt{|\det g|}$. The Hamiltonian satisfies
$2P=g^{ab}\xi_a\xi_b$. Define
\begin{equation*}
\mathcal G_\xi(k)_b=\xi^ak_{ab}-\f12\tr_gk\,\xi_b,\qquad
\mathsf C_g(k,l)=\f14(k_{ab}\overline l^{ab}
                       -\f12\tr_gk\,\overline{\tr_gl}).
\end{equation*}
The Fourier principal symbol $A(x,\xi)$ of $E'_g$, with
$\partial\mapsto \ii \xi$, is the trace reversal and raising of the Ricci
symbol:
\begin{equation*}
(Ak)^{ab}=w\bigl(\mathcal R_\xi(k)^{ab}
                  -\f12g^{ab}\tr_g\mathcal R_\xi(k)\bigr).
\end{equation*}
Direct contraction of \eqref{eq:ricci-symbol} gives the exact Hermitian
identity
\begin{equation}\label{eq:var-hessian-form}
k\cdot Ak
=w\bigl(4P\mathsf C_g(k,k)
        -g^{ab}\mathcal G_\xi(k)_a
                    \overline{\mathcal G_\xi(k)_b}\bigr).
\end{equation}
For a frozen metric $m$, write $\mathsf C_m$ and $A_m$ for these
objects with $g$ replaced by $m$; all raising operations and traces
then use $m$. The same notation denotes the associated map
\begin{equation*}
(\mathsf C_m J)^{ab}=\f14
\bigl(J^{ab}-\f12m^{ab}\tr_mJ\bigr).
\end{equation*}
The second term in \eqref{eq:var-hessian-form} is a Lorentzian
contraction, with no positivity assertion.
Varying density and raising factors in $E$ only produces order-zero terms
times background curvature, so none is missing from this principal symbol.

On the null Ricci kernel, $P=0$ and $\mathcal G_\xi(k)=0$.
Differentiate \eqref{eq:var-hessian-form} in any phase-space direction,
holding the coordinate components of $k$ fixed. All differentiated terms
except $4w(\dd P)\mathsf C_g(k,k)$ vanish. Thus
\begin{equation*}
k\cdot\{A,a\}k=4w H_Pa\,\mathsf C_g(k,k)
\quad(P=0,\ \mathcal R_\xi(k)=0).
\end{equation*}
Polarization gives this identity on the whole Hermitian null kernel,
including physical/gauge mixed terms. Off the cone the Ricci kernel
consists of gauge tensors $K_\xi(v)=\xi\otimes v+v\otimes\xi$.
The identity $A(x,\xi)K_\xi(v)=0$ holds for every $(x,\xi)$.
Differentiate it and pair with another vector in this gauge kernel;
self-adjointness of $A$ removes the differentiated gauge vector. Hence
the compression of $\{A,a\}$ to the off-null Ricci kernel is zero.
This accounts for the full raw defect, including its off-null part.

Choose nested real compact cutoffs $\theta,\rho$ with $\theta=1$
near the base support of $a$ and $\rho=1$ near $\supp\theta$.
Take $B_a=\theta\Op^w(a)\theta$, with a fixed infrared regularization.
It is scalar and self-adjoint, has principal symbol $a$, and is scalar
at every symbol order. Replacing $h_n$ by $\rho h_n$ in the following
pairings is exact: $B_a(\rho h_n)=B_ah_n$, and the local differential
expressions agree near the support of $B_ah_n$. This also applies to
the cubic action pairing. We suppress this outer compact cutoff below.
Then
\begin{equation*}
2\Im\ip{E'_gh_n}{B_ah_n}
=\ip{\ii [E'_g,B_a]h_n}{h_n}.
\end{equation*}
The order-three product symbols commute because $B_a$ is scalar, and
the order-two symbol is $\{A,a\}$. The order-one part of $E'_g$, and
every scalar lower-order part of $B_a$, contributes at most order one
to this commutator; its pairing tends to zero by $h_n\to0$ in $L^2(\R^4)$
and boundedness in $H^1(\R^4)$. Thus no subprincipal propagation term is omitted.
The defect with Lebesgue density is $w^{-1}\sigma$ for the
$\dd V_g$ normalization, and consequently
\begin{equation}\label{eq:var-commutator}
\lim_n2\Im\ip{E'_gh_n}{B_ah_n}=4\int_{S_r^*U} H_Pa\,\dd\mu.
\end{equation}
The factor four follows from \eqref{eq:var-hessian-form}; it differs
from the factor two in the reduced scalar-wave identity.

We justify the nonlinear expansion in this pairing. Choose a fixed
Euclidean orthonormal basis of symmetric tensors. In its field coordinates
$I,J\in\{1,\ldots,10\}$, with coordinate indices $a,b\in\{0,1,2,3\}$,
a first-order representative of the action is
\begin{equation*}
\mathcal S(k)=\f12\int_U C_{IJ}^{ab}(k)
                       \partial_ak^I\partial_bk^J\,\dd x,
\qquad C_{IJ}^{ab}=C_{JI}^{ba}.
\end{equation*}
Its Euler expression is
\begin{equation*}
E_I(k)=-\partial_a(C_{IJ}^{ab}(k)\partial_bk^J)
+\f12\f{\partial C_{JK}^{ab}}{\partial k^I}(k)
                    \partial_ak^J\partial_bk^K.
\end{equation*}
Taylor expansion of these smooth coefficient functions gives
\begin{equation*}
E(g+h_n)=E(g)+E'_gh_n+\f12E''_g[h_n,h_n]+R_n.
\end{equation*}
Separate the full remainder from its principal coefficient matrices:
write $C^{ab}=(C_{IJ}^{ab})$ and
\begin{equation*}
\begin{split}
R_n&=A_n^{ab}\partial_{ab}h_n+r_n,\\
A_n^{ab}
&=-\bigl(C^{ab}(g+h_n)-C^{ab}(g)
                    -\mathrm d C^{ab}_g[h_n]\bigr)\\
&=-\int_0^1(1-t)\,
       \mathrm d^2C^{ab}_{g+t h_n}[h_n,h_n]\,\dd t.
\end{split}
\end{equation*}
Here $\mathrm d C$ differentiates the metric field variables.
After extracting $A_n^{ab}\partial_{ab}h_n$, Taylor's formula gives
\begin{equation*}
\begin{split}
|r_n|\leq C\bigl(&\|h_n\|_{L^\infty_x}^3
                     (|\partial^2g|+|\partial g|^2)\\
&+\|h_n\|_{L^\infty_x}^2|\partial g||\partial h_n|
 +\|h_n\|_{L^\infty_x}|\partial h_n|^2\bigr).
\end{split}
\end{equation*}
Smoothness of $C$ and the original first-derivative bound therefore give
\begin{equation*}
\begin{gathered}
\|A_n\|_{L^\infty_x}=O(\|h_n\|_{L^\infty_x}^2),\\
\|\partial A_n\|_{L^\infty_x}+\|r_n\|_{L^\infty_x}
=O(\|h_n\|_{L^\infty_x}).
\end{gathered}
\end{equation*}
The symmetry $(A_n^{ab})^*=A_n^{ba}$ makes
$-\partial_a(A_n^{ab}\partial_b)$ self-adjoint.
Since $B_a$ is scalar, its divergence commutator is controlled entry
by entry by Lemma~\ref{lem:commutator}. Therefore
\begin{equation*}
\left|2\Im\ip{A_n^{ab}\partial_{ab}h_n}{B_ah_n}\right|
\leq C_a\|A_n\|_{W^{1,\infty}_x}\|h_n\|_{H^1_x}^2\longrightarrow0.
\end{equation*}
This argument uses no second derivative of $A_n$ or $h_n$.
The $r_n$ pairing tends to zero by its local $L^\infty$ bound and
the bounded $L^2$ norm of $B_ah_n$. The full remainder $R_n$ is
controlled only in this imaginary energy pairing; a small norm for
$R_n$ itself is neither needed nor asserted.
The smooth $E(g)$ pairing also vanishes. Combining this with
\eqref{eq:var-commutator} gives the ungauged identity
\begin{equation}\label{eq:var-cubic-residual}
4\int_{S_r^*U} H_Pa\,\dd\mu
=-\lim_n\Im\ip{E''_g[h_n,h_n]}{B_ah_n}.
\end{equation}

For real fields the scalar measure $\mu$ is invariant under
$\xi\mapsto-\xi$ (the matrix defect is conjugated), and
$H_Pa$ is odd when $a$ is even. It suffices to treat real odd $a$.
Choose scalar Weyl quantization with real odd full symbol and symmetric
localization, so that $B_a$ anticommutes with complex conjugation.
Then $-\ii B_ah$ is real for real $h$. If
$\mathcal T_g=\mathcal S'''_g$ is the symmetric real third variation,
extended complex multilinearly, the right side of
\eqref{eq:var-cubic-residual} is
$\lim_n\mathcal T_g(h_n,h_n,-\ii B_ah_n)$.

For $j\in\{1,2,3\}$, symmetric tensor fields $f_1,\ldots,f_j$,
and tensor-valued covectors $A,B$, define
\begin{equation*}
C_g^{(j)}[f_1,\ldots,f_j](A,B)
=\sum_{a,b=0}^3\sum_{I,J,K_1,\ldots,K_j=1}^{10}
\f{\partial^j C_{IJ}^{ab}}{\partial k^{K_1}\cdots\partial k^{K_j}}(g)
f_1^{K_1}\cdots f_j^{K_j} A_a^I B_b^J.
\end{equation*}
Commuting three variations of the first-order action gives
\begin{equation*}
\begin{split}
\mathcal T_g(f_1,f_2,f_3)=\int_U\Bigl[&
 \sum_{\mathrm{cyc}} C_g^{(1)}[f_i](\partial f_j,\partial f_k)\\
&+\sum_{\mathrm{cyc}} C_g^{(2)}[f_j,f_k](\partial f_i,\partial g)
 +\f12 C_g^{(3)}[f_1,f_2,f_3](\partial g,\partial g)\Bigr]\,\dd x.
\end{split}
\end{equation*}
Here the cyclic triples are $(1,2,3),(2,3,1),(3,1,2)$.
Put $Q=-\ii B_a$ and $(f_1,f_2,f_3)=(h,h,Qh)$.
For the sole term containing $\partial(Qh)$, namely
$C_g^{(2)}[h,h](\partial(Qh),\partial g)$, integrate once by parts.
The derivative then falls only on the fixed background coefficients,
the two $h$ factors, or $\partial g$, leaving $Qh$ times
$O_g(h\partial h)+O_g(h^2)$. The outer compact cutoff removes
boundary terms. H\"older's inequality and the order-one bound for $Q$ give
\begin{equation*}
\begin{aligned}
&\left|\mathcal T_g(h,h,Qh)
-\int_U\sum_{\mathrm{cyc}}
 C_g^{(1)}[f_i](\partial f_j,\partial f_k)\,\dd x\right|\\
&\quad\leq C_g\|h\|_{L^\infty_x}
 \bigg(\|\partial h\|_{L^2_x}+\|h\|_{L^2_x}\bigg)
 \|Qh\|_{L^2_x}\\
&\quad\leq C_{g,a}\|h\|_{L^\infty_x}\|h\|_{H^1_x}^2.
\end{aligned}
\end{equation*}
Thus all background-derivative contributions vanish
for the admissible sequence, without estimating $\partial(Qh)$ or
assuming that the smooth background is vacuum. The first cyclic sum,
with coefficients frozen at a constant metric $m$, is exactly the
principal third-variation vertex used next.

For frozen $m$, define the translation-invariant trilinear form on
Schwartz tensor fields by the whole-space integral
\begin{equation*}
\mathcal T_m(f_1,f_2,f_3)
=\int_{\R^4}\sum_{\mathrm{cyc}}
 C_m^{(1)}[f_i](\partial f_j,\partial f_k)\,\dd x.
\end{equation*}
It agrees with the localized frozen third variation on compactly
supported fields, and also on $(h,h,Qh)$ when $h\in C_c^\infty(U)$.
For $(p_1,H_1)=(p,H)$, $(p_2,H_2)=(q,J)$,
$(p_3,H_3)=(r,L)$, define the full vertex by
\begin{equation*}
\Gamma_m(p,H;q,J;r,L)
=-\sum_{\mathrm{cyc}}C_m^{(1)}[H_i]
                  (p_j\otimes H_j,p_k\otimes H_k).
\end{equation*}
It is symmetric in the three pairs and has frequency degree two.
Set $r=-p-q$ and $\alpha_a(p,q,r)=a(p)+a(q)+a(r)$. Fourier inversion gives
\begin{equation}\label{eq:var-cyclic}
\begin{split}
\mathcal T_m(h,h,-\ii a(D)h)
=-\f{\ii}{3(2\pi)^8}\int_{\R^4\times\R^4}
 \alpha_a(p,q,-p-q)\\
\qquad\qquad\qquad\times
 \Gamma_m(p,\widehat h(p);q,\widehat h(q);-p-q,\widehat h(-p-q))
 \,\dd p\,\dd q.
\end{split}
\end{equation}
The factor $(2\pi)^{-8}$ is the product-integral normalization in four
dimensions; $1/3$ averages the three equal test placements.

To make the polarization precise, for $Q=-\ii B_a$ define
\begin{equation*}
\mathcal V_Q(f_1,f_2,f_3)=\f13\bigl(
\mathcal T_g(Qf_1,f_2,f_3)+\mathcal T_g(f_1,Qf_2,f_3)
+\mathcal T_g(f_1,f_2,Qf_3)\bigr).
\end{equation*}
It is symmetric and satisfies
$\mathcal V_Q(h,h,h)=\mathcal T_g(h,h,Qh)$. For frozen coefficients
its Fourier representation is \eqref{eq:var-cyclic} with distinct fields.
All subsequent replacements and expansions use this complete polarized
form. It is continuous in three $W^{1,3}(\R^4)$ factors in the sense of
Lemma~\ref{lem:cubic-continuity}; no such estimate is asserted for an
individual test placement. To verify the extension to the present
matrix equation, polarize the coefficient-linear principal term. Its two
high slots form a symmetric matrix divergence operator. The scalar test
allows the finite-exponent commutator estimate in that lemma entry by
entry, with $(3,3,3/2)$ exponents. Its drift and semilinear terms have
three $L^3$ factors. Fixed inhomogeneous or differentiated background
terms add undifferentiated factors, which vanish for this sequence. In
particular the continuity bound has no largest-frequency parameter.

\subsection{Physical and gauge fields and the frozen Ward identity}

Use the smooth Euclidean gauge-complement projection of
Proposition~\ref{prop:general-projection}:
\begin{equation*}
\Pi_E(\xi)H=P_E(\xi)HP_E(\xi),\qquad
P_E(\xi)=\Id-\f{\xi\otimes\xi}{|\xi|_E^2}.
\end{equation*}
Choose a fixed real even high-frequency cutoff $\chi_{\rm hi}$,
vanishing near zero and equal to one outside a fixed ball. Its support
is chosen outside the infrared transition of the frozen multiplier
test, whose symbol there is exactly homogeneous. Lower-order terms of
the variable test have already been included in the freezing error.
Choose a second real even cutoff $\chi_{\rm reg}$ which vanishes in a
smaller ball and equals one on $\supp\chi_{\rm hi}$. Every infrared
regularization of a correction at either high-frequency slot below uses
$\chi_{\rm reg}$. In particular
$\chi_{\rm reg}\chi_{\rm hi}=\chi_{\rm hi}$: regularizing a correction
does not insert an unwanted second factor of $\chi_{\rm hi}$ into the
exact cubic identity. More precisely, its symbols $t$ and $\sigma$
are both multiplied by $\chi_{\rm reg}(q)\chi_{\rm reg}(p+q)$.
Thus $\ii [P_0,T_{\rm reg}]=S_{\rm reg}$ is an operator identity, while
the pairing on $u$ agrees with the unregularized energy because
$\chi_{\rm reg}(D)u=u$. The same nested choice applies to the original
high-frequency gauge field when it occupies a low slot.
For a localized
field $z$ put $z^{\rm hi}=\chi_{\rm hi}(D)z$ and
$u=\Pi_E(D)z^{\rm hi}$, $G=(\Id-\Pi_E(D))z^{\rm hi}$, with
$\Pi_E$ homogeneous in these two formulas. Then $G$ is exactly pure gauge
and $u$ is exactly the regularized projector of
Proposition~\ref{prop:general-projection} applied to the original compactly supported
field $z$. The omitted $z-z^{\rm hi}$ tends strongly to zero in $W^{1,p}(\R^4)$ for every fixed $1<p<\infty$. The gauge part satisfies, for $p\ne0$,
\begin{equation}\label{eq:var-gauge-potential}
\widehat G(p)=K_p(v(p)),\qquad
K_p(v)=p\otimes v+v\otimes p,\qquad
v(p)=\f{\widehat G(p)p}{|p|^2}
-\f{p\,p^{\mathsf T}\widehat G(p)p}{2|p|^4}.
\end{equation}
Thus $|p|v(p)$ depends smoothly on $p/|p|$ and linearly on $\widehat G(p)$.
The gauge factor need not be null. Both $u,G$ tend to zero in
$L^p(\R^4)$ and have bounded gradients there for each $1<p<\infty$.

For a frozen Lorentz metric $m$ choose a smooth degree-zero screen
projection $\Pi_T(\xi)$ on a fixed wide neighborhood of its null cone.
Its range consists of tensors $J$ satisfying
\begin{equation}\label{eq:var-screen}
J\xi=0,\qquad Jm^{-1}\xi=0,\qquad\tr_mJ=0.
\end{equation}
Here the first equation uses the coordinate Euclidean identification.
An explicit construction will also quantify its dependence on $m$.
Set $B=m^{-1}$, $e=\xi/|\xi|_E$, and
\begin{equation}\label{eq:var-screen-matrix}
\begin{gathered}
w=Be-(e^TBe)e,\qquad
Q=\Id-ee^T-\f{ww^T}{|w|_E^2},\\
\mathsf A=QBQ+(\Id-Q),\qquad
\mathsf R=Q\mathsf A^{-1}Q,\\
\Pi_{T,m}(\xi)H=QHQ-\f12\operatorname{tr}(BQHQ)\mathsf R.
\end{gathered}
\end{equation}
On the null cone $|w|_E=|Be|_E$ is nonzero. The range
$W=\operatorname{ran}Q=(\operatorname{span}\{e,Be\})^{\perp_E}$
is two-dimensional, and $B|_W$ is positive definite: the restriction of
$B$ to $e^{\perp_B}$ is positive semidefinite with radical $\R e$,
whereas $W\cap\R e=\{0\}$. On a compact family of nondegenerate
Lorentz metrics, $|w|_E$ and the least eigenvalue of $B|_W$ have
uniform positive lower bounds on the unit null set. They retain such
bounds on a fixed wide neighborhood. Thus \eqref{eq:var-screen-matrix}
is smooth, even, and degree zero there, with all fixed-order symbol and
metric derivatives uniformly bounded.
Since $Q\mathsf R=\mathsf RQ=\mathsf R$ and
$QB\mathsf R=Q$, we have $\operatorname{tr}(B\mathsf R)=2$.
Consequently \eqref{eq:var-screen-matrix} has range
\eqref{eq:var-screen} and fixes every tensor in that range, so
$\Pi_{T,m}^2=\Pi_{T,m}$. Multiplying by a smooth cutoff equal to one
on a smaller fixed wide cone gives the global order-zero multiplier
used below; its finite-$L^p$ bounds are uniform in $m$.
We write $\Pi_T=\Pi_{T,m}$ when $m$ is fixed. On a null covector its range equals
$\Pi_E(\xi)\ker\mathcal R_{m,\xi}$: Euclidean transversality and the
Lorenz constraint force first the trace and then the second contraction
in \eqref{eq:var-screen} to vanish. Consequently the screen for the
actual metric fixes the physical defect at each point. Replacing it
by the screen for a frozen metric need not fix that defect.
Section~\ref{sec:completion} quantifies this error and uses the complete
$W^{1,3}(\R^4)$ continuity above to insert the frozen screen in physical slots.

Write $J_q,L_r$ for actual Fourier tensors in the ranges
\eqref{eq:var-screen}, even when $q,r$ are slightly off cone. They obey
the exact identity
\begin{equation*}
A_m(q)J_q=\f12w_m q_m(q)J_q^{\sharp\sharp}
        =2w_m q_m(q)\mathsf C_mJ_q,
\qquad w_m=\sqrt{|\det m|}.
\end{equation*}
This is the phase-extraction identity used below. The underlying scalar
equation will always be applied to the original $u$, with additional
screen/cone projections retained inside the multiplier in its high-frequency slots.

For a fixed compactly supported vector field $X$, diffeomorphism
invariance gives $\mathcal S'_g[\mathcal L_Xg]=0$.
Twice differentiating in fixed metric directions $J,L$ gives the exact
off-shell Ward identity
\begin{equation}\label{eq:var-ward}
\mathcal S'''_g[J,L,\mathcal L_Xg]
+\mathcal S''_g[J,\mathcal L_XL]
+\mathcal S''_g[L,\mathcal L_XJ]=0.
\end{equation}
There is no third term involving $\mathcal S'_g$: the generator is
linear in the metric, so its second variation vanishes. In particular
this identity does not assume $E(g)=0$ for the smooth limit.

For frozen $m$, put $G_p=K_p(v)$ and set $V=m^{-1}v$ for the
corresponding vector. Thus the potential in
\eqref{eq:var-gauge-potential} has lower components, whereas Lie
derivatives below use $V$. Put $\Delta=q_m(r)-q_m(q)$. The highest-derivative Fourier component of
\eqref{eq:var-ward} gives
\begin{equation*}
\begin{split}
\Gamma_m(p,G;q,J;r,L)=-J\cdot A_m(q)
 \bigl((V\cdot r)L+K_p(L(V,\cdot))\bigr)\\
{}-L\cdot A_m(r)
 \bigl((V\cdot q)J+K_p(J(V,\cdot))\bigr).
\end{split}
\end{equation*}
This formula is bilinearly complexified, not Hermitian, with
$p+q+r=0$. It follows by using a Fourier vector generator; all factors
$\ii$ cancel from its three terms. For two screen tensors, use the bilinear extension of $\mathsf C_m$. The exact expression
symmetric in the two physical slots is
\begin{equation}\label{eq:var-low-gauge-identity}
\begin{split}
\Gamma_m(p,G;q,J;r,L)
={}&-w_m\Delta\,V\cdot(q-r)\mathsf C_m(J,L)\\
&+w_m(V\cdot p)(q_m(q)+q_m(r))\mathsf C_m(J,L)\\
&-2w_mq_m(q)\mathsf C_m(J,K_p(L(V,\cdot)))\\
 -2w_mq_m(r)\mathsf C_m(L,K_p(J(V,\cdot))).
\end{split}
\end{equation}
To see every coefficient without suppressing a physical phase, the
transport part of the Ward formula is
$-2w_m[q_m(q)(V\cdot r)+q_m(r)(V\cdot q)]\mathsf C_m(J,L)$. Since $p+q+r=0$,
\begin{equation*}
-2[q_m(q)(V\cdot r)+q_m(r)(V\cdot q)]
=-\Delta V\cdot(q-r)+(q_m(q)+q_m(r))V\cdot p.
\end{equation*}
The remaining two terms are obtained directly from
$A_m(q)J=2w_mq_m(q)\mathsf C_mJ$ and its $r$ counterpart.
This exact symmetric form isolates a Hermitian correction.

For the initial low-slot calculation, relabel the triad so that $p$ is
its low covector, and set
$\lambda_p=|p|$, $\lambda_q=|q|$, $\lambda_r=|r|$.
The two high magnitudes are comparable. Use the partition
\eqref{eq:analytic-triad-partition}. On a low piece $\lambda_p\leq2\kappa \lambda_q$,
put $p=\lambda_p\omega$, $q=\lambda_q\theta$, $z=p/\lambda_q$. Then
\begin{equation}\label{eq:var-cyclic-extension}
\f{\alpha_a}{\lambda_p}
=a(\omega)-\int_0^1\omega\cdot\partial a(\theta+tz)\,\dd t.
\end{equation}
The formula is smooth in independent $(\omega,z,\theta)$, including
$z=0$. In particular $\alpha_a=O(\lambda_p)$, with all normalized derivatives.
For comparable frequencies the same bound and derivative statement
hold on a compact partition of annular variables; the nonzero third
frequency is bounded below by the choice of the smallest slot.

For a low gauge leg, \eqref{eq:var-gauge-potential} gives
\begin{equation*}
|V|\ls\f{|G|}{\lambda_p},\qquad
|V\cdot p|\ls |G|,\qquad
|K_p(L(V,\cdot))|\ls |G||L|.
\end{equation*}
Thus each of the last three terms in \eqref{eq:var-low-gauge-identity},
after multiplication by $\alpha_a$, is
$\lambda_p\lambda_q^2$ times a smooth bounded normalized symbol with one factor
$q_m(q)/\lambda_q^2$ or $q_m(r)/|r|^2$ on its own physical factor.
There is no negative power of $\lambda_p/\lambda_q$.
The first term contains the exact scalar phase difference $\Delta$
and will be handled as a correction below.

For one high gauge leg, relabel the factors as physical $(p,H),(q,J)$
and gauge $(r,K_r(v))$, with $V=m^{-1}v$. Screen transversality gives
\begin{equation*}
\begin{gathered}
\mathsf C_m(H,K_r(J(V,\cdot)))
=-\f12H^{ab}q_aJ_{bc}V^c,\\
\mathsf C_m(J,K_r(H(V,\cdot)))
=-\f12J^{ab}p_aH_{bc}V^c.
\end{gathered}
\end{equation*}
Substituting these contractions into the Ward identity yields
\begin{equation*}
\begin{split}
\Gamma_m(p,H;q,J;r,K_r(v))
={}&-2w_m(q_m(p)V\cdot q+q_m(q)V\cdot p)\mathsf C_m(H,J)\\
&+w_mq_m(p)H^{ab}q_aJ_{bc}V^c
 +w_mq_m(q)J^{ab}p_aH_{bc}V^c.
\end{split}
\end{equation*}
If $p$ is low, write $\rho=\lambda_p/\lambda_q$. Since $|V|\ls \lambda_q^{-1}|G|$,
the coefficient $\alpha_a\Gamma_m/(\lambda_p\lambda_q^2)$ is a sum of smooth
bounded normalized symbols times $\rho^2q_m(p)/\lambda_p^2$ and
$\rho q_m(q)/\lambda_q^2$, respectively. Formula
\eqref{eq:var-cyclic-extension} makes these symbols smooth in the
independent angles and $\rho$ down to zero.

For two gauge legs the corresponding exact identity is
\begin{equation*}
\Gamma_m(p,H;q,K_q(v_q);r,K_r(v_r))
=w_mq_m(p)\bigl(v_r\cdot m^{-1}v_q\bigr)H^{ab}q_aq_b.
\end{equation*}
To obtain it, use the $r$ leg as generator. The other gauge Hessian
vanishes, while its Lie derivative is
$-(m^{-1}v_r\cdot q)K_p(v_q)
 +(v_r\cdot m^{-1}v_q)K_r(q)$; the first term pairs to zero with
the physical Hessian. If the physical leg is low, set
$\rho=\lambda_p/\lambda_q$; the normalized coefficient gains
$\rho^2$ times its own phase $q_m(p)/\lambda_p^2$.
If a gauge leg is low, interchange $q,r$ so that $q$ is low and set
$\rho=\lambda_q/\lambda_p$. The identity
$H^{ab}q_aq_b=H^{ab}r_ar_b$ puts the low covector in both factors.
After normalization by $\lambda_q\lambda_p^2$, the coefficient gains
$\rho$ times $q_m(p)/\lambda_p^2$. These formulas have
the same smooth angular and radial extensions. With three gauge
legs both Hessians vanish, so the frozen vertex is exactly zero.
On the comparable piece every frequency is bounded below by
$c_\kappa \lambda_q$; a gauge generator may therefore be chosen in any fixed
gauge slot, including the least one, with smooth annular bounds.

The right side of \eqref{eq:var-cyclic} defines a symmetric trilinear
energy functional in its three fields. Substituting $h=u+G$ gives its
terms with zero, one, two and three gauge fields with coefficients
$1,3,3,1$. Each term was classified above; no reduced-equation or
selected semilinear vertex is substituted for the action vertex.

Every normalized physical-phase term obeys an all-scale trilinear
estimate by three gradient $L^3$ norms, one of them replaced by
\begin{equation*}
\left\|\f{q_m(D)}{|D|^2}\,\partial v_{\rm phys}\right\|_{L^3_x}.
\end{equation*}
The precise multiplier interface is the first assertion of
Lemma~\ref{lem:analytic-near}.
After relabeling the least-frequency factor, the remaining symbol is
smooth in two angular variables and the ratio $\rho=\lambda_p/\lambda_q$ down to zero.
Its value at $\rho=0$ is a separated angular expansion of a lowpass/high
paraproduct. Its remainder is $\rho R$, summed over gaps using
$\sum_s2^{-s}<\infty$ and annular Fourier expansion. Comparable-frequency factors
have finitely many relative shifts and the three square-function bound.
The physical phase is moved to its own factor before these operations.
Consequently this estimate has no maximum-frequency dependence and
requires no smallness of a cone-multiplier norm. Strong phase localization
therefore removes these terms at fixed frozen parameters, up to the
explicit freezing error below.

\subsection{The two corrections and the physical near sum}

For the remaining one-low-gauge term, consider first a single symmetric
trilinear energy in \eqref{eq:var-cyclic}, with the gauge field put in
the first slot. Its transport coefficient from
\eqref{eq:var-low-gauge-identity} is
\begin{equation*}
\sigma_G=\f{\ii w_m}{3}\alpha_a\Delta\,
              V\cdot(q-r)\mathsf C_m.
\end{equation*}
The low Fourier amplitude is understood in $V=m^{-1}v$, and the two
screen projections are understood on the high arguments of $\mathsf C_m$.
Thus in a fixed tensor basis the high-slot matrix is
$\Pi_{T,m}(-r)^T\mathsf C_m\Pi_{T,m}(q)$, multiplied by the displayed
scalar; the projection at $p+q$ is transposed even when the screen
projection is not Euclidean self-adjoint.
Define its correction by the symbol
\begin{equation}\label{eq:var-gauge-correction}
t_G=\f{w_m}{3}\alpha_a V\cdot(q-r)\mathsf C_m.
\end{equation}
With high factor frequency $q$ and product frequency $p+q=-r$,
$\sigma_G=\ii(q_m(p+q)-q_m(q))t_G$. Consequently
\begin{equation*}
S_G=\ii [P_0,T_G],\qquad P_0=-m^{ab}\partial_{ab},
\end{equation*}
\emph{including at $\Delta=0$}; the factor is cancelled algebraically
and no nonresonance division occurs. The other physical-phase terms in
\eqref{eq:var-low-gauge-identity} need not vanish at $\Delta=0$;
only the transport term is represented by this commutator.
The contribution with exactly one
gauge field is three copies of this symmetric trilinear form, so its
correction has the same factor three. No further Taylor factorial is used.

We claim that these high-slot operators are Hermitian.
Indeed, under high adjunction
$(p,q)\mapsto(-p,p+q)$ the signed triad changes to $(-p,-r,-q)$.
Oddness of $a$ changes the sign of $\alpha_a$. Reality of $G$ and
\eqref{eq:var-gauge-potential} give $V(-p)=-\overline{V(p)}$, while
$q-r$ remains the same. The two minus signs cancel, and the real
contraction, the even high projections, and the symmetric cutoffs
transpose to one another. Thus $T_G=T_G^*$ exactly. This verification
uses the already symmetrized $V\cdot(q-r)$; it does not infer the
adjunction of an isolated summand from that of the entire energy.

The symbol in \eqref{eq:var-gauge-correction} is order one in the high
frequency with a regular
low/high extension. In the variables of
\eqref{eq:var-cyclic-extension}, its coefficient on the low tensor is
$\lambda_q M_a(\omega,z,\theta)$: $\alpha_a/\lambda_p$ has that exact integral
extension, $\lambda_p V$ is smooth angular in the low tensor, and
$(q-r)/\lambda_q=2\theta+z$. High projections depend smoothly on $\theta$ and
$\theta+z$. This holds for \emph{every} low direction $\omega$;
there is no low-gauge cone restriction. Choose radial cutoffs smooth
in $|z|^2$, with their high-slot symmetric counterparts.

To state precisely why the bilinear commutator theorem applies, expand
only the independent low angle:
\begin{equation*}
M_a(\omega,z,\theta)
=\sum_\nu Y_\nu(\omega)c_\nu(z,\theta).
\end{equation*}
All required $z,\theta$ derivatives of $c_\nu$ decay faster than any
power of the spherical-harmonic index. Move $Y_\nu(p/\lambda_p)$ to the low
field $G$. Its fixed finite-$L^p$ norms grow only polynomially in that
index. The remaining symbol $\lambda_q c_\nu(p/\lambda_q,q/\lambda_q)$ has Cartesian low and
high derivatives bounded by
\begin{equation*}
C_{N,\alpha,\beta}(1+\nu)^{-N}
(1+|p|+|q|)^{1-|\alpha|-|\beta|},
\end{equation*}
after fixed high-frequency infrared regularization. Every fixed
\emph{low}-frequency cutoff is placed on $G$, before its angular
multipliers, and is not part of the Cartesian symbol. This is essential:
a cutoff derivative at fixed $|p|$ would not gain $\lambda_q^{-1}$ as
$\lambda_q\to\infty$. With the smooth $z=0$ extension the remaining symbol
belongs to $BS^1_{1,0}$, including at $p=0$. The series is summable in every seminorm used in
Lemma~\ref{lem:analytic-correction-bounds}. In particular, no low angular
factor with derivatives $|p|^{-1}$ is left inside a claimed Cartesian
$BS^1$ symbol.

There remain the three-physical terms. Fix a cap width $\delta$ and
cone thickness $\eta=c\kappa\delta^2$ as in Section~\ref{sec:analytic}.
Triad geometry in \eqref{eq:analytic-phase-gap} places every comparable
physical triad in adjacent projective null caps. Split low/high
physical triads into near and separated directions.

On the near region the test factor alone has the quadratic angular
gain
\begin{equation}\label{eq:var-near-test-gain}
\alpha_a=\delta^2\lambda_p M,
\end{equation}
where $M$ has uniform smooth normalized cap derivatives and is smooth
at $\lambda_p/\lambda_q=0$. To prove this at the symbol level, use the parabolic map
$\mathcal A$ from the transfer lemma and take its central ray
$e_\gamma=E_\gamma(1,0,0,0)$. Enlarge its fixed cone once more, if
needed, to contain the low/high interpolation segments
$\theta'+t p'/|q'|$ for $0\leq t\leq1$. Low/high separation keeps their
longitudinal coordinate away from zero, so this enlargement is fixed
independently of $\delta$. On that cone set
\begin{equation*}
\widetilde a_\delta(\zeta)=\delta^{-2}
 \bigl[a(\mathcal A\zeta)-\partial_\xi a(e_\gamma)\cdot\mathcal A\zeta\bigr].
\end{equation*}
Oddness and homogeneity give the same linear first jet on both signed
central rays. More explicitly, for $\zeta=(s,t,z)$ put
$v_\delta=E_\gamma(0,\delta t,z)$, so that
$\mathcal A\zeta=s e_\gamma+\delta v_\delta$. Taylor's integral formula
and Euler's homogeneous identity give the exact extension
\begin{equation*}
\widetilde a_\delta(\zeta)
=\int_0^1(1-\tau)D^2a(s e_\gamma+\tau\delta v_\delta)
                         [v_\delta,v_\delta]\,\dd\tau.
\end{equation*}
The longitudinal component $s$ is bounded away from zero on every
normalized annulus of the enlarged cone, for either sign of $s$.
The segment therefore stays away from the origin. This formula makes
$\widetilde a_\delta$ an odd degree-one symbol with uniform normalized
seminorms and a smooth extension to $\delta=0$. For
$p=\mathcal A p'$, $q=\mathcal A q'$, $r=\mathcal A r'$,
the subtracted linear terms sum to zero exactly, so
\begin{equation*}
\alpha_a=\delta^2\bigl[
\widetilde a_\delta(p')+\widetilde a_\delta(q')
                         +\widetilde a_\delta(r')\bigr].
\end{equation*}
Applying the integral identity \eqref{eq:var-cyclic-extension} to
this uniformly bounded symbol gives the low factor $|p'|$ and its
smooth extension to zero radial ratio. The transfer lemma's uniform
bound for $|\mathcal A p'|/|p'|$ converts it to $\lambda_p$.
The comparable pieces lie on compact annular boxes with all three
magnitudes bounded below. This proves every required normalized
derivative in \eqref{eq:var-near-test-gain}, not only a pointwise bound.

Here is the full normalized kernel, including the actual screen tensors.
Set $\lambda_p'=|p'|$, $\lambda_q'=|q'|$,
$p'=\lambda_p'\omega$, $q'=\lambda_q'\theta$, $\rho'=\lambda_p'/\lambda_q'$ and
$v=\theta+\rho'\omega$, and put
$\mathsf S_\delta(\zeta)=\Pi_{T,m}(\mathcal A\zeta)$. For independent
tensors $H,J,L$ define
\begin{equation*}
\mathfrak G_\delta=
\Gamma_m(\rho'\mathcal A\omega,\mathsf S_\delta(\omega)H;
         \mathcal A\theta,\mathsf S_\delta(\theta)J;
         -\mathcal A v,\mathsf S_\delta(v)L).
\end{equation*}
On a sufficiently small fixed $\rho'$ interval the normalized kernel is
\begin{equation}\label{eq:var-near-normalized-kernel}
\f{\alpha_a\Gamma_m}{\delta^2\lambda_p\lambda_q\lambda_r}
=\f{\mathfrak G_\delta}
 {|\mathcal A\omega|\,|\mathcal A\theta|\,|\mathcal A v|}
 \bigg[\widetilde a_\delta(\omega)
 -\int_0^1 D\widetilde a_\delta(\theta+t\rho'\omega)
                         \cdot\omega\,\dd t\bigg],
\end{equation}
The vertex on the left has those same screen tensors
at the original frequencies. Its vertex factor is
$\lambda_q'^2\mathfrak G_\delta$, its cyclic factor is $\delta^2\lambda_p'$ times
the bracket, and $\lambda_p\lambda_q\lambda_r=\lambda_p'\lambda_q'^2|\mathcal A\omega|\,
|\mathcal A\theta|\,|\mathcal A v|$. All three denominators and
the longitudinal coordinate of $\theta+t\rho'\omega$ are bounded
away from zero on the enlarged fixed cone. The residual transition
region $\rho'\geq\rho_0>0$ has only finitely many relative radial
shifts and is handled by the normalized annular boxes.

The actual angular cutoffs have the same uniform extension.
For fixed transformed covectors $u,v$ set
$\mathfrak D_\delta(u,v)=\delta^{-1}(\mathcal A u\wedge\mathcal A v)$.
It is polynomial in $\delta$, since the powers in the wedge of
$\operatorname{diag}(1,\delta^2,\delta,\delta)$ are all at least one.
Writing $q_{\rm std}(s,t,z)=-st+|z|^2$, we have
\begin{equation*}
\f{\mathfrak d(\mathcal A u,\mathcal A v)}{\delta^2}
=\f{|\mathfrak D_\delta(u,v)|^2}
 {|\mathcal A u|^2|\mathcal A v|^2},\qquad
\f{q_m(\mathcal A u)}{c\kappa\delta^2|\mathcal A u|^2}
=\f{q_{\rm std}(u)}{c\kappa|\mathcal A u|^2}.
\end{equation*}
The cone-restricted lower bounds for $|\mathcal A u|/|u|$, and the
uniform matrix bounds in \eqref{eq:var-screen-matrix}, therefore give
smooth extensions to $\delta=0$ for the near cutoffs, the thin-cone
profiles, and the screen symbols. Together with
\eqref{eq:var-near-normalized-kernel}, this gives their simultaneous
normalized derivative bounds, including at $\rho'=0$.

On comparable annular blocks the third cap and the annulus of the product frequency
are retained before any extension of the bilinear symbol. With
normalized covectors $\zeta,\upsilon$, their support gives
$|\zeta+\upsilon|\geq c_\kappa$ and
$|s(\zeta+\upsilon)|\geq c|\zeta+\upsilon|$. Thus the third screen
and the third test symbol are evaluated away from their zero-frequency
singularity, uniformly in $\delta$. Choose slightly larger fixed
cutoffs preserving both lower bounds, define the symbols there, and
extend their products by zero to the annular boxes. This supplies
the extension required by Lemma~\ref{lem:analytic-near}; it does not
infer product-frequency localization from the two factor caps alone. The same
argument treats the fixed low/high transition blocks.
For the chosen profiles and a compact family $\mathcal K$ of frozen
metrics, each cap-localized block in its normalized variables obeys
\begin{equation*}
\mathfrak A_{100}(\delta^{-2}m_{\mathrm{near}})
\leq C_{\mathcal K,\kappa}\|a\|_{C^{103}(\mathbb S^3)},
\end{equation*}
where $m_{\mathrm{near}}$ is that block's full multiplier after removing
its three first-derivative weights, and the norm on $a$ is that of its
degree-one homogeneous restriction to the Euclidean unit sphere.
The two derivatives in Taylor's formula and the additional derivative
in the cyclic integral explain this sufficient, nonoptimal order.
These are frequency bounds uniform over the cap centers; the discrete
cap covering is not differentiated with respect to $\delta$.
Fixed infrared cutoffs on the factors are kept outside this homogeneous core.

The action vertex $\Gamma_m/\lambda_q^2$ is smooth in these variables: it is a
finite sum of metric contractions with two differentiated covectors
and three tensor factors, and the fixed-wide screen symbols are smooth.
Moreover, $\Gamma_m$ has degree two and $\alpha_a$ degree one under
simultaneous frequency dilation. After the three gradient weights are
removed, the core symbol therefore has degree zero, as required by
Lemma~\ref{lem:analytic-near}. All projectors, cap cutoffs and normalized
physical phases have degree zero as well. Fixed absolute-frequency
cutoffs act on the fields and do not change this homogeneous core.
Thus the full energy has the $\delta^2\lambda_p\lambda_q^2$ symbol hypothesis of
Lemma~\ref{lem:analytic-near}. Its transfer proof applies unchanged: in
one cap the common parabolic linear map preserves the trilinear norm;
the $\rho=0$ part is a lowpass/high paraproduct, the $\rho R$ part has
a summable gap factor, and comparable blocks have finitely many radial
shifts. There are $O(\delta^{-2})$ caps. Each individual cap has uniform
$L^p$ bounds, while its $L^3(\dd x\times\#)$ norm loses at most
$N_\delta^{1/(3(p_*-2))}$ per factor by the $L^2/L^{p_*}$ interpolation
in \eqref{eq:analytic-cap-bounds}. Hence
\begin{equation}\label{eq:var-near-bound}
\begin{split}
|\mathcal E^{\mathrm{near}}_{\rm physical}|
&\leq C_{\kappa,p_*}\delta^2N_\delta^{1/(p_*-2)}
                         \|\partial u\|_{L^3_x}^3\\
&\leq C_{\kappa,p_*}\delta^{\beta_*}\|\partial u\|_{L^3_x}^3.
\end{split}
\end{equation}
Only the fixed-wide projected field occurs on the right; the thin
cutoff is inside individual caps. There is no norm of the full
thin-cone projection multiplying $\delta^{\beta_*}$.

For separated physical low/high triads, the complete high-slot
Hermitian numerator has normalized size $\lambda_p\lambda_q^2$ by
\eqref{eq:var-cyclic-extension}. Divide it by the exact
$\ii \Delta=\ii(q_m(p+q)-q_m(q))$ using
\eqref{eq:analytic-exact-correction}. The support obeys
$|\Delta|\geq c_\delta \lambda_p\lambda_q$. The required symbol extension is explicit:
\begin{equation*}
\f\Delta{\lambda_p\lambda_q}
=2m(\omega,\theta)+(\omega\cdot z)q_m(\omega).
\end{equation*}
The numerator divided by $\lambda_p\lambda_q^2$ extends smoothly in independent
$\omega,z,\theta$: its cyclic factor is
\eqref{eq:var-cyclic-extension}, its vertex is polynomial in the
normalized covectors, and every symbol at the product frequency is smooth at
$\theta+z$. The actual denominator bound is initially known only on
the physical set $z=\rho\omega$. To justify an independent-variable
extension, write
\begin{equation*}
d_m(\omega,z,\theta)
=2m(\omega,\theta)+(\omega\cdot z)q_m(\omega).
\end{equation*}
The low cutoff extends as
$\vartheta(|z|^2/[\kappa^2(1+|\theta+z|^2)])$; the two angular
arguments extend as
\begin{equation*}
1-(\omega\cdot\theta)^2,\qquad
1-\f{[\omega\cdot(\theta+z)]^2}{|\theta+z|^2}.
\end{equation*}
They are smooth on a fixed compact extension domain with $|z|<1/2$.
By \eqref{eq:analytic-phase-gap}, including its $\rho=0$ limit, there
is $c_0>0$ with $|d_m|\geq c_0\delta^2$ on the entire physical far
support. Set
\begin{equation*}
\psi_\delta=1-\vartheta\!\left(\f{d_m^2}{c_0^2\delta^4}\right).
\end{equation*}
Then $\psi_\delta=1$ on that support and vanishes where
$|d_m|\leq c_0\delta^2/2$. Replacing the reciprocal denominator in
the extension by $\psi_\delta/d_m$, extended by zero at its zeros,
gives a smooth function with finite $C_\delta$ seminorms. It agrees
with the exact quotient on every physical frequency used by the
operator. Therefore the exact correction is
$\lambda_q M_\delta(\omega,z,\theta)$ and has the same convergent low-angular
and Cartesian $BS^1$ expansion just proved for $T_G$.
Hermitian symmetry follows from the Hermitian numerator and the sign
change of the real $\Delta$ under high adjunction. The other two physical
low-slot placements follow by permutation; their equal contributions
are counted explicitly in \eqref{eq:var-full-assembly} below.

Finite angular sums need not retain the original far support; they
retain the symmetric low/high support needed for the operator bounds.
The exact numerator is identified only after the full correction is
recovered, and no further division by $\Delta$ is performed on these sums.

Both corrections can now be paired with the scalar equation from
Proposition~\ref{prop:general-projection}. On a patch take the original
$u=\Pi_E(D)\chi_{\rm hi}(D)(\chi h_n)$ in the high-frequency slot and write its global extended
operator as $P_n=P_0+L_d$, where $L_d=-\partial_a(d^{ab}\partial_b)$.
The low field $H$ is a fixed angular projection of either $G$ or $u$.
For either correction, the preceding symbol bounds and Lemma~\ref{lem:analytic-correction-bounds} give
\begin{equation*}
\|T_Hu\|_{L^2_x}\leq C_\delta\|H\|_{L^4_x}\|\partial u\|_{L^4_x}
\end{equation*}
and the complete estimate \eqref{eq:analytic-quartic-bound}. Specifically,
\cite[Theorem 1]{BenyiOh} applies to each $BS^1_{1,0}$ term in the
high-frequency slot with $L^4\times L^4\to L^2$; the other high adjoint uses
$L^4\times L^2\to L^{4/3}$ and duality. Translation invariance gives
$[\partial,T_H]=T_{\partial H}$. Summing the rapidly decaying angular
series gives all three divergence commutator terms.
Lemma~\ref{lem:analytic-correction-bounds} removes angular sums in
operator norm and upper-frequency cutoffs strongly at fixed $n$.
The sequence limit is taken only after both cutoffs have been removed.
The source identity and the three commutator estimates give
\begin{equation}\label{eq:var-correction-energy}
\begin{split}
|\ip{S_Hu}{u}|
\leq{}&2\|P_nu\|_{L^2_x}\|T_Hu\|_{L^2_x}\\
&+C_\delta\|d\|_{L^\infty_x}\|\partial u\|_{L^2_x}
                         \|\partial H\|_{L^4_x}\|\partial u\|_{L^4_x}\\
&+C_\delta\|\partial d\|_{L^\infty_x}\|\partial u\|_{L^2_x}
                         \|H\|_{L^4_x}\|\partial u\|_{L^4_x}.
\end{split}
\end{equation}
The second line controls the first and third divergence commutators,
and the last controls the middle one. This is applied to the full
Hermitian correction. In particular no equation is applied to a gauge
field, and the matrix Einstein Hessian is not used as the nonlinear
wave operator in this step.

We record the complete frozen assembly explicitly. Let
$\mathcal F_\chi(f_1,f_2,f_3)$ denote the right side of
\eqref{eq:var-cyclic} with distinct fields and an additional multiplier
$\chi(p,q,r)$. Put $v=S(D)u$, where $S$ is the frozen screen and
thin-cone multiplier used in the physical replacement, and set
\begin{equation*}
\mathcal N_a=\mathcal F_{L_{\rm comp}+L_pN_p+L_qN_q+L_rN_r}(v,v,v).
\end{equation*}
The near profiles $N_q,N_r$ are defined cyclically. Let
$\mathcal F^{\rm ph}_{L_p}(G,v,v)$ retain the same factor
$-\ii \alpha_a/3$ and cutoff, but replace the vertex by the last three
terms of \eqref{eq:var-low-gauge-identity}. Define
\begin{equation*}
\mathcal P_a=3\bigl[
 \mathcal F^{\rm ph}_{L_p}(G,v,v)
 +\mathcal F_{L_q+L_r+L_{\rm comp}}(G,v,v)
 +\mathcal F_1(v,G,G)\bigr].
\end{equation*}
Choose $T_F(u)$ to represent the single physical-low far piece
$\mathcal F_{L_p(1-N_p)}(v,v,v)$, and $T_G(G)$ to represent the
transport part of $\mathcal F_{L_p}(G,v,v)$.
In $T_F(u)$ the low amplitude is $\widehat u(p)$ and the symbol
contains $S$ exactly once in each of its three slots; in $T_G(G)$
only the two high slots contain $S$. This convention avoids applying
the generally nonidempotent thin-cone profile twice to a low factor.
All infrared cutoffs on the high factor and product frequencies remain the nested ones
specified above. The exact identity is
\begin{equation}\label{eq:var-full-assembly}
\begin{split}
\mathcal F_1(v+G,v+G,v+G)
={}&\mathcal N_a+3\ip{\ii [P_0,T_F(u)]u}{u}\\
&+3\ip{\ii [P_0,T_G(G)]u}{u}+\mathcal P_a.
\end{split}
\end{equation}
The first factor three counts the three equivalent physical low slots;
the second comes from the one-gauge term in the field expansion.
The two-gauge term is entirely in $\mathcal P_a$, and the three-gauge
vertex vanishes. For the comparable one-gauge part, including a least
gauge leg, use the original Ward formula to extract the physical
phases directly; no low/high correction is applied there.
The preceding phase classification and all-scale transfer give
\begin{equation*}
|\mathcal P_a|\leq C_{\delta,\kappa,a,m}
 \left\|\f{q_m(D)}{|D|^2}\partial v\right\|_{L^3_x}
 \bigg(\|\partial v\|_{L^3_x}+\|\partial G\|_{L^3_x}\bigg)^2.
\end{equation*}
Thus every frozen principal cubic term is accounted for by a near
term, one of the two corrections, or a physical phase on its own
factor. These are identities and estimates at fixed fields and
parameters; no separate limit of the individual summands is assumed.

\section{Proof of Theorem~\ref{thm:main}}\label{sec:completion}

\subsection{The transport equation on a variable background}

\begin{proof}[Proof of Theorem~\ref{thm:variational-transport}]
Fix compact sets $K\Subset K^+$ in the chart, with the test based in
$K$. In this proof $L^p_x=L^p(\R^4,\dd x)$ and
$W^{1,p}_x=W^{1,p}(\R^4)$ refer to zero-extended coordinate fields.
The Euclidean tensor and derivative component norms are used throughout. Fix a finite $p_*>3$ and the exponents in
\eqref{eq:optimized-exponents}. Hold $\kappa$ fixed, then fix $\delta$ and
$\eta=c\kappa\delta^2$. Choose $\ell$ sufficiently small relative to
this fixed cone thickness. Take nonnegative
bounded-overlap cutoffs with
\begin{equation*}
\sum_j\chi_j^3=1\quad\hbox{near }K,\qquad
\supp\chi_j\subset B(x_j,C\ell),\qquad
|\partial^\alpha\chi_j|\leq C_\alpha\ell^{-|\alpha|}.
\end{equation*}
Choose nested nonnegative cutoffs $\chi_j^+$ equal to one on a
$c\ell$-neighborhood of $\supp\chi_j$, with bounded-overlap supports
$K_j^+\Subset K^+$ and
\begin{equation*}
K_j^+\subset B(x_j,C_+\ell),\qquad
|\partial^\alpha\chi_j^+|\leq C_\alpha^+\ell^{-|\alpha|}.
\end{equation*}
These choices are obtained from fixed nested cutoffs at scale $\ell$;
all constants are independent of $j$ and $\ell$. Put
$z_{n,j}=\chi_jh_n$ and $m_j=g(x_j)$.

Insert the three equal factors in the complete cubic action energy in
\eqref{eq:var-cubic-residual}. At fixed $\ell$ this changes it by
$o_n(1)$. To account for the nonlocal test outside the partition region,
put $\tau=1-\sum_j\chi_j^3$. Since $\tau a=0$, after the fixed outer
compact localization one has $\tau Q\in\Psi^0$ and
$\tau\partial_bQ\in\Psi^1$, where the derivative acts on the function obtained by applying $Q$. Their remaining action placements are bounded by
$C_\ell\|h_n\|_{L^3_x}\|h_n\|_{W^{1,3}_x}^2=o_n(1)$.
The action's principal third variation has two derivatives
distributed over its three factors. A cutoff derivative in it gives an
undifferentiated $h_n$; commuting $B_a$ through a cutoff gives an
order-zero operator on $h_n$. If its differentiated test factor appears
to give a second derivative, use the symmetric principal divergence
form before estimating. More explicitly, $[Q,\chi_j]$ is order zero;
its three placements are controlled as in Lemma~\ref{lem:freezing}.
A remaining term of type $(D\chi_j)h_n^2D(Qh_n)$ is integrated once
by parts, leaving terms $h_nDh_nQh_n$ or $h_n^2Qh_n$ with fixed
coefficient derivatives. The resulting terms contain an undifferentiated
$h_n$ and bounded finite-exponent gradient factors. They vanish as
$n\to\infty$. This is the same finite-exponent commutator argument
used above, now with fixed cutoff derivatives. Lower terms in the third
variation have fewer than two field derivatives and obey this estimate
directly. Properly supported off-diagonal kernels with separated nested
cutoffs are smoothing. There are finitely many patches at each fixed
$\ell$, so all these $o_n(1)$ terms can be summed before changing
$\ell$.

Freeze the principal action coefficients, density, and scalar test
symbol at $x_j$. The error is
\begin{equation}\label{eq:var-freezing-error}
C_\delta\ell\|\partial z_{n,j}\|_{L^3_x}^3+o_n(1).
\end{equation}
Apply Lemma~\ref{lem:freezing} to the complete third variation, with
$Q=-\ii B_a$. At fixed $\ell$ its second term satisfies
\begin{equation*}
B_\ell\|z_{n,j}\|_{L^3_x}\|\partial z_{n,j}\|_{L^3_x}^2=o_n(1).
\end{equation*}
Its first term is precisely the leading bound in
\eqref{eq:var-freezing-error}. The lemma includes the density,
quantization, and lower background terms and gives constants uniform
in the patch index. All differentiated coefficients belong to the
fixed smooth background or test. No second derivative of $g_n$ is used.

Now decompose the \emph{already localized} field globally in Fourier
space:
\begin{equation*}
z_{n,j}^{\rm hi}=\chi_{\rm hi}(D)z_{n,j},\qquad
u_{n,j}=\Pi_E(D)z_{n,j}^{\rm hi},\qquad
G_{n,j}=(\Id-\Pi_E(D))z_{n,j}^{\rm hi}.
\end{equation*}
Finite-$L^p$ multiplier bounds and the original gradient bound give
\begin{equation}\label{eq:var-local-volume}
\limsup_n\bigl(\|\partial u_{n,j}\|_{L^p_x}+
                      \|\partial G_{n,j}\|_{L^p_x}\bigr)
\leq C_p|K_j^+|^{1/p},\qquad 1<p<\infty.
\end{equation}
Both undifferentiated fields tend to zero in these norms. Their
noncompact tails are part of these global multiplier estimates; no
copy of a global field norm is charged to each patch.

All $L^2$ defect bounds below are global and follow directly from the
same quadratic measure. More generally, let $\chi$ be a fixed real
compactly supported cutoff and $T(D)$ a fixed classical order-zero
matrix multiplier, smooth at zero, with homogeneous principal symbol
$t_0$. Write $\dd\sigma_{dx}=w^{-1}\dd\sigma$, where
$w=\sqrt{|\det g|}$, and express this covariance in a fixed Euclidean
orthonormal frame of symmetric tensors. Then
\begin{equation}\label{eq:var-global-defect-norm}
\lim_n\|\partial T(D)(\chi h_n)\|_{L^2(\mathbb R^4)}^2
=\int_{S_r^*M}\chi^2|\xi|_E^2
       \operatorname{tr}\bigl(t_0\,\dd\sigma_{dx}\,t_0^*\bigr).
\end{equation}
The integrand is supported in the coordinate chart; the derivative
and the norm on the left are Euclidean coordinate ones. Indeed,
$\partial T(D)(\chi h_n)=T(D)(\chi\partial h_n)+o_{L^2_x}(1)$.
Its squared norm is the quadratic pairing with
$\chi T(D)^*T(D)\chi$, whose kernel is supported in
$\supp\chi\times\supp\chi$. This is a properly supported classical
order-zero operator with principal symbol $\chi^2t_0^*t_0$;
its lower-order remainder has zero defect limit. Apply
\eqref{eq:sigma}, with the density conversion above, and sum over
the derivative components. This proves the formula without
discarding any nonlocal tail. In particular it applies to
every wide-screen, normalized-phase and thin-cone-complement error
used below. At each fixed $\ell,\delta$ all these multipliers and
cutoffs are fixed tests of the already extracted measure. No
uniform symbol bound as $\ell,\delta\to0$, and no further subsequence,
is required for this quadratic-limit identity.

We quantify the screen and phase errors for the frozen metric. At an
actual $g(x)$-null covector over $K_j^+$ the range of the defect of
$u_{n,j}$ is the Euclidean gauge-complement representative of the
actual Ricci kernel. More precisely, for
$k\in\ker\mathcal R_{g(x),\xi}$ on this cone, the matrix construction
\eqref{eq:var-screen-matrix} and its uniform metric derivative bound give
\begin{equation*}
\begin{split}
|(\Id-\Pi_{T,m_j}(\xi))\Pi_E(\xi)k|_E
&=|(\Pi_{T,g(x)}(\xi)-\Pi_{T,m_j}(\xi))\Pi_E(\xi)k|_E\\
&\leq C|g(x)^{-1}-m_j^{-1}|_E|k|_E
\leq C\ell|k|_E.
\end{split}
\end{equation*}
The wide cutoffs equal one on these directions for small $\ell$.
This is an operator bound on the original defect range, so its
quadratic compression is $O(\ell^2)$. Also
$q_{m_j}(\xi)/|\xi|^2=O(\ell)$ there. Use the global wide extension of $\Pi_{T,m_j}$ from
Section~\ref{sec:variational-transport}. Define the two gradient errors by
\begin{equation*}
\begin{gathered}
E_{n,j}^{(1)}=(\Id-\Pi_{T,m_j}(D))\partial u_{n,j},\\
E_{n,j}^{(2)}=\f{q_{m_j}(D)}{|D|^2}
                  \Pi_{T,m_j}(D)\partial u_{n,j}.
\end{gathered}
\end{equation*}
Their symbols have size $O(\ell)$ on the defect range. The local
gradient bound and \eqref{eq:var-global-defect-norm} therefore give,
for $r\in\{1,2\}$,
\begin{equation*}
\begin{gathered}
\limsup_n\|E_{n,j}^{(r)}\|_{L^2_x}
\leq C\ell|K_j^+|^{1/2},\\
\limsup_n\|E_{n,j}^{(r)}\|_{L^{p_*}_x}
\leq C_{p_*}|K_j^+|^{1/p_*}.
\end{gathered}
\end{equation*}
The second bound follows from the fixed-wide multiplier estimates and
\eqref{eq:var-local-volume}. Choose the interpolation exponent by
\begin{equation*}
\f13=\f{\alpha_*}{2}+\f{1-\alpha_*}{p_*},
\qquad \alpha_* = \f{2(p_*-3)}{3(p_*-2)}.
\end{equation*}
H\"older's inequality then yields
\begin{equation}\label{eq:var-projection-error}
\begin{aligned}
\limsup_n\|E_{n,j}^{(r)}\|_{L^3_x}
&\leq\limsup_n
 \|E_{n,j}^{(r)}\|_{L^2_x}^{\alpha_*}
 \|E_{n,j}^{(r)}\|_{L^{p_*}_x}^{1-\alpha_*}\\
&\leq C_{p_*}\ell^{\alpha_*}
 |K_j^+|^{\alpha_*/2+(1-\alpha_*)/p_*}
 =C_{p_*}\ell^{\alpha_*}|K_j^+|^{1/3}.
\end{aligned}
\end{equation}
This argument uses defect compression for the $L^2$ estimate and
finite-$L^{p_*}$ bounds for the upgrade. It does not claim an $O(\ell)$
$L^3$ error from quadratic localization alone.

For sufficiently small $\ell$ the frozen thin-cone cutoff equals one
on every actual null direction over $K_j^+$. Its complementary defect
is then zero, and its gradient complement is strongly zero in $L^3$
by the same fixed $L^2/L^{p_*}$ interpolation, with constants allowed to
depend on $\delta$. Set $v_{n,j}=S_j(D)u_{n,j}$, where $S_j$ is the frozen screen and
thin-cone symbol $S$ in \eqref{eq:var-full-assembly} for $m=m_j$.
Combining
\eqref{eq:var-projection-error}, these facts, and fixed
multiplier bounds gives
\begin{equation*}
\begin{gathered}
\limsup_n\|\partial(u_{n,j}-v_{n,j})\|_{L^3_x}
\leq C_{\delta,p_*}\ell^{\alpha_*}|K_j^+|^{1/3},\\
\limsup_n\left\|\f{q_{m_j}(D)}{|D|^2}
                            \partial v_{n,j}\right\|_{L^3_x}
\leq C_{\delta,p_*}\ell^{\alpha_*}|K_j^+|^{1/3}.
\end{gathered}
\end{equation*}
The bounded-frequency regularizations add $o_n(1)$ only. The complete
$W^{1,3}$ continuity estimate therefore permits replacement of every
physical factor at cost $C_{\delta,p_*}\ell^{\alpha_*}|K_j^+|$ in the limsup.
Every phase term in the Ward classification has the same bound by
the all-scale normalized phase estimate proved above.

Estimate the three-physical near part by
\eqref{eq:var-near-bound}, with the fixed-wide screen projection on
the right. Its limsup is at most
$C_{\kappa,p_*}\delta^{\beta_*}|K_j^+|$. In particular the constants in this
bound are uniform in $\ell$ on the compact smooth family of frozen
metrics. Projection errors have already been treated at fixed
$\delta$; a possibly large norm of the total thin-cone multiplier
does not appear in this near estimate.

For both corrections use the actual global scalar coefficient
\begin{equation*}
a_{n,j}^{ab}=m_j^{ab}+d_{n,j}^{ab},\qquad
d_{n,j}^{ab}=\chi_j^+(g_n^{ab}-m_j^{ab}),\qquad
P_{n,j}=-\partial_a(a_{n,j}^{ab}\partial_b).
\end{equation*}
The coefficient $a_{n,j}$ equals the inverse metric near
$\supp z_{n,j}$. Proposition~\ref{prop:general-projection} gives a global
source bound, including the projection tails.
The projection annihilates the symmetrized derivative of the lowered
gauge source; its commutator uses only the coefficient's Lipschitz norm.
In particular the equation is applied to $u_{n,j}$ itself, while every
additional screen or cone projection remains in the correction.
By \eqref{eq:general-explicit-source} and the scaled cutoff bounds,
for $\|h_n\|_{L^\infty(K^+)}\leq1$,
\begin{equation*}
\|P_{n,j}u_{n,j}\|_{L^2_x}
\leq C|K_j^+|^{1/2}
 \bigl(1+\ell^{-1}+\ell^{-2}\|h_n\|_{L^\infty(K^+)}\bigr).
\end{equation*}
Here $C$ is independent of $j,n,\ell$. One cutoff derivative multiplies
a bounded gauge source or first derivative; two multiply $h_n$.
The global commutator and divergence drift obey the same bound, since
\begin{equation*}
\bigg(1+\f{\|h_n\|_{L^\infty(K^+)}}{\ell}\bigg)^2
\leq C\bigl(1+\ell^{-1}+\ell^{-2}\|h_n\|_{L^\infty(K^+)}\bigr).
\end{equation*}
The extensions satisfy
\begin{equation*}
\begin{gathered}
\|d_{n,j}\|_{L^\infty_x}
 \leq C\bigl(\ell+\|h_n\|_{L^\infty(K^+)}\bigr),\\
\|\partial d_{n,j}\|_{L^\infty_x}
 \leq C\bigg(1+\f{\|h_n\|_{L^\infty(K^+)}}\ell\bigg).
\end{gathered}
\end{equation*}
For the physical and gauge corrections, let $H_{n,j}$ denote the
corresponding low tensor in \eqref{eq:var-correction-energy}.
After summing its angular expansion, the multiplier bounds give
\begin{equation*}
\begin{gathered}
\|H_{n,j}\|_{L^4_x}
 \leq C_\delta\|h_n\|_{L^\infty(K^+)}|K_j^+|^{1/4},\\
\|\partial u_{n,j}\|_{L^p_x}+\|\partial H_{n,j}\|_{L^p_x}
 \leq C_{\delta,p}
 \bigg(1+\f{\|h_n\|_{L^\infty(K^+)}}\ell\bigg)|K_j^+|^{1/p},
 \qquad 1<p<\infty.
\end{gathered}
\end{equation*}
For each correction the three terms of
\eqref{eq:var-correction-energy} carry the same volume factor
$|K_j^+|^{1/2}|K_j^+|^{1/4}|K_j^+|^{1/4}=|K_j^+|$.
Denote their complete quadratic forms by
\begin{equation*}
\mathcal C_{n,j}^{\sigma}
 =\bigl\langle
 \ii[P_{0,j},T_{\sigma,n,j}]u_{n,j},u_{n,j}\bigr\rangle_{L^2(\R^4)},
\qquad \sigma\in\{F,G\},
\end{equation*}
where $P_{0,j}=-m_j^{ab}\partial_a\partial_b$ and
$T_{F,n,j}$, $T_{G,n,j}$ are the actual two corrections from
\eqref{eq:var-full-assembly}, with frozen metric $m_j$.
Substitution in \eqref{eq:var-correction-energy}, followed by
$\sum_j|K_j^+|\leq C|K^+|$, gives
\begin{equation*}
\begin{aligned}
\sum_j\sum_{\sigma\in\{F,G\}}|\mathcal C_{n,j}^{\sigma}|
&\leq C_\delta\sum_j|K_j^+|
 \bigg[\|h_n\|_{L^\infty(K^+)}
 \bigl(1+\ell^{-1}+\ell^{-2}\|h_n\|_{L^\infty(K^+)}\bigr)\\
&\hspace{42mm}\times
 \bigg(1+\f{\|h_n\|_{L^\infty(K^+)}}\ell\bigg)\\
&\hspace{19mm}
 +\bigl(\ell+2\|h_n\|_{L^\infty(K^+)}\bigr)
 \bigg(1+\f{\|h_n\|_{L^\infty(K^+)}}\ell\bigg)^3\bigg]\\
&\leq C_\delta|K^+|
 \bigg[\|h_n\|_{L^\infty(K^+)}
 \bigl(1+\ell^{-1}+\ell^{-2}\|h_n\|_{L^\infty(K^+)}\bigr)\\
&\hspace{42mm}\times
 \bigg(1+\f{\|h_n\|_{L^\infty(K^+)}}\ell\bigg)\\
&\hspace{19mm}
 +\bigl(\ell+2\|h_n\|_{L^\infty(K^+)}\bigr)
 \bigg(1+\f{\|h_n\|_{L^\infty(K^+)}}\ell\bigg)^3\bigg].
\end{aligned}
\end{equation*}
At fixed $\ell,\delta$ its limsup is at most
$C_\delta\ell|K^+|$. In particular all negative powers of $\ell$
in the source term disappear before the spatial scale changes.
The constants have no dependence on the number of patches, the
angular truncation, or the number of radial scales.

Finally bounded overlap gives $\sum_j|K_j^+|\leq C|K^+|$.
Sum all near, phase, correction, replacement, and freezing bounds
after taking the sequence limit. Equation~\eqref{eq:var-cubic-residual}
then gives
\begin{equation}\label{eq:var-final-bound}
\left|4\int_{S_r^*M}H_Pa\,\dd\mu\right|
\leq C_{\kappa,p_*}\delta^{\beta_*}
   +C_{\delta,p_*}\ell^{\alpha_*}+C_\delta\ell.
\end{equation}
All constants may depend on the fixed compact set, the smooth limit,
the test, and the original first-derivative bound. The first constant
is uniform in $\ell$, while the other constants may depend on the
fixed $\delta$. Every quadratic limit used here is a test of the
single defect measure already extracted in Section~\ref{sec:quadratic}.
Changing the fixed test or the fixed localization parameters does not
require another subsequence. Individual cubic summands need not converge:
their limsup bounds suffice because the total residual has the limit in
\eqref{eq:var-cubic-residual}. Thus the left side of
\eqref{eq:var-final-bound} is the same number for every allowed fixed
$\ell,\delta$. In that numerical inequality let $\ell\to0$, then
$\delta\to0$, with $\kappa$ fixed. This proves the local
transport identity for odd real
tests. Scalar parity and complex linearity give all degree-one tests,
and a base partition together with the invariant construction gives
the stated global identity. The full-momentum conclusion follows
from Lemma~\ref{lem:lift} after this transport step.
\end{proof}

The strengthened exponents do not give a convergence rate in $n$:
the hypotheses prescribe no such rate. They improve the auxiliary
angular and spatial errors after the sequence limit. Every
$0<\beta_*<2$ is available, while $\alpha_*\uparrow2/3$ as
$p_*\to\infty$; neither endpoint is claimed, because the finite-exponent
multiplier constants are not controlled at infinity.

The order of limits is $n\to\infty$, then $\ell\to0$, then
$\delta\to0$. There is no sequence-dependent cone width, largest
frequency, or second-derivative assumption. The source estimates used here come from the physical projection and
require no coordinate condition. The normal forms act on the scalar
divergence equation for the projected physical field, with the original
inverse-metric coefficient still uniformly Lipschitz. The density
Hessian is used for the variational residual and its smooth-background
commutator.

\subsection{Full-momentum realization and completion}

\begin{lemma}[Local-moment version of the globally hyperbolic lift]\label{lem:lift}
Suppose a smooth globally hyperbolic spacetime has a nonnegative locally finite measure on $S_r^*M$, supported on $P=0$, satisfying both \eqref{eq:moment} and \eqref{eq:transport}. Then it has a full null cotangent representation $\nu$ with $\int_{\pi^{-1}(K)}|\xi|_r^2\,\dd\nu<\infty$ for every $K\Subset M$.
\end{lemma}
\begin{proof}
This is the construction of \cite[Lemma 2.5]{HLelliptic} with its local
finiteness made explicit. Fix a smooth spacelike Cauchy hypersurface
$\Sigma$. Every nonzero null Hamilton orbit meets $T^*M|_\Sigma$ exactly
once along its inextendible projected geodesic. Denote its covector at
that crossing by $\xi_\Sigma$, and put
$\mathcal E(x,\xi)=|\xi_\Sigma|_r$. Transversality to $\Sigma$ and
smooth dependence of the geodesic equation make $\mathcal E$ smooth
and positive on the nonzero null cone. Hamilton-flow scaling gives
$\mathcal E(x,t\xi)=t\mathcal E(x,\xi)$ for $t>0$, and its definition
gives $H_P\mathcal E=0$. Thus
$s(x,\omega)=\mathcal E(x,\omega)^{-1}$ defines a smooth positive scale
on the null unit-ray bundle, and the section
$L(x,\omega)=(x,s(x,\omega)\omega)$ has invariant image
$\{\mathcal E=1\}$. No affine completeness or compactness of
$\Sigma$ is used.

Define $\nu=L_*(s^{-2}\mu)$. This is a nonnegative Radon measure: $s$ and its reciprocal are bounded on each compact base unit-null bundle. Its second moments equal those of $\mu$, and, exactly,
\begin{equation*}
\int_{\pi^{-1}(K)}|\xi|_r^2\,\dd\nu=\mu(\pi^{-1}(K)\cap S_r^*M)<\infty.
\end{equation*}
For $b\in C_c^\infty(T^*M\setminus0)$, define the degree-one
extension on the null cone by
\begin{equation*}
a(x,\rho\omega)=\f{\rho}{s(x,\omega)}
 b\bigl(x,s(x,\omega)\omega\bigr),
\qquad \rho>0,\quad |\omega|_r=1,\quad P(x,\omega)=0.
\end{equation*}
A smooth angular extension off the null cone preserves compact base
support. A partition of unity extends this same function on the cone.
Since $a=b$ on the invariant section $L$, tangency and homogeneity give
\begin{equation*}
\int_{T^*M\setminus 0}H_Pb\,\dd\nu
=\int_{S_r^*M}s^{-2}(H_Pa)\circ L\,\dd\mu
=\int_{S_r^*M}H_Pa\,\dd\mu=0.
\end{equation*}
One can also verify the weight directly in radial coordinates on the
null cone. Write $\xi=\rho\omega$, $|\omega|_r=1$, and
$H_P=\rho X+\rho^2 k\partial_\rho$, where $X$ is tangent to the
unit-null bundle and $k$ is smooth there. Tangency of $\rho=s$ means
$Xs=ks$. For $B=b\circ L$ and $f=B/s$, the homogeneous extension
$a=\rho f$ therefore satisfies
\begin{equation*}
 (H_Pa)|_{\rho=1}=Xf+kf=s^{-1}XB,
 \qquad (H_Pb)\circ L=sXB.
\end{equation*}
These formulas recover exactly the factor $s^{-2}$ and include the
radial drift caused by the varying unit normalization.
This proves full transport and all local moment requirements.
\end{proof}

Lemma~\ref{lem:lift} is conditional on transport. Applying it to Theorem~\ref{thm:quadratic} alone would be invalid. The formal radial lift $\rho^{-3}\dd\rho\,\dd\mu$ has divergent second moment $\int_0^\infty\dd\rho/\rho$ and cannot replace this construction on a general spacetime.


\begin{proof}[Proof of Theorem~\ref{thm:main}]
Apply Theorem~\ref{thm:quadratic} to obtain one global positive null
measure $\mu$, with the moment identity and $\Scal(g)=0$.
Theorem~\ref{thm:variational-transport} establishes transport for this
same measure. Its proof uses the same subsequence and does not choose
new measures for different test functions or charts. For a test whose
base support meets several charts, decompose the test by a smooth base
partition of unity. Since the resulting tests sum to the original one,
their Hamilton derivatives also sum to its Hamilton derivative.
This proves the global equation \eqref{eq:transport}.

Restrict $\mu$ to a globally hyperbolic open subspacetime $U$.
Lemma~\ref{lem:lift} supplies $\nu$, proves
\eqref{eq:full-moment-bound}, and yields both equations in
\eqref{eq:full-equations}. These are all the assertions of the theorem.
\end{proof}

\begin{corollary}[Vacuum limits and strong physical compactness]
\label{cor:strong}
Assume \eqref{eq:admissible}, and let $U\subset M$ be open. Then
$\Ric(g)=0$ on $U$ if and only if, for every coordinate chart in $U$
and every $\chi\in C_c^\infty$ of that chart,
\begin{equation}\label{eq:physical-strong-criterion}
\|\partial\Pi_E(D)(\chi h_n)\|_{L^2(\R^4)}\longrightarrow0.
\end{equation}
Here the compactly supported tensor $\chi h_n$ is extended by zero, and $\Pi_E$ is
the homogeneous Euclidean gauge-complement projection from
Proposition~\ref{prop:general-projection}; its value at frequency zero
is immaterial. The convergence holds for the entire sequence.
In particular, $Dh_n\to0$ strongly in $L^2_{\loc}(U)$ implies
$\mu|_{S_r^*U}=0$ and $\Ric(g)=0$ on $U$. The converse need only hold after the physical
projection: pure coordinate oscillations may remain.
\end{corollary}
\begin{proof}
Suppose first that $\Ric(g)=0$ on $U$. By
Corollary~\ref{cor:base-marginal}, every subsequential physical measure
vanishes over $U$. Off the null cone the Ricci-symbol kernel is pure
gauge and is annihilated by $\Pi_E$. On the null cone, $\Pi_E$ identifies
the quotient by gauge with a complementary subspace on which $F$ in
\eqref{eq:B-null} is positive definite. On a compact unit-null bundle,
finite-dimensional compactness gives a uniform constant $c>0$ such that
\begin{equation*}
F(k,k)=F(\Pi_E k,\Pi_E k)\geq c|\Pi_E k|_E^2
\quad\text{for }k\in\ker\mathcal R_\xi.
\end{equation*}
Consequently the projected tensor-gradient defect is zero everywhere
over $\supp\chi$.

There is no loss at spatial infinity from the nonlocal projection.
Since $\Pi_E^*\Pi_E=\Pi_E$ and it commutes with derivatives, Parseval gives
\begin{equation*}
\|\partial\Pi_E(D)(\chi h_n)\|_{L^2(\R^4)}^2
=\sum_{a=0}^3\bigl\langle\Pi_E(D)\partial_a(\chi h_n),
                         \partial_a(\chi h_n)\bigr\rangle_{L^2(\R^4)}.
\end{equation*}
The second factor has fixed compact support. Its derivatives of $\chi$
multiply $h_n$ and converge strongly to zero; the remaining limit is
exactly the vanishing projected defect just computed. Every subsequence
has a further subsequence with this zero limit, proving
\eqref{eq:physical-strong-criterion} for the full sequence.

Conversely, \eqref{eq:physical-strong-criterion} annihilates the physical
part of every local derivative defect. Gauge invariance of $F$ then
gives $\mu=0$ on $U$, and the quadratic moment identity gives
$\Ric(g)=0$ there. Taking strong convergence of all derivatives is a
sufficient special case. This proof uses only
Theorem~\ref{thm:quadratic}, and is independent of the transport argument.
\end{proof}

\appendix
\section{A harmonic-coordinate example with an unbounded source}\label{sec:shear}

In this appendix $m=\operatorname{diag}(-1,1,1,1)$ is the Minkowski
metric, and $\Box_m=m^{ab}\partial_a\partial_b$.

\begin{proposition}[Failure of a background $L^2$ source estimate]\label{prop:shear}
There are globally smooth flat metrics $g_n$ on $\R^4$, satisfying \eqref{eq:admissible} with $g=m$, and satisfying $H(g_n)=H(m)=0$ exactly, for which $\|\Box_mh_n\|_{L^2(K)}\to\infty$ on a fixed compact set $K$.
\end{proposition}
\begin{proof}
In Minkowski coordinates $y$, put $\varepsilon=n^{-1}$, $\delta=n^{-3}$,
\begin{equation*}
f(p)=\varepsilon^2\sin(p/\varepsilon),\quad j(q)=\delta^2\sin(q/\delta),
\end{equation*}
and define the harmonic coordinate map
\begin{equation*}
x^0=y^0,\quad x^3=y^3,\quad
x^1=y^1+f(y^0-y^3),\quad x^2=y^2+j(y^0-y^1).
\end{equation*}
Each $x^a(y)$ is Minkowski harmonic. The triangular map has determinant one and a global inverse. With $t=x^0$, $z=x^3$, $p=t-z$ and $s=t-x^1+f(p)$, that inverse is
\begin{equation*}
y^0=t,\quad y^3=z,\quad y^1=x^1-f(p),\quad y^2=x^2-j(s).
\end{equation*}
Set $g_n=(y(x))^*m$. It is flat and has $H(g_n)=0$, since the inverse coordinate functions $x^a(y)$ are harmonic. Moreover
\begin{equation*}
g_n=-\dd t^2+\bigl(\dd x^1-f'(p)\dd p\bigr)^2
+\bigl(\dd x^2-j'(s)\dd s\bigr)^2+\dd z^2.
\end{equation*}
First derivatives of $f,j$ have sizes $\varepsilon,\delta$, and their second derivatives are bounded by one. The inverse Jacobian tends uniformly to the identity; its first derivative is uniformly bounded, including the term $j' f''$. Thus $g_n\to m$ uniformly and its first derivatives are uniformly bounded.

The $21$ component is exactly $h_{n,21}=j'(s)$. Directly,
\begin{equation*}
m^{-1}(\dd s,\dd s)=-2f'(p),\quad \Box_ms=0,
\end{equation*}
so
\begin{equation*}
\Box_mh_{n,21}=j'''(s)(-2f'(p))
=\f{2\varepsilon}{\delta}\cos(p/\varepsilon)\cos(s/\delta).
\end{equation*}
Use coordinates $(t,x^2,p,q)$, with $q=t-x^1$, whose Jacobian
has absolute determinant one. Let $K$ be the image of
$[0,1]^2\times[0,2\pi]^2$ under this fixed linear coordinate change.
Since $\varepsilon^{-1}=n$ and $\delta^{-1}=n^3$ are integers,
\begin{equation*}
\begin{aligned}
\|\Box_mh_{n,21}\|_{L^2(K)}^2
&=\f{4\varepsilon^2}{\delta^2}
 \int_0^1\int_0^1\int_0^{2\pi}\int_0^{2\pi}
 \cos^2(p/\varepsilon)\cos^2((q+f(p))/\delta)
 \,\dd q\,\dd p\,\dd x^2\,\dd t\\
&=\f{4\varepsilon^2}{\delta^2}\pi^2=4\pi^2n^4.
\end{aligned}
\end{equation*}
This proves the claim on a fixed compact parallelepiped, and hence on a larger compact coordinate box.
\end{proof}

The limit has the kinetic representation $\mu=0$. Thus the bound \cite[(3.16)]{HLwave} need not hold after its derivative-scale hypothesis is removed; it does not refute Burnett's conjecture. In particular it explains why merely assuming an exact wave source does not restore the hypothesis of that paper's Lemma~5.2.

\section{Characteristic localization and cubic interactions}\label{sec:cubic}

\begin{proposition}[An auxiliary trilinear counterexample]\label{prop:cubic}
There are smooth scalar $f_n,w_n\to0$ locally uniformly with uniformly bounded first derivatives and $\Box_m f_n,\Box_mw_n\to0$ strongly in $H^{-1}_{\loc}(\R^4)$ such that
\begin{equation*}
Q_0(\dd f_n,\dd w_n)\partial_tf_n\rightharpoonup-\f14
\end{equation*}
in distributions. Here $Q_0(\dd f,\dd w)=m^{ab}\partial_af\partial_bw$.
\end{proposition}
\begin{proof}
In coordinates $(t,x,y,z)\in\R^4$, set $u=t-x$, $v=t+x$,
$\lambda_{1,n}=n^2$, $\lambda_{2,n}=n$, and
\begin{equation*}
f_n=\lambda_{1,n}^{-1}\cos(\lambda_{2,n}v)\sin(\lambda_{1,n}u),\qquad
w_n=(2\lambda_{2,n})^{-1}\sin(2\lambda_{2,n}v).
\end{equation*}
The asserted amplitude and gradient bounds follow by differentiation. Since $\Box_m=-4\partial_u\partial_v$ on these functions,
\begin{equation*}
\Box_mf_n=4\lambda_{2,n}\sin(\lambda_{2,n}v)\cos(\lambda_{1,n}u),\qquad \Box_mw_n=0.
\end{equation*}
Each localized Fourier mode of the first expression has frequency comparable to $\lambda_{1,n}$, hence its $H^{-1}(\R^4)$ norm is $O(\lambda_{2,n}/\lambda_{1,n})$. The tails introduced by a fixed smooth cutoff decay faster than every power. Cutoff commutators are bounded in $L^2(\R^4)$ and weakly zero, hence strongly zero in $H^{-1}_{\loc}(\R^4)$. For every fixed compact localization, all fixed off-null conic projections of the localized gradients therefore tend strongly to zero in $L^2(\R^4)$.

Nevertheless
\begin{equation*}
Q_0(\dd f_n,\dd w_n)=-2\cos(\lambda_{2,n}v)\cos(\lambda_{1,n}u)\cos(2\lambda_{2,n}v),
\end{equation*}
Differentiation also gives
\begin{equation*}
\partial_tf_n=\cos(\lambda_{2,n}v)\cos(\lambda_{1,n}u)
 +O(\lambda_{2,n}/\lambda_{1,n}).
\end{equation*}
Thus the trilinear product equals
\begin{equation*}
-2\cos^2(\lambda_{2,n}v)\cos(2\lambda_{2,n}v)\cos^2(\lambda_{1,n}u)+o(1).
\end{equation*}
The means of the $v$ and $u$ factors are $1/4$ and $1/2$. Every other trigonometric mode has frequency tending to infinity, and vanishes against a fixed compact smooth test by integration by parts. The limit is $-1/4$.
\end{proof}

These scalar functions have not been realized as exact Einstein solutions. The proposition rules out a generic null-form lemma based only on quadratic characteristic localization. The coupled Einstein equations and their entire physical contraction must supply information beyond those hypotheses.

\section{Conservation and positive null moments do not give transport}\label{sec:rotating}

On Minkowski space fix $c>0$ and $v(t)=(\cos t,\sin t,0)$. Define
\begin{equation*}
T_{00}=2c,\quad T_{0i}=0,\quad T_{ij}=2c v_i(t)v_j(t).
\end{equation*}
This tensor is smooth, positive semidefinite as an ordinary quadratic form, trace-free in $m$, and divergence-free. Pointwise it is the sum $c(k_+^2+k_-^2)$ with $k_\pm=(-1,\pm v(t))$ null. Nevertheless it has no transported positive null ray representation on a time interval about zero.

To prove this, fold a proposed measure to the future sheet. With Euclidean unit cotangent normalization, all null covectors have $\xi_0=-1/\sqrt2$. Vanishing moments on spatial vectors perpendicular to $v$ force $\xi_{\rm sp}\parallel v$ almost everywhere, by positivity. Thus only $\omega_\pm=(-1,\pm v)/\sqrt2$ remain. The $00$ and $0i$ moments force each branch to have mass $2c\dd t\dd x$. These are measure identities: the $00$ identity fixes the base marginal, after which the mixed identities fix both conditional branch weights.

Fix an interval $I$ about zero and nonnegative
$\chi\in C_c^\infty(I)$, $\eta\in C_c^\infty(\R^3)$.
Choose a smooth degree-one angular extension of
$F(\xi)=\xi_1\xi_2/(-\xi_0)$ from the future null shell and set
$a=\chi\eta F$. The spatial terms cancel between the two branches, so
\begin{equation*}
\begin{aligned}
\int_{S_r^*(I\times\R^3)}H_Pa\,\dd\mu
&=2c\bigg(\int_{\R^3}\eta(x)\,\dd x\bigg)
 \int_I\chi'(t)\cos t\sin t\,\dd t\\
&=-2c\bigg(\int_{\R^3}\eta(x)\,\dd x\bigg)
 \int_I\chi(t)\cos(2t)\,\dd t.
\end{aligned}
\end{equation*}
It is nonzero for nonnegative nonzero cutoffs near $t=0$. Folding is legitimate because the antipodal map preserves quadratic moments and reverses the Hamilton field, so a zero transport integral remains zero. This contradicts transport.

Here $T$ is not $G(m)$, and no vacuum sequence with this tensor as its effective stress has been constructed. The example establishes precisely the failure of the proposed implication from tensor conservation and all pointwise positivity properties to kinetic transport.

\subsection{A smooth metric whose actual Einstein tensor is not kinetic}

The preceding obstruction can be realized by the Einstein tensor of a smooth metric. The missing condition remains vacuum approximation, so the following proposition is still not a vacuum-limit counterexample to Theorem~\ref{thm:main}.

\begin{proposition}[Geometric nonrepresentability without vacuum approximation]\label{prop:bianchi}
There exists a smooth Lorentzian metric on a product interval times $\R^3$ whose Einstein tensor is trace-free, positive semidefinite, and pointwise a positive sum of two null squares, but which admits no nonnegative null ray measure satisfying both its Einstein identity and Hamilton transport.
\end{proposition}
\begin{proof}
Let $p(t)=(\cos t,\sin t,0)^T$. For symmetric positive definite $\gamma$ and symmetric $A$, put
\begin{equation*}
K=\gamma^{-1}A,\quad\theta=\tr K,\quad d=p^T\gamma^{-1}p,\quad
W=\f{\gamma^{-1}pp^T}{d}.
\end{equation*}
Fix $H>0$ and solve the smooth finite-dimensional system
\begin{equation*}
\begin{gathered}
\gamma'=2A,\\
A'=2A\gamma^{-1}A-\theta A+\rho\f{pp^T}{d},\\
\rho'=-\rho\bigl(\theta+\tr(KW)\bigr),\\
\gamma(0)=I,\quad A(0)=HI,\quad\rho(0)=3H^2.
\end{gathered}
\end{equation*}
The vector field preserves symmetry and is smooth where $\gamma$ is positive definite. Ordinary local existence supplies an interval $|t|<\tau<\pi/4$ with $\gamma>0$ and $\rho>0$. On this interval $K'=\rho W-\theta K$ and $\theta'=\rho-\theta^2$, since $\tr W=1$. Direct differentiation gives
\begin{equation*}
\f{\dd}{\dd t}\bigl(\theta^2-\tr(K^2)-2\rho\bigr)
=-2\theta\bigl(\theta^2-\tr(K^2)-2\rho\bigr).
\end{equation*}
Its initial value is zero, so the constraint persists.

For $g=-\dd t^2+\gamma_{ij}(t)\dd x^i\dd x^j$, the nonzero connection coefficients are $\Gamma^0_{ij}=A_{ij}$ and $\Gamma^i_{0j}=K^i{}_j$. Inserting these in the Ricci formula yields
\begin{equation*}
\Ric_{00}=-\theta'-\tr(K^2)=\rho,\quad
\Ric_{0i}=0,\quad
\Ric_{ij}=\bigl[\gamma(K'+\theta K)\bigr]_{ij}=\rho\f{p_ip_j}{d}.
\end{equation*}
The scalar curvature is $-\rho+\rho\tr W=0$, so these are also the components of $G$. With $k_\pm=(-\sqrt d,\pm p)$,
\begin{equation*}
G=\f{\rho}{2d}\bigl(k_+\otimes k_++k_-\otimes k_-\bigr),
\end{equation*}
which proves the stated pointwise properties. The usual Bianchi identity also gives divergence zero.

Suppose a transported positive null ray representation existed, and fold it to the future sheet. Spatial vectors perpendicular in the ordinary Euclidean sense to $p$ have zero $G$ moment. Nonnegativity therefore forces $\xi_{\rm sp}\parallel p(t)$ almost everywhere. In particular $\xi_1\ne0$ on the support, since $\cos t>0$ and a nonzero null covector cannot have zero spatial part.

Set $I=(-\tau,\tau)$ and $M=I\times\R^3$.
For nonnegative nonzero $\chi\in C_c^\infty(I)$ and
$\eta\in C_c^\infty(\R^3)$, choose a degree-one test equal near
the future null shell to
\begin{equation*}
a(t,x,\xi)=\chi(t)\eta(x)
\f{\xi_1\xi_2-\tan(t)\xi_1^2}{-\xi_0}.
\end{equation*}
A smooth homogeneous angular cutoff extends it to the entire nonzero cotangent bundle. On the forced support its numerator vanishes, as do all terms obtained by differentiating the base cutoffs and the denominator. Spatial translation invariance of $g$ means the Hamilton field has no $\partial_{\xi_i}$ components. Therefore, on this support,
\begin{equation*}
H_Pa=-\chi\eta\sec^2(t)\xi_1^2.
\end{equation*}
Using the $11$ Einstein moment, the transport integral is exactly
\begin{equation*}
\int_{S_r^*M}H_Pa\,\dd\mu=-\int_{I\times\R^3}\chi(t)\eta(x)\f{\rho(t)}{d(t)}\dd V_g<0,
\end{equation*}
a contradiction. This excludes every positive representation, including those initially using both time sheets.
\end{proof}

Proposition~\ref{prop:bianchi} shows that the second-moment conclusion of Theorem~\ref{thm:quadratic}, even together with an actual smooth Einstein tensor and the Bianchi identity, does not imply the missing transport. No sequence satisfying \eqref{eq:admissible} and converging to the metric of Proposition~\ref{prop:bianchi} has been constructed.

\section{The Rosen plane-wave class}\label{sec:rosen}

\begin{proposition}[Two-polarization Rosen closure]\label{prop:rosen}
Let $I\subset\R$ be an interval and
$g_n=-2\dd u\dd v+\gamma_{n,AB}(u)\dd y^A\dd y^B$ on
$M=I\times\R^3$, with $A,B\in\{1,2\}$ and $\gamma_n$ smooth
positive definite on $I$. Suppose \eqref{eq:admissible} holds with smooth positive definite $\gamma$. Then $g$ has a full massless kinetic representation, without any second-derivative scale assumption.
\end{proposition}
\begin{proof}
Put $L_n=\gamma_n^{-1}\gamma_n'/2$ and $\tau_n=\tr L_n$. The nonzero connection coefficients are
\begin{equation*}
\Gamma^A_{uB}=(L_n)^A{}_B,\qquad \Gamma^v_{AB}=\gamma'_{n,AB}/2.
\end{equation*}
The only possibly nonzero Ricci component is
\begin{equation*}
\Ric(g_n)_{uu}=-\tau_n'-\tr(L_n^2).
\end{equation*}
Define the real symmetric matrices
\begin{equation*}
S_n=\f12\gamma_n^{-1/2}\gamma_n'\gamma_n^{-1/2},\qquad
S=\f12\gamma^{-1/2}\gamma'\gamma^{-1/2}.
\end{equation*}
Uniform convergence of $\gamma_n$ and the derivative bound imply
\begin{equation*}
S_n\rightharpoonup^* S\quad\hbox{in }L^\infty_{\loc}(I),\qquad
-\tau_n'=\tr(L_n^2)=|S_n|_F^2.
\end{equation*}
Pass to a weak-* limit $|S_n|_F^2\rightharpoonup^*Q$ in
$L^\infty_{\loc}(I)$. The vacuum identity gives $Q=-\tau'$.
For $0\leq\varphi\in C_c^\infty(I)$, weak lower semicontinuity yields
\begin{equation*}
\int_I\varphi(u)|S(u)|_F^2\,\dd u
\leq\liminf_n\int_I\varphi(u)|S_n(u)|_F^2\,\dd u
=\int_I\varphi(u)Q(u)\,\dd u.
\end{equation*}
Consequently
\begin{equation*}
\Ric(g)=\rho(u)\dd u^2,\qquad \rho=-\tau'-|S|_F^2\geq0.
\end{equation*}
Smoothness of the limit makes $\rho$ smooth. At $\xi=\dd u$, one has $P=0$ and $H_P=-\partial_v$, with no vertical component. Thus
\begin{equation*}
\nu=\rho(u)\dd V_g(x)\,\delta_{\dd u(x)}(\dd\xi)
\end{equation*}
has the required second moment. For $b\in C_c^\infty(T^*M\setminus0)$,
integration in $v$ gives
\begin{equation*}
\int_{T^*M\setminus0}H_Pb\,\dd\nu
=-\int_{I\times\R^3}\rho(u)
 \partial_vb(x,\dd u(x))\sqrt{\det\gamma(u)}
 \,\dd u\,\dd v\,\dd y^1\,\dd y^2=0.
\end{equation*}
The density and $r$-length of $\dd u$ are bounded on compact sets,
so all local second moments are finite.
\end{proof}

For an explicit nonzero example take $\lambda_n=e^n$,
\begin{equation*}
\psi_n(u)=\f{\sin(nu)}n+\f{\sin(\lambda_n u)}{\lambda_n},\quad
B_n''+(\psi_n')^2B_n=0,\quad B_n(0)=1,\ B_n'(0)=0.
\end{equation*}
On $I=(-1/4,1/4)$, comparison with $\cos(2u)$ gives $B_n\geq\cos(2u)>0$, and $|B_n'|,|B_n''|$ are uniformly bounded. Since $(\psi_n')^2\rightharpoonup^*1$ in $L^\infty(I)$,
compactness in $C^1(\overline I)$ and uniqueness for the limiting ODE
give $B_n\to\cos u$ in $C^1(\overline I)$. Hence the exact vacuum metrics
\begin{equation*}
g_n=-2\dd u\dd v+B_n^2(e^{2\psi_n}\dd x^2+e^{-2\psi_n}\dd z^2)
\end{equation*}
converge to $g=-2\dd u\dd v+\cos^2u(\dd x^2+\dd z^2)$, with $\Ric(g)=2\dd u^2$ and $\nu=2\dd V_g\delta_{\dd u}$.

Fix a closed interval $I_0\Subset I$ containing zero in its interior.
The screen ratio $g_{n,xx}/g_{n,zz}=e^{4\psi_n}$, evaluated at $u=\pi/(2n)$, gives $\|g_n-g\|_{L^\infty(I_0\times[-1,1]^3)}\geq c/n$ for all sufficiently large $n$. At $u=\pi/(2\lambda_n)$, $|\psi_n''|\geq \lambda_n$, whereas the remaining terms in $\partial_u^2g_{n,xx}$ are bounded. Thus $\|\partial^2g_n\|_{L^\infty(I_0\times[-1,1]^3)}\geq c\lambda_n-C$. No common amplitude parameter can satisfy the pair of estimates in \eqref{eq:literature-scale}.

This example also lies in an exact zero-source gauge. Set $v=B_n(u)V$ and pull back to the fixed product coordinates $(u,V,x,z)$:
\begin{equation*}
\widetilde g_n=-2B_n\dd u\dd V-2B_n'V\dd u^2
+B_n^2(e^{2\psi_n}\dd x^2+e^{-2\psi_n}\dd z^2).
\end{equation*}
Its volume density is $B_n^3$ and $\widetilde g_n^{uV}=-B_n^{-1}$, $\widetilde g_n^{VV}=2B_n'V/B_n^2$, $\widetilde g_n^{uu}=0$. Therefore
\begin{equation*}
\partial_u(B_n^3\widetilde g_n^{Vu})+\partial_V(B_n^3\widetilde g_n^{VV})
=-2B_nB_n'+2B_nB_n'=0,
\end{equation*}
and every other contracted source is zero as well. Convergence and the first-derivative bound persist because $B_n\to B$ in $C^1(\overline I)$ and $B_n''$ is bounded. The same ratio and second-derivative arguments rule out a fixed scale. The limit has $H_P=-B^{-1}\partial_V$ at $\dd u$, so $2\dd V_{\widetilde g}\delta_{\dd u}$ again gives full transport.

\subsection{Rigidity of flat sequences}

\begin{proposition}[Flat admissible sequences have flat smooth limits]\label{prop:flat-rigidity}
If, in addition to \eqref{eq:admissible}, the full curvature of every $g_n$ vanishes, then the full curvature of $g$ vanishes.
\end{proposition}
\begin{proof}
Let $B\Subset\R^4$ be a simply connected coordinate ball.
Uniform inverse-metric convergence and weak-* convergence of first
derivatives give
\begin{equation*}
\sup_n\|\Gamma_n\|_{L^\infty(B)}<\infty,\qquad
\Gamma_n\rightharpoonup^*\Gamma\quad\hbox{in }L^\infty(B).
\end{equation*}
Flatness gives a parallel frame $E_n$ with a fixed invertible value
at the center. Gronwall's inequality along radial segments gives
\begin{equation*}
\sup_n\bigl(\|E_n\|_{L^\infty(B)}
             +\|E_n^{-1}\|_{L^\infty(B)}\bigr)<\infty.
\end{equation*}
Path independence gives the derivative equation and hence
\begin{equation*}
\partial_aE_n=-\Gamma_{n,a}E_n,\qquad
\sup_n\|\partial E_n\|_{L^\infty(B)}<\infty.
\end{equation*}
A subsequence converges uniformly to an invertible frame $E$.
Uniform convergence of $E_n$ and weak-* convergence of $\Gamma_n$
allow passage to the distributional equation
\begin{equation*}
\partial_aE=-\Gamma_aE.
\end{equation*}
Since $\Gamma$ is smooth, this equation makes $E$ smooth.
Commuting derivatives gives $\operatorname{Rm}(g)E=0$; invertibility
of $E$ implies $\operatorname{Rm}(g)=0$ on $B$. Such balls cover $M$.
\end{proof}

This excludes all flat pure-coordinate sequences as sources of a nonvacuum smooth limit in the admissible class; no near-identity assumption on the original coordinate maps is needed.

\section{Auxiliary identities and coordinate oscillations}

We record elementary examples and identities used to distinguish the
physical measure from the unprojected derivative defect.

\paragraph{Multiscale flat pullback.}
For $n\geq3$, set $F_n(s)=s+n^{-2}\sin(ns)+e^{-2n}\sin(e^ns)$. Then
\begin{equation*}
\begin{gathered}
F_n'=1+n^{-1}\cos(ns)+e^{-n}\cos(e^ns)>0,\\
F_n''=-\sin(ns)-\sin(e^ns),\qquad
F_n'''=-n\cos(ns)-e^n\cos(e^ns).
\end{gathered}
\end{equation*}
The bound $F_n^\prime\geq1-1/3-e^{-3}>0$ and bounded displacement make $F_n$ a global diffeomorphism. The pullback metric differs from Minkowski only in $g_{11}=(F_n')^2$, is exactly flat, has amplitude $O(n^{-1})$ and bounded first derivatives. At zero its amplitude is at least $2/n$ and its second derivative is $-2(1+n^{-1}+e^{-n})(n+e^n)$. Thus no common parameter can bound amplitude by $C\lambda_n$ and the second derivative by $C/\lambda_n$.

\paragraph{Spacelike raw defect.}
For the same sequence, $\partial_1h_{11}=-2(\sin(nx_1)+\sin(e^nx_1))+o(1)$ uniformly on compact sets. Each sine square has weak mean $1/2$, while the mixed frequencies $e^n\pm n$ tend to infinity. Hence $|\partial_1h_{11}|^2\rightharpoonup4$ and its raw directional defect lies at $\pm\dd x^1$, which are spacelike. Formula~\eqref{eq:mu-invariant} gives zero physical measure, as it must.

\paragraph{Weak energy is insufficient even algebraically.}
The constant covariant tensor $\operatorname{diag}(3,-1,2,2)$ has Lorentz trace zero and zero divergence. For timelike $(t,x)$, $3t^2-x_1^2+2x_2^2+2x_3^2>0$. It is not a positive real-covector second moment because its value on $\partial_{x_1}$ is negative.

\paragraph{The second-derivative product identity.}
If $f_n\to0$ uniformly and $\partial f_n$ is bounded in $L^\infty_{\loc}(\R^4)$, then
\begin{equation*}
f_n\partial_{ab}f_n=\partial_a(f_n\partial_bf_n)-\partial_af_n\partial_bf_n.
\end{equation*}
The divergence term tends strongly to zero in $H^{-1}_{\loc}(\R^4)$ since $f_n\partial_bf_n\to0$ in $L^2(K)$ for every compact $K\subset\R^4$. For $\chi\in C_c^\infty(\R^4)$, put $K=\supp\chi$.
H\"older's inequality gives
\begin{equation*}
\left|\int_{\R^4}(\partial_a\chi)f_n\partial_bf_n\,\dd x\right|
\leq\|\partial_a\chi\|_{L^1(\R^4)}
 \|f_n\|_{L^\infty(K)}\|\partial_bf_n\|_{L^\infty(K)}\longrightarrow0.
\end{equation*}
No conclusion follows for arbitrary sequence-dependent high-frequency tests from this fact alone.

\paragraph{Change of unit section.}
For $F_0(x,\xi)=(x,\xi/s_0)$ as in \eqref{eq:norm-change}, every degree-two test $b$ obeys $b\circ F_0=s_0^{-2}b$. Thus $F_{0*}(s_0^2\mu)$ preserves both the stress tests and the tests $H_Pa$. Unweighted pushforward generally does not.

\clearpage
\phantomsection
\addcontentsline{toc}{section}{References}
\begin{thebibliography}{9}\raggedright
\bibitem{BenyiOh} \'A. B\'enyi and T. Oh. Smoothing of commutators for a H\"ormander class of bilinear pseudodifferential operators. \emph{Journal of Fourier Analysis and Applications} 20 (2014), 282-300.
\bibitem{Burnett1989} G. A. Burnett. The high-frequency limit in general relativity. \emph{Journal of Mathematical Physics} 30 (1989), 90-96.
\bibitem{GW2011} S. R. Green and R. M. Wald. New framework for analyzing the effects of small scale inhomogeneities in cosmology. \emph{Physical Review D} 83 (2011), 084020.
\bibitem{GTC} A. Guerra and R. Teixeira da Costa. Oscillations in wave map systems and homogenization of the Einstein equations in symmetry. \emph{Archive for Rational Mechanics and Analysis} 249 (2025), Article 9.
\bibitem{HLelliptic} C. Huneau and J. Luk. Trilinear compensated compactness and Burnett's conjecture in general relativity. \emph{Annales scientifiques de l'\'Ecole normale sup\'erieure} (4) 57 (2024), 385-472.
\bibitem{HLwave} C. Huneau and J. Luk. Burnett's conjecture in generalized wave coordinates. arXiv:2403.03470v1 (2024).
\end{thebibliography}
\end{document}
